% The generated directory contains evidence-derived results and tables.
% The original project proposal remains in PROJECT/proposal.
\documentclass[conference,onecolumn]{IEEEtran}
\usepackage[T1]{fontenc}
\usepackage[utf8]{inputenc}
\usepackage{amsmath,amssymb}
\usepackage{array,booktabs}
\usepackage{cite}
\usepackage{textcomp}
\usepackage{url}
\usepackage[hidelinks]{hyperref}
\urlstyle{same}
\interdisplaylinepenalty=2500
\pagestyle{plain}

\begin{document}
\title{Baseline evaluation for speculative decoding and drafter feed-forward compression}
\author{\IEEEauthorblockN{Ayo Adetayo, Anh Nguyen, Pilar Maldonado, and Quang Nguyen}
\IEEEauthorblockA{Columbia University}}
\maketitle
\thispagestyle{plain}

\begin{abstract}
Speculative decoding uses a small drafter to propose tokens that a target language model verifies together. This study measures EAGLE-3, P-EAGLE, DFlash, and three pilot feed-forward widths on 30 prompts from ten benchmarks. Nine configurations complete 810 timed responses across three repetitions, using greedy decoding and one NVIDIA L40S GPU per run. DFlash leads the five initial configurations in output tokens per second and timed GPU efficiency on six benchmarks. The 75\% and 50\% pilots are slower than their matched autoregressive (AR) reference on every benchmark. The observed width results provide no consistent speed or compute-cost advantage from narrowing. Mean acceptance length is approximately 1.1--1.25 tokens per pilot verification round. Exact response mismatches remain unresolved, and task accuracy remains unmeasured. These findings motivate further drafter optimization and architecture experiments against ordinary generation and the released methods.
\end{abstract}

\begin{IEEEkeywords}
speculative decoding, parallel drafting, language-model inference, feed-forward compression, reproducible evaluation
\end{IEEEkeywords}

\section{Principal findings}
DFlash has the highest observed mean output tokens per second (OTPS) on six of the ten benchmarks in the initial comparison. The five compared rows are EAGLE-3, DFlash, P-EAGLE, AR SGLang, and AR vLLM ModelRunnerV1. P-EAGLE leads on MATH-500 and AIME 2025. EAGLE-3 leads on MT-Bench and CNN/DailyMail. DFlash is therefore the strongest current released-method speed reference in this trial. The conclusion concerns these 30 selected prompts and their observed responses. The response differences limit an exact-preservation interpretation of the speed gains.

Compute cost describes the GPU time and memory used to produce the responses. Within each one-GPU benchmark repetition, GPU-seconds per delivered token equal the reciprocal of OTPS. Across the recorded repetitions, DFlash has the lowest mean timed GPU time per token on six benchmarks. P-EAGLE and EAGLE-3 each lead on two. Table~\ref{tab:gpu-time} reports that compute-cost measure in the main results. The measured interval includes the benchmark's requests and client gaps. Model startup, warmups, and time outside the benchmark have separate accounting scopes.

The pilot results answer two separate questions. Across the observed width means, 100\% leads on six benchmarks and 75\% leads on four. The 50\% width leads on none. Within the matched continuation, both narrowed pilots are slower than ordinary AR generation on all ten benchmarks. Narrowing therefore provides no consistent observed speed or compute-cost advantage in this pilot.

\section{Research objective}
An autoregressive language model predicts the next token from the preceding context. Speculative decoding adds a cheaper model that proposes several tokens. The target then verifies the proposals in a common pass. A suitable acceptance procedure preserves the target distribution under its stated assumptions~\cite{leviathan2023}. Practical acceleration depends on drafting cost, verification cost, and the number of retained tokens per round.

The project examines whether feed-forward compression improves that tradeoff. A feed-forward network transforms each hidden vector within a transformer layer. Reducing its intermediate width reduces matrix sizes while preserving the layer's input and output dimensions. A smaller drafter can also change proposal quality. The experiment therefore measures both execution speed and accepted output.

The current study has two parts. The first part measures three existing speculative methods and their matching autoregressive references. The second part measures the P-EAGLE-style pilot at 100\%, 75\%, and 50\% feed-forward width. The pilot uses a Transformers reference. The observations characterize the current checkpoints and identify requirements for a controlled width study.

\section{Methods and implementation scope}
EAGLE-3 predicts tokens from target-model features and trains the drafter with a training-time test procedure~\cite{eagle3}. P-EAGLE drafts multiple future positions in parallel~\cite{peagle}. DFlash uses a block-diffusion drafter conditioned on target features~\cite{dflash}.

\subsection{Five initial configurations}
The autoregressive (AR) baseline runs the target model one token at a time. Its recorded responses and timings provide the reference for each speculative method. Matching the inference engine keeps the target's serving implementation consistent within each pair.

The current initial trial contains EAGLE-3, DFlash, and P-EAGLE, together with two autoregressive references. EAGLE-3 and DFlash run in SGLang and share the SGLang AR reference. P-EAGLE uses the H-Spec reproduction with vLLM ModelRunnerV1 and its corresponding AR reference~\cite{hspec}. These five configurations contain 450 completed timed request executions.

Each speedup uses the reference from the same inference engine and model runner. The comparison therefore measures the implemented method together with its serving path. Cross-engine differences remain part of the whole-implementation comparison. The local H-Spec configurations retain their reproduction labels because their target models and serving settings differ from the original papers' complete experimental setups.

The EAGLE-3 configuration uses three drafting steps, a top-$k$ value of four, and 16 total draft-tree tokens. The tree contains one previously generated root token and 15 new candidate nodes. The DFlash configuration uses a block size of ten, which supplies nine draft proposals per verification round. The P-EAGLE reproduction proposes seven linear draft tokens per round.

\subsection{Pilot width comparison}
The pilot compares four configurations: the Transformers AR reference and the 100\%, 75\%, and 50\% feed-forward variants. Each speculative variant uses the team's custom four-layer P-EAGLE-style drafter and proposes eight tokens per round. The released P-EAGLE reproduction keeps its own result row.

The supplied pilot dataset contains 200 training rows. Its verified SHA-256 checksum begins with \texttt{d13b9f2c7948d88d}. The import manifest records the full checksum. The supplied file contains no held-out examples. The imported checkpoints' original held-out data, training seed, and complete training history remain unresolved. The notebook describes one epoch and 50 optimizer updates per width. The imported artifacts retain their own provenance records.

The evaluation loader validates each checkpoint's weights, model dimensions, vocabulary, and target compatibility. All three pilot widths use the imported original checkpoints. Their run records preserve the checkpoint hashes and available training provenance. The three widths and the active Transformers AR reference contribute 360 completed timed responses.

\section{Evaluation procedure}
\subsection{Prompts and repetitions}
The trial takes the first three prepared prompts from each of ten benchmark families. The mathematics benchmarks are GSM8K, MATH-500, and AIME 2025. The code benchmarks are HumanEval, MBPP, and LiveCodeBench. MT-Bench, Alpaca, and Arena-Hard supply conversational or instruction-following prompts. CNN/DailyMail supplies article-summarization prompts.

The prepared bundle contains 30 unique prompts. Each prompt uses one turn and recurs in three timing repetitions. A complete configuration therefore contains 90 timed request executions. A benchmark block contains the three prompts for one benchmark in one repetition. Every configuration uses the same prepared input token IDs and prompt order. The saved bundle and request evidence preserve the exact text and token sequence.

The first-three selection supports a small, reproducible trial. Its observations describe these selected prompts. Three timing repetitions measure variation across repeated execution of the same inputs. Broader claims about benchmark performance require a larger prompt sample. In particular, a percentile calculated from three requests has limited resolution.

\subsection{Target, hardware, and generation settings}
The target is Meta Llama-3.1-8B-Instruct at the recorded model revision. The five initial configurations and the complete 100\% pilot run on one NVIDIA L40S on Insomnia node \texttt{ins061}. The 75\%, 50\%, and active Transformers AR configurations run on another L40S on \texttt{ins060}. Run records identify the physical GPU, driver, host, Slurm job, model revision, and prepared-bundle hash. A paired comparison requires matching target, inputs, output rules, repetitions, and physical GPU evidence.

The active tables use the complete original 100\% run because its continuation replay ends during launch preparation. Its different physical L40S limits conclusions about small width effects. Its paired AR speedup and agreement cells remain empty. The complete continuation AR supports the 75\% and 50\% comparisons. The original partial AR remains separate historical evidence.

Generation uses greedy decoding with temperature zero. The client submits one request at a time, so concurrency is one. The output ends at a configured terminal token or the 2,048-token limit. The primary stream records preserve native output token IDs, delivered token IDs, and stopping reasons. Delivered counts exclude the terminal stop token. Warmups, stream qualification, training, and diagnostic requests retain their own records and accounting scopes.

\subsection{Timing and throughput}
The measurement client runs on the compute node and connects to the model server through the loopback interface. A monotonic nanosecond clock records submission and streaming events. Time to first token (TTFT) measures the interval from submission to the first nonempty decoded output chunk. Metadata-only events contribute no first-content observation. A chunk can contain several tokens, so TTFT describes delivery to this client. Browser refresh time lies outside the measurement interval.

End-to-end request latency measures submission through stream completion. Output tokens per second (OTPS) divides the total delivered tokens in a benchmark block by the complete block interval. That interval includes the sequential requests and their client-side gaps. The following equation defines OTPS for benchmark $b$ and repetition $r$:
\begin{equation}
\mathrm{OTPS}_{b,r} =
\frac{\sum_{q=1}^{3} N_{b,r,q}}
{t^{\mathrm{end}}_{b,r}-t^{\mathrm{start}}_{b,r}},
\end{equation}
where $N_{b,r,q}$ is the delivered output-token count for request $q$. The paired speedup uses the matching AR measurement:
\begin{equation}
S_{b,r}=\frac{\mathrm{OTPS}^{\mathrm{method}}_{b,r}}
{\mathrm{OTPS}^{\mathrm{AR}}_{b,r}}.
\end{equation}
The tables summarize the repetition estimates with their mean, sample standard deviation, and eligible repetition count. A single eligible repetition contributes its observed value and count. Repetitions with missing or invalid evidence leave a documented gap.

\subsection{Acceptance and response agreement}
Mean acceptance length, denoted by $\tau$, measures committed progress per target verification round. Let $A_j$ count accepted draft tokens and $B_j$ count retained target correction or bonus tokens in round $j$. The project uses
\begin{equation}
\tau = \frac{\sum_{j=1}^{R}(A_j+B_j)}{R}.
\end{equation}
The numerator counts committed tokens before final output trimming. Any initial target token generated outside verification rounds has a separate count. Pooling the rounds gives each verification round equal weight. A supporting request-averaged measure first calculates each request's mean acceptance length, then averages those means equally.

For the pinned SGLang runtime, validated raw counters provide accepted draft tokens and verification rounds. The implementation retains one target token per round before output trimming, giving a committed count of accepted drafts plus rounds~\cite{sglangsource}. The exporter checks count aliases, proposal totals, histograms, retractions, and terminal trimming. DFlash's fixed-width histogram also supports positions one through nine.

Draft-token acceptance rate divides accepted draft tokens by all proposed draft tokens. The recorded EAGLE-3 setting contributes 15 candidate tree nodes to the denominator per round. DFlash contributes nine linear proposals, and P-EAGLE contributes seven linear proposals. EAGLE-3 counts candidates on alternative branches in its denominator and counts the accepted path in its numerator. Each linear method counts its proposed sequence. These distinct proposal structures limit direct comparisons of the acceptance percentages. The exported evidence preserves each native counter interval and its source-file hash.

EAGLE-3's branching proposal tree has a separate tree-prefix diagnostic. Prefix survival counts rounds whose tree reaches the stated depth and measures the fraction whose accepted path reaches that depth. Conditional acceptance counts rounds whose accepted prefix has a proposed child at the next depth. A leaf has no eligible next proposal. The exporter validates the captured tree, accepted path, verification result, and native round identity. EAGLE-3's paper defines its $n$-$\alpha$ study with a chain and a draft-feature-history convention~\cite{eagle3}. The report preserves that distinction in separate table variants.

Exact output agreement compares each speculative response with its matching AR response. A match requires identical delivered token IDs, output length, and stopping reason. This measure checks response preservation on the selected inputs. Task accuracy requires a separate scoring procedure for mathematics answers, code tests, summaries, or conversational responses.

\section{Results and interpretation}
% This file requires the array and booktabs packages and uses full-width benchmark panels.
% The three benchmark panels share one caption for each metric.
\begin{table*}[!tp]
\centering
\small
\setlength{\tabcolsep}{3pt}
\caption{Output tokens per second (OTPS) (output tokens/s).}
\label{tab:otps}
\textit{Mathematics}\par\smallskip
\begin{tabular}{@{}>{\raggedright\arraybackslash}p{3.7cm}rrr@{}}
\toprule
Method & GSM8K & MATH-500 & AIME 2025 \\
\midrule
EAGLE-3 & \(106.46 \pm 1.1\;[3]\) & \(103.23 \pm 0.086\;[3]\) & \(104.17 \pm 0.26\;[3]\) \\
P-EAGLE & \(117.50 \pm 0.82\;[3]\) & \(\mathbf{153.49} \pm 0.012\;[3]\) & \(\mathbf{157.94} \pm 0.0071\;[3]\) \\
DFlash & \(\mathbf{213.37} \pm 0.29\;[3]\) & \(138.76 \pm 0.31\;[3]\) & \(154.89 \pm 0.062\;[3]\) \\
AR \textperiodcentered{} SGLang & \(43.27 \pm 0.046\;[3]\) & \(43.44 \pm 0.0078\;[3]\) & \(43.34 \pm 0.058\;[3]\) \\
AR \textperiodcentered{} vLLM ModelRunnerV1 & \(46.28 \pm 0.0013\;[3]\) & \(46.47 \pm 0.002\;[3]\) & \(46.30 \pm 0.00078\;[3]\) \\
Pilot FFN 100\% & \(39.02 \pm 0.036\;[3]\) & \(33.96 \pm 0.022\;[3]\) & \(32.42 \pm 0.026\;[3]\) \\
Pilot FFN 75\% & \(39.02 \pm 0.045\;[3]\) & \(31.92 \pm 2.8\;[3]\) & \(32.74 \pm 0.054\;[3]\) \\
Pilot FFN 50\% & \(37.59 \pm 0.082\;[3]\) & \(32.85 \pm 0.22\;[3]\) & \(31.56 \pm 0.035\;[3]\) \\
AR \textperiodcentered{} Transformers & \(40.17 \pm 0.15\;[3]\) & \(40.01 \pm 0.027\;[3]\) & \(39.12 \pm 0.05\;[3]\) \\
\bottomrule
\end{tabular}
\par\medskip
\textit{Code}\par\smallskip
\begin{tabular}{@{}>{\raggedright\arraybackslash}p{3.7cm}rrr@{}}
\toprule
Method & HumanEval & MBPP & LiveCodeBench \\
\midrule
EAGLE-3 & \(100.29 \pm 0.28\;[3]\) & \(110.69 \pm 1.1\;[3]\) & \(99.48 \pm 0.39\;[3]\) \\
P-EAGLE & \(114.10 \pm 0.011\;[3]\) & \(124.91 \pm 0.0019\;[3]\) & \(107.91 \pm 0.01\;[3]\) \\
DFlash & \(\mathbf{137.18} \pm 1.4\;[3]\) & \(\mathbf{155.12} \pm 0.35\;[3]\) & \(\mathbf{117.21} \pm 0.21\;[3]\) \\
AR \textperiodcentered{} SGLang & \(42.88 \pm 0.081\;[3]\) & \(43.27 \pm 0.056\;[3]\) & \(43.27 \pm 0.043\;[3]\) \\
AR \textperiodcentered{} vLLM ModelRunnerV1 & \(45.72 \pm 0.0035\;[3]\) & \(46.17 \pm 0.029\;[3]\) & \(46.21 \pm 0.021\;[3]\) \\
Pilot FFN 100\% & \(36.12 \pm 0.052\;[3]\) & \(37.10 \pm 0.057\;[3]\) & \(35.49 \pm 0.18\;[3]\) \\
Pilot FFN 75\% & \(35.42 \pm 0.064\;[3]\) & \(35.39 \pm 0.044\;[3]\) & \(35.10 \pm 0.082\;[3]\) \\
Pilot FFN 50\% & \(35.66 \pm 0.13\;[3]\) & \(36.68 \pm 0.06\;[3]\) & \(34.45 \pm 0.03\;[3]\) \\
AR \textperiodcentered{} Transformers & \(39.73 \pm 0.077\;[3]\) & \(40.03 \pm 0.042\;[3]\) & \(39.87 \pm 0.13\;[3]\) \\
\bottomrule
\end{tabular}
\par\medskip
\textit{Conversation, instruction following, and summarization}\par\smallskip
\begin{tabular}{@{}>{\raggedright\arraybackslash}p{3.7cm}rrrr@{}}
\toprule
Method & MT-Bench & Alpaca & Arena-Hard & CNN/DailyMail \\
\midrule
EAGLE-3 & \(\mathbf{110.29} \pm 0.77\;[3]\) & \(116.71 \pm 0.37\;[3]\) & \(94.66 \pm 0.49\;[3]\) & \(\mathbf{97.98} \pm 1.2\;[3]\) \\
P-EAGLE & \(82.75 \pm 0.37\;[3]\) & \(105.33 \pm 0.0069\;[3]\) & \(93.02 \pm 0.12\;[3]\) & \(83.56 \pm 0.012\;[3]\) \\
DFlash & \(109.73 \pm 0.13\;[3]\) & \(\mathbf{117.62} \pm 0.39\;[3]\) & \(\mathbf{112.30} \pm 0.14\;[3]\) & \(97.67 \pm 0.067\;[3]\) \\
AR \textperiodcentered{} SGLang & \(43.37 \pm 0.067\;[3]\) & \(43.38 \pm 0.011\;[3]\) & \(43.25 \pm 0.018\;[3]\) & \(43.00 \pm 0.044\;[3]\) \\
AR \textperiodcentered{} vLLM ModelRunnerV1 & \(46.43 \pm 0.0034\;[3]\) & \(46.37 \pm 0.0016\;[3]\) & \(46.19 \pm 0.0067\;[3]\) & \(45.80 \pm 0.0037\;[3]\) \\
Pilot FFN 100\% & \(36.30 \pm 0.038\;[3]\) & \(36.34 \pm 0.055\;[3]\) & \(34.89 \pm 0.1\;[3]\) & \(34.28 \pm 0.019\;[3]\) \\
Pilot FFN 75\% & \(36.33 \pm 0.0059\;[3]\) & \(36.42 \pm 0.045\;[3]\) & \(34.36 \pm 0.012\;[3]\) & \(34.48 \pm 0.044\;[3]\) \\
Pilot FFN 50\% & \(35.71 \pm 0.11\;[3]\) & \(35.54 \pm 0.13\;[3]\) & \(33.96 \pm 0.095\;[3]\) & \(33.25 \pm 0.078\;[3]\) \\
AR \textperiodcentered{} Transformers & \(39.97 \pm 0.093\;[3]\) & \(40.02 \pm 0.029\;[3]\) & \(39.76 \pm 0.039\;[3]\) & \(39.39 \pm 0.013\;[3]\) \\
\bottomrule
\end{tabular}
\par\medskip
\begin{minipage}{\linewidth}\footnotesize
Brackets give the number of eligible repetitions. With two or more repetitions, entries give the mean and sample standard deviation. A single repetition supplies one observed value. Each repetition uses the same three prompts. Empty cells have no validated measurement. The target is Meta Llama-3.1-8B-Instruct on an NVIDIA L40S GPU. Generation uses greedy decoding, one active request, and a 2,048-token output limit. Pilot FFN rows identify the notebook's P-EAGLE-style drafter at feed-forward widths of 100\%, 75\%, and 50\%. Run metadata identifies each checkpoint's origin. Their AR reference uses Transformers. The P-EAGLE row uses the H-Spec reproduction. Acceptance uses validated native counters or matched diagnostic runs. The evidence records identify each measurement scope. Committed tokens are counted before output trimming.

The bold mean is the best observed value in each benchmark column. The comparison covers the recorded implementations under their stated configurations. At least two eligible method entries are required. Ties at the displayed precision share bold formatting. Bold formatting conveys an observed comparison and makes no statistical-significance claim. Brackets show the available repetition count. Response differences remain part of the interpretation. Higher values are preferred.
\end{minipage}
\end{table*}

\begin{table*}[!tp]
\centering
\small
\setlength{\tabcolsep}{3pt}
\caption{Speedup over AR (\(\times\)).}
\label{tab:speedup}
\textit{Mathematics}\par\smallskip
\begin{tabular}{@{}>{\raggedright\arraybackslash}p{3.7cm}rrr@{}}
\toprule
Method & GSM8K & MATH-500 & AIME 2025 \\
\midrule
EAGLE-3 & \(2.46 \pm 0.029\;[3]\) & \(2.38 \pm 0.0016\;[3]\) & \(2.40 \pm 0.0057\;[3]\) \\
P-EAGLE & \(2.54 \pm 0.018\;[3]\) & \(\mathbf{3.30} \pm 0.00034\;[3]\) & \(3.41 \pm 0.00021\;[3]\) \\
DFlash & \(\mathbf{4.93} \pm 0.0036\;[3]\) & \(3.19 \pm 0.0075\;[3]\) & \(\mathbf{3.57} \pm 0.0062\;[3]\) \\
Pilot FFN 100\% &  &  &  \\
Pilot FFN 75\% & \(0.97 \pm 0.0048\;[3]\) & \(0.80 \pm 0.07\;[3]\) & \(0.84 \pm 0.0015\;[3]\) \\
Pilot FFN 50\% & \(0.94 \pm 0.0024\;[3]\) & \(0.82 \pm 0.006\;[3]\) & \(0.81 \pm 0.00092\;[3]\) \\
\bottomrule
\end{tabular}
\par\medskip
\textit{Code}\par\smallskip
\begin{tabular}{@{}>{\raggedright\arraybackslash}p{3.7cm}rrr@{}}
\toprule
Method & HumanEval & MBPP & LiveCodeBench \\
\midrule
EAGLE-3 & \(2.34 \pm 0.0059\;[3]\) & \(2.56 \pm 0.025\;[3]\) & \(2.30 \pm 0.0098\;[3]\) \\
P-EAGLE & \(2.50 \pm 0.00022\;[3]\) & \(2.71 \pm 0.0017\;[3]\) & \(2.34 \pm 0.0012\;[3]\) \\
DFlash & \(\mathbf{3.20} \pm 0.028\;[3]\) & \(\mathbf{3.59} \pm 0.0073\;[3]\) & \(\mathbf{2.71} \pm 0.0022\;[3]\) \\
Pilot FFN 100\% &  &  &  \\
Pilot FFN 75\% & \(0.89 \pm 0.0016\;[3]\) & \(0.88 \pm 0.0016\;[3]\) & \(0.88 \pm 0.0021\;[3]\) \\
Pilot FFN 50\% & \(0.90 \pm 0.0026\;[3]\) & \(0.92 \pm 0.0024\;[3]\) & \(0.86 \pm 0.0029\;[3]\) \\
\bottomrule
\end{tabular}
\par\medskip
\textit{Conversation, instruction following, and summarization}\par\smallskip
\begin{tabular}{@{}>{\raggedright\arraybackslash}p{3.7cm}rrrr@{}}
\toprule
Method & MT-Bench & Alpaca & Arena-Hard & CNN/DailyMail \\
\midrule
EAGLE-3 & \(\mathbf{2.54} \pm 0.014\;[3]\) & \(2.69 \pm 0.0091\;[3]\) & \(2.19 \pm 0.011\;[3]\) & \(\mathbf{2.28} \pm 0.027\;[3]\) \\
P-EAGLE & \(1.78 \pm 0.0078\;[3]\) & \(2.27 \pm 0.00021\;[3]\) & \(2.01 \pm 0.0028\;[3]\) & \(1.82 \pm 0.00018\;[3]\) \\
DFlash & \(2.53 \pm 0.0027\;[3]\) & \(\mathbf{2.71} \pm 0.0083\;[3]\) & \(\mathbf{2.60} \pm 0.0027\;[3]\) & \(2.27 \pm 0.0038\;[3]\) \\
Pilot FFN 100\% &  &  &  &  \\
Pilot FFN 75\% & \(0.91 \pm 0.0021\;[3]\) & \(0.91 \pm 0.00065\;[3]\) & \(0.86 \pm 0.00082\;[3]\) & \(0.88 \pm 0.00085\;[3]\) \\
Pilot FFN 50\% & \(0.89 \pm 0.00097\;[3]\) & \(0.89 \pm 0.0031\;[3]\) & \(0.85 \pm 0.0028\;[3]\) & \(0.84 \pm 0.0019\;[3]\) \\
\bottomrule
\end{tabular}
\par\medskip
\begin{minipage}{\linewidth}\footnotesize
Brackets give the number of eligible repetitions. With two or more repetitions, entries give the mean and sample standard deviation. A single repetition supplies one observed value. Each repetition uses the same three prompts. Empty cells have no validated measurement. The target is Meta Llama-3.1-8B-Instruct on an NVIDIA L40S GPU. Generation uses greedy decoding, one active request, and a 2,048-token output limit. Pilot FFN rows identify the notebook's P-EAGLE-style drafter at feed-forward widths of 100\%, 75\%, and 50\%. Run metadata identifies each checkpoint's origin. Their AR reference uses Transformers. The P-EAGLE row uses the H-Spec reproduction. Acceptance uses validated native counters or matched diagnostic runs. The evidence records identify each measurement scope. Committed tokens are counted before output trimming.

The bold mean is the best observed value in each benchmark column. The comparison covers the recorded implementations under their stated configurations. At least two eligible method entries are required. Ties at the displayed precision share bold formatting. Bold formatting conveys an observed comparison and makes no statistical-significance claim. Brackets show the available repetition count. Response differences remain part of the interpretation. Higher values are preferred.
\end{minipage}
\end{table*}

\begin{table*}[!tp]
\centering
\small
\setlength{\tabcolsep}{3pt}
\caption{Allocated GPU time per million output tokens (GPU-seconds/million output tokens).}
\label{tab:gpu-time}
\textit{Mathematics}\par\smallskip
\begin{tabular}{@{}>{\raggedright\arraybackslash}p{3.7cm}rrr@{}}
\toprule
Method & GSM8K & MATH-500 & AIME 2025 \\
\midrule
EAGLE-3 & \(9394 \pm 100\;[3]\) & \(9687 \pm 8.1\;[3]\) & \(9600 \pm 24\;[3]\) \\
P-EAGLE & \(8511 \pm 60\;[3]\) & \(\mathbf{6515} \pm 0.5\;[3]\) & \(\mathbf{6331} \pm 0.29\;[3]\) \\
DFlash & \(\mathbf{4687} \pm 6.3\;[3]\) & \(7206 \pm 16\;[3]\) & \(6456 \pm 2.6\;[3]\) \\
AR \textperiodcentered{} SGLang & \(23110 \pm 25\;[3]\) & \(23020 \pm 4.1\;[3]\) & \(23080 \pm 31\;[3]\) \\
AR \textperiodcentered{} vLLM ModelRunnerV1 & \(21610 \pm 0.59\;[3]\) & \(21520 \pm 0.95\;[3]\) & \(21600 \pm 0.36\;[3]\) \\
Pilot FFN 100\% & \(25630 \pm 23\;[3]\) & \(29440 \pm 19\;[3]\) & \(30850 \pm 25\;[3]\) \\
Pilot FFN 75\% & \(25630 \pm 30\;[3]\) & \(31500 \pm 2900\;[3]\) & \(30540 \pm 50\;[3]\) \\
Pilot FFN 50\% & \(26600 \pm 58\;[3]\) & \(30440 \pm 210\;[3]\) & \(31690 \pm 35\;[3]\) \\
AR \textperiodcentered{} Transformers & \(24900 \pm 94\;[3]\) & \(24990 \pm 17\;[3]\) & \(25560 \pm 32\;[3]\) \\
\bottomrule
\end{tabular}
\par\medskip
\textit{Code}\par\smallskip
\begin{tabular}{@{}>{\raggedright\arraybackslash}p{3.7cm}rrr@{}}
\toprule
Method & HumanEval & MBPP & LiveCodeBench \\
\midrule
EAGLE-3 & \(9971 \pm 28\;[3]\) & \(9035 \pm 93\;[3]\) & \(10050 \pm 40\;[3]\) \\
P-EAGLE & \(8764 \pm 0.87\;[3]\) & \(8006 \pm 0.12\;[3]\) & \(9267 \pm 0.87\;[3]\) \\
DFlash & \(\mathbf{7290} \pm 73\;[3]\) & \(\mathbf{6447} \pm 15\;[3]\) & \(\mathbf{8532} \pm 15\;[3]\) \\
AR \textperiodcentered{} SGLang & \(23320 \pm 44\;[3]\) & \(23110 \pm 30\;[3]\) & \(23110 \pm 23\;[3]\) \\
AR \textperiodcentered{} vLLM ModelRunnerV1 & \(21870 \pm 1.7\;[3]\) & \(21660 \pm 14\;[3]\) & \(21640 \pm 9.7\;[3]\) \\
Pilot FFN 100\% & \(27690 \pm 40\;[3]\) & \(26950 \pm 41\;[3]\) & \(28180 \pm 140\;[3]\) \\
Pilot FFN 75\% & \(28230 \pm 51\;[3]\) & \(28260 \pm 35\;[3]\) & \(28490 \pm 66\;[3]\) \\
Pilot FFN 50\% & \(28040 \pm 100\;[3]\) & \(27260 \pm 45\;[3]\) & \(29020 \pm 25\;[3]\) \\
AR \textperiodcentered{} Transformers & \(25170 \pm 49\;[3]\) & \(24980 \pm 26\;[3]\) & \(25080 \pm 82\;[3]\) \\
\bottomrule
\end{tabular}
\par\medskip
\textit{Conversation, instruction following, and summarization}\par\smallskip
\begin{tabular}{@{}>{\raggedright\arraybackslash}p{3.7cm}rrrr@{}}
\toprule
Method & MT-Bench & Alpaca & Arena-Hard & CNN/DailyMail \\
\midrule
EAGLE-3 & \(\mathbf{9067} \pm 63\;[3]\) & \(8568 \pm 27\;[3]\) & \(10560 \pm 55\;[3]\) & \(\mathbf{10210} \pm 130\;[3]\) \\
P-EAGLE & \(12080 \pm 54\;[3]\) & \(9494 \pm 0.62\;[3]\) & \(10750 \pm 13\;[3]\) & \(11970 \pm 1.8\;[3]\) \\
DFlash & \(9113 \pm 11\;[3]\) & \(\mathbf{8502} \pm 28\;[3]\) & \(\mathbf{8904} \pm 11\;[3]\) & \(10240 \pm 7\;[3]\) \\
AR \textperiodcentered{} SGLang & \(23060 \pm 36\;[3]\) & \(23050 \pm 5.9\;[3]\) & \(23120 \pm 9.8\;[3]\) & \(23260 \pm 24\;[3]\) \\
AR \textperiodcentered{} vLLM ModelRunnerV1 & \(21540 \pm 1.6\;[3]\) & \(21570 \pm 0.74\;[3]\) & \(21650 \pm 3.2\;[3]\) & \(21830 \pm 1.8\;[3]\) \\
Pilot FFN 100\% & \(27550 \pm 29\;[3]\) & \(27520 \pm 42\;[3]\) & \(28660 \pm 83\;[3]\) & \(29170 \pm 16\;[3]\) \\
Pilot FFN 75\% & \(27520 \pm 4.5\;[3]\) & \(27460 \pm 34\;[3]\) & \(29100 \pm 10\;[3]\) & \(29000 \pm 37\;[3]\) \\
Pilot FFN 50\% & \(28000 \pm 86\;[3]\) & \(28140 \pm 110\;[3]\) & \(29440 \pm 82\;[3]\) & \(30070 \pm 71\;[3]\) \\
AR \textperiodcentered{} Transformers & \(25020 \pm 58\;[3]\) & \(24990 \pm 18\;[3]\) & \(25150 \pm 25\;[3]\) & \(25380 \pm 8.1\;[3]\) \\
\bottomrule
\end{tabular}
\par\medskip
\begin{minipage}{\linewidth}\footnotesize
Brackets give the number of eligible repetitions. With two or more repetitions, entries give the mean and sample standard deviation. A single repetition supplies one observed value. Each repetition uses the same three prompts. Empty cells have no validated measurement. The target is Meta Llama-3.1-8B-Instruct on an NVIDIA L40S GPU. Generation uses greedy decoding, one active request, and a 2,048-token output limit. Pilot FFN rows identify the notebook's P-EAGLE-style drafter at feed-forward widths of 100\%, 75\%, and 50\%. Run metadata identifies each checkpoint's origin. Their AR reference uses Transformers. The P-EAGLE row uses the H-Spec reproduction. Acceptance uses validated native counters or matched diagnostic runs. The evidence records identify each measurement scope. Committed tokens are counted before output trimming.
Compute cost uses each timed benchmark interval and the GPU count from its final Slurm allocation. The interval includes client gaps between requests. Startup, warmups, and time outside the benchmark have separate whole-job accounting. Per-prompt GPU-seconds multiply the allocated GPU count by that prompt's completion time in seconds.
The bold mean is the best observed value in each benchmark column. The comparison covers the recorded implementations under their stated configurations. At least two eligible method entries are required. Ties at the displayed precision share bold formatting. Bold formatting conveys an observed comparison and makes no statistical-significance claim. Brackets show the available repetition count. Response differences remain part of the interpretation. Lower values are preferred.
\end{minipage}
\end{table*}

\begin{table*}[!tp]
\centering
\small
\setlength{\tabcolsep}{3pt}
\caption{Mean acceptance length (\(\tau\)) (committed tokens/round).}
\label{tab:acceptance-length}
\textit{Mathematics}\par\smallskip
\begin{tabular}{@{}>{\raggedright\arraybackslash}p{3.7cm}rrr@{}}
\toprule
Method & GSM8K & MATH-500 & AIME 2025 \\
\midrule
EAGLE-3 & \(3.443 \pm 0\;[3]\) & \(3.306 \pm 0\;[3]\) & \(3.328 \pm 0\;[3]\) \\
P-EAGLE &  &  &  \\
DFlash & \(\mathbf{6.656} \pm 0\;[3]\) & \(\mathbf{4.297} \pm 0\;[3]\) & \(\mathbf{4.803} \pm 0\;[3]\) \\
Pilot FFN 100\% & \(1.245 \pm 0\;[3]\) & \(1.095 \pm 0\;[3]\) & \(1.121 \pm 0\;[3]\) \\
Pilot FFN 75\% & \(1.248 \pm 0\;[3]\) & \(1.085 \pm 0\;[3]\) & \(1.133 \pm 0\;[3]\) \\
Pilot FFN 50\% & \(1.216 \pm 0\;[3]\) & \(1.072 \pm 0\;[3]\) & \(1.106 \pm 0\;[3]\) \\
\bottomrule
\end{tabular}
\par\medskip
\textit{Code}\par\smallskip
\begin{tabular}{@{}>{\raggedright\arraybackslash}p{3.7cm}rrr@{}}
\toprule
Method & HumanEval & MBPP & LiveCodeBench \\
\midrule
EAGLE-3 & \(3.365 \pm 0\;[3]\) & \(3.624 \pm 0\;[3]\) & \(3.218 \pm 0\;[3]\) \\
P-EAGLE &  &  &  \\
DFlash & \(\mathbf{4.429} \pm 0\;[3]\) & \(\mathbf{4.909} \pm 0\;[3]\) & \(\mathbf{3.67} \pm 0\;[3]\) \\
Pilot FFN 100\% & \(1.159 \pm 0\;[3]\) & \(1.181 \pm 0\;[3]\) & \(1.157 \pm 0\;[3]\) \\
Pilot FFN 75\% & \(1.136 \pm 0\;[3]\) & \(1.13 \pm 0\;[3]\) & \(1.148 \pm 0\;[3]\) \\
Pilot FFN 50\% & \(1.163 \pm 0\;[3]\) & \(1.184 \pm 0\;[3]\) & \(1.144 \pm 0\;[3]\) \\
\bottomrule
\end{tabular}
\par\medskip
\textit{Conversation, instruction following, and summarization}\par\smallskip
\begin{tabular}{@{}>{\raggedright\arraybackslash}p{3.7cm}rrrr@{}}
\toprule
Method & MT-Bench & Alpaca & Arena-Hard & CNN/DailyMail \\
\midrule
EAGLE-3 & \(\mathbf{3.558} \pm 0\;[3]\) & \(\mathbf{3.753} \pm 0\;[3]\) & \(3.045 \pm 0\;[3]\) & \(\mathbf{3.145} \pm 0\;[3]\) \\
P-EAGLE &  &  &  &  \\
DFlash & \(3.399 \pm 0\;[3]\) & \(3.66 \pm 0\;[3]\) & \(\mathbf{3.503} \pm 0\;[3]\) & \(3.085 \pm 0\;[3]\) \\
Pilot FFN 100\% & \(1.178 \pm 0\;[3]\) & \(1.158 \pm 0\;[3]\) & \(1.153 \pm 0\;[3]\) & \(1.154 \pm 0\;[3]\) \\
Pilot FFN 75\% & \(1.181 \pm 0\;[3]\) & \(1.164 \pm 0\;[3]\) & \(1.139 \pm 0\;[3]\) & \(1.163 \pm 0\;[3]\) \\
Pilot FFN 50\% & \(1.175 \pm 0\;[3]\) & \(1.15 \pm 0\;[3]\) & \(1.139 \pm 0\;[3]\) & \(1.136 \pm 0\;[3]\) \\
\bottomrule
\end{tabular}
\par\medskip
\begin{minipage}{\linewidth}\footnotesize
Brackets give the number of eligible repetitions. With two or more repetitions, entries give the mean and sample standard deviation. A single repetition supplies one observed value. Each repetition uses the same three prompts. Empty cells have no validated measurement. The target is Meta Llama-3.1-8B-Instruct on an NVIDIA L40S GPU. Generation uses greedy decoding, one active request, and a 2,048-token output limit. Pilot FFN rows identify the notebook's P-EAGLE-style drafter at feed-forward widths of 100\%, 75\%, and 50\%. Run metadata identifies each checkpoint's origin. Their AR reference uses Transformers. The P-EAGLE row uses the H-Spec reproduction. Acceptance uses validated native counters or matched diagnostic runs. The evidence records identify each measurement scope. Committed tokens are counted before output trimming.

The bold mean is the best observed value in each benchmark column. The comparison covers the recorded implementations under their stated configurations. At least two eligible method entries are required. Ties at the displayed precision share bold formatting. Bold formatting conveys an observed comparison and makes no statistical-significance claim. Brackets show the available repetition count. Response differences remain part of the interpretation. Higher values are preferred.
\end{minipage}
\end{table*}

\subsection{Available observations}

The initial record set contains 150 of 150 planned benchmark blocks. Each complete block contains three request executions.

The team record set contains 120 of 120 planned benchmark blocks. Each complete block contains three request executions.

The current initial scope contains 450 completed timed request executions out of 450 planned executions.

The 30 unique prompts recur across three timing repetitions. A complete configuration therefore contains 90 timed request executions.

Pilot FFN 100\% uses an imported original team checkpoint. The run evidence preserves available training provenance and leaves unknown fields unresolved.

Pilot FFN 75\% uses an imported original team checkpoint. The run evidence preserves available training provenance and leaves unknown fields unresolved.

Pilot FFN 50\% uses an imported original team checkpoint. The run evidence preserves available training provenance and leaves unknown fields unresolved.

\subsection{Saved diagnostic coverage}

EAGLE-3 has validated complete-round timing for 30 benchmark repetitions across 10 benchmarks. Those intervals cover 90 diagnostic request executions.

DFlash has validated complete-round timing for 4 benchmark repetitions across 4 benchmarks. Those intervals cover 12 diagnostic request executions.

These separate diagnostic executions include synchronization and trace-construction overhead. Primary requests supply the throughput and response-time tables. The appendix reports the eligible diagnostic measurements and their repetition counts.

\subsection{Throughput and paired speedup}

EAGLE-3 has benchmark mean speedups from 2.189\(\times\) on Arena-Hard to 2.690\(\times\) on Alpaca across 10 measured benchmarks. Table entries retain the repetition counts and standard deviations.

P-EAGLE has benchmark mean speedups from 1.782\(\times\) on MT-Bench to 3.411\(\times\) on AIME 2025 across 10 measured benchmarks. Table entries retain the repetition counts and standard deviations.

DFlash has benchmark mean speedups from 2.271\(\times\) on CNN/DailyMail to 4.931\(\times\) on GSM8K across 10 measured benchmarks. Table entries retain the repetition counts and standard deviations.

Pilot FFN 75\% has benchmark mean speedups from 0.798\(\times\) on MATH-500 to 0.971\(\times\) on GSM8K across 10 measured benchmarks. Table entries retain the repetition counts and standard deviations.

Pilot FFN 50\% has benchmark mean speedups from 0.807\(\times\) on AIME 2025 to 0.936\(\times\) on GSM8K across 10 measured benchmarks. Table entries retain the repetition counts and standard deviations.

The throughput ratio describes delivered output per elapsed second under the recorded prompts and settings. Response length and output-agreement results define the scope of each comparison.

\subsection{Team pilot observations}

Pilot FFN 100\% has observed benchmark mean OTPS from 32.42 tokens/s on AIME 2025 to 39.02 tokens/s on GSM8K across 10 benchmarks. These measurements use the team's recorded Transformers implementation.

The active AR reference supports paired speedup on 0 of these 10 benchmarks. Compatibility checks leave the remaining speedup and exact-agreement cells empty. The measured OTPS remains available for each completed method block.

Pilot FFN 100\% has benchmark mean acceptance length from 1.095 to 1.245 committed tokens per verification round. The calculation uses the saved primary round counts before final output trimming.

Pilot FFN 75\% has observed benchmark mean OTPS from 31.92 tokens/s on MATH-500 to 39.02 tokens/s on GSM8K across 10 benchmarks. These measurements use the team's recorded Transformers implementation.

Pilot FFN 75\% has benchmark mean acceptance length from 1.085 to 1.248 committed tokens per verification round. The calculation uses the saved primary round counts before final output trimming.

Pilot FFN 50\% has observed benchmark mean OTPS from 31.56 tokens/s on AIME 2025 to 37.59 tokens/s on GSM8K across 10 benchmarks. These measurements use the team's recorded Transformers implementation.

Pilot FFN 50\% has benchmark mean acceptance length from 1.072 to 1.216 committed tokens per verification round. The calculation uses the saved primary round counts before final output trimming.

\subsection{Exact response agreement}

EAGLE-3 matches its AR reference on 33 of 90 checked request executions (36.7\%). All observed repetitions match for 11 of 30 distinct prompts. Each prompt has the same agreement outcome across its observed repetitions. The earliest recorded token divergence occurs at returned token 4. Output lengths match in every observed repetition for 13 of 30 distinct prompts. Across one response per distinct prompt, the method delivers 10,534 tokens and its AR reference delivers 12,513. Each observed repetition has these same per-prompt lengths.

P-EAGLE matches its AR reference on 42 of 90 checked request executions (46.7\%). All observed repetitions match for 14 of 30 distinct prompts. Each prompt has the same agreement outcome across its observed repetitions. The earliest recorded token divergence occurs at returned token 4. Output lengths match in every observed repetition for 15 of 30 distinct prompts. Across one response per distinct prompt, the method delivers 11,210 tokens and its AR reference delivers 10,560. Each observed repetition has these same per-prompt lengths.

DFlash matches its AR reference on 39 of 90 checked request executions (43.3\%). All observed repetitions match for 13 of 30 distinct prompts. Each prompt has the same agreement outcome across its observed repetitions. The earliest recorded token divergence occurs at returned token 4. Output lengths match in every observed repetition for 15 of 30 distinct prompts. Across one response per distinct prompt, the method delivers 12,993 tokens and its AR reference delivers 12,513. Each observed repetition has these same per-prompt lengths.

Pilot FFN 75\% matches its AR reference on 36 of 90 checked request executions (40.0\%). All observed repetitions match for 12 of 30 distinct prompts. Each prompt has the same agreement outcome across its observed repetitions. The earliest recorded token divergence occurs at returned token 4. Output lengths match in every observed repetition for 13 of 30 distinct prompts. Across one response per distinct prompt, the method delivers 11,579 tokens and its AR reference delivers 9,816. Each observed repetition has these same per-prompt lengths.

Pilot FFN 50\% matches its AR reference on 36 of 90 checked request executions (40.0\%). All observed repetitions match for 12 of 30 distinct prompts. Each prompt has the same agreement outcome across its observed repetitions. The earliest recorded token divergence occurs at returned token 4. Output lengths match in every observed repetition for 13 of 30 distinct prompts. Across one response per distinct prompt, the method delivers 11,579 tokens and its AR reference delivers 9,816. Each observed repetition has these same per-prompt lengths.

The response mismatches are a material finding. Their cause remains unresolved. The measured throughput gains describe the recorded executions. An exact-response-preservation claim requires resolving these mismatches. A causal explanation requires checking target scores, input token IDs, verification behavior, and cache state at the first divergent position.

Benchmark task accuracy remains unmeasured. Exact token agreement evaluates reproduction of the AR response. Math correctness, code pass@1, summarization quality, and conversational judging require separate scoring procedures.

\begin{table*}[!htp]
\centering\small
\caption{Allocated compute resources for the evaluation jobs.}
\label{tab:allocations}
\begin{tabular}{@{}lp{5.0cm}lrr@{}}
\toprule
Job ID & Purpose & GPU model / count & GPU-hours & CPU-hours \\
\midrule
14056875 & Released-method evaluation and the original pilot attempt. & NVIDIA L40S / 1 & 3.0053 & 48.0844 \\
14058516 & Pilot width evaluation and its AR reference. & NVIDIA L40S / 1 & 1.5064 & 24.1022 \\
\bottomrule\end{tabular}
\par\smallskip\begin{minipage}{\linewidth}\footnotesize
Each parent allocation appears once. These totals include all elapsed allocation time, including setup, model loading, qualification, warmups, recorded requests, and incomplete attempts. The table covers the evaluation allocations used by these result rows. The project compute ledger separately preserves earlier preparation jobs. Final Slurm records and their hashes accompany the exported evidence. CPU-hours count allocated CPU cores multiplied by elapsed hours.
\end{minipage}\end{table*}

\subsection{Results by benchmark category}
\paragraph*{Code throughput}
DFlash has the highest observed mean OTPS on HumanEval, MBPP, and LiveCodeBench: 137.18, 155.12, and 117.21 tokens/s, respectively (Table~\ref{tab:otps}). Both narrowed pilots remain below their matched Transformers AR reference on all three code benchmarks.

\paragraph*{Code GPU time}
DFlash also uses the least timed GPU time per million output tokens on HumanEval, MBPP, and LiveCodeBench: 7{,}290, 6{,}447, and 8{,}532 GPU-seconds, respectively (Table~\ref{tab:gpu-time}). Both narrowed pilots use more GPU time per output token than their matched AR reference on all three benchmarks.

\paragraph*{Conversation and instruction-following throughput}
EAGLE-3 leads MT-Bench at 110.29 tokens/s, while DFlash leads Alpaca and Arena-Hard at 117.62 and 112.30 tokens/s (Table~\ref{tab:otps}). The 75\% pilot exceeds the 50\% pilot on these three benchmarks, and both remain below their matched AR reference.

\paragraph*{Conversation and instruction-following GPU time}
EAGLE-3 uses the least GPU time per million output tokens on MT-Bench at 9{,}067 GPU-seconds, and DFlash leads Alpaca and Arena-Hard at 8{,}502 and 8{,}904 GPU-seconds (Table~\ref{tab:gpu-time}). The 75\% pilot uses less GPU time per token than the 50\% pilot on these benchmarks, and both use more than their matched AR reference.

\paragraph*{Summarization throughput}
On CNN/DailyMail, EAGLE-3 has the highest observed mean OTPS at 97.98 tokens/s, followed by DFlash at 97.67 tokens/s (Table~\ref{tab:otps}). The 75\% and 50\% pilots deliver 34.48 and 33.25 tokens/s, respectively, compared with 39.39 tokens/s for their matched AR reference.

\paragraph*{Summarization GPU time}
EAGLE-3 has the lowest observed mean GPU time per million output tokens on CNN/DailyMail at 10{,}210 GPU-seconds, followed by DFlash at 10{,}240 GPU-seconds (Table~\ref{tab:gpu-time}). Both narrowed pilots use more GPU time per output token than their matched AR reference.

\subsection{Interpretation of the initial comparison}
Across the five initial configurations, DFlash has the highest observed mean OTPS on six of the ten benchmarks in Table~\ref{tab:otps}. P-EAGLE leads on MATH-500 and AIME 2025. EAGLE-3 leads on MT-Bench and CNN/DailyMail. These rankings describe the recorded implementations and selected prompts. The three repetitions quantify execution variation on the same fixed prompt sample.

Absolute throughput and paired speedup answer different questions. On AIME 2025, P-EAGLE delivers approximately 158 tokens/s and DFlash delivers approximately 155 tokens/s. Their paired speedups are approximately $3.41\times$ and $3.57\times$, respectively. P-EAGLE uses the faster vLLM ModelRunnerV1 AR reference in its ratio. Each speedup therefore retains its own reference, and OTPS supports the direct comparison of observed output delivery.

DFlash's mean acceptance length exceeds EAGLE-3's on seven benchmarks. EAGLE-3 has the higher value on MT-Bench, Alpaca, and CNN/DailyMail. Throughput reflects both acceptance length and execution cost. For example, DFlash has slightly higher OTPS on Alpaca even though EAGLE-3 retains more tokens per verification round. The drafting and verification costs also contribute to elapsed time. Complete diagnostic observations can quantify those costs where validated evidence is available. The report leaves unavailable component measurements empty.

The output-length differences materially affect the interpretation of speedup. On the third GSM8K prompt, EAGLE-3 returns 329 tokens and stops at a terminal token. Its SGLang AR reference reaches the 2,048-token limit. Their three-prompt blocks therefore contain 596 and 2,315 delivered tokens, respectively.\footnote{The saved records are \path{eagle3/01-gsm8k.json} and \path{ar-sglang/01-gsm8k.json}, relative to \path{measurements/trial-first30-20261122}. Their SHA-256 prefixes are \texttt{39aa3d823455e6d2} and \texttt{30e40dc9e6cd872b}. The export's \path{measurement-evidence.json} preserves the full source hashes.} OTPS describes the delivered tokens and elapsed time of these observed responses. An identical-response acceleration claim requires resolving the response differences.

The two AR serving paths also produce different responses. SGLang AR and vLLM ModelRunnerV1 AR match exactly on 11 of the 30 distinct prompts. Their outputs remain stable across the three repetitions.\footnote{The comparison pairs the delivered token IDs, lengths, and stopping reasons in the ten \texttt{01-*} benchmark files under \path{ar-sglang} and \path{ar-vllm-v1}. The corresponding \texttt{02-*} and \texttt{03-*} files establish repetition stability. All paths are relative to the same trial directory, and the export preserves their hashes.} Each speculative method therefore retains its matched serving reference. The AR difference also motivates checking the target's serving and numerical execution paths. The current records establish the observed differences. Their cause remains unresolved.

\subsection{Pilot results and implications}
The width comparison and the AR comparison have different conclusions. Among the observed width means, 100\% has the highest OTPS on six benchmarks. The 75\% width leads on MT-Bench, Alpaca, CNN/DailyMail, and AIME 2025. The 50\% width leads on none. On GSM8K, 100\% and 75\% differ by less than 0.01 tokens/s and both display as 39.02. The table preserves that displayed tie.

Relative to the observed 100\% means, the 75\% width delivers 94.0\%--101.0\% of its throughput across benchmarks. The 50\% width delivers 96.3\%--98.9\%. The timed GPU-cost comparison has the same per-benchmark leaders. Within the same-device comparison of the narrowed widths, 75\% leads on seven benchmarks. The 50\% width leads on MATH-500, HumanEval, and MBPP. These observations provide no consistent throughput or compute-cost improvement from reducing width.

Both narrowed pilots have a complete matching AR reference. The 75\% width achieves $0.798\times$--$0.971\times$ the AR throughput, and the 50\% width achieves $0.807\times$--$0.936\times$. Each pilot is slower than AR on all ten benchmarks. Their normalized GPU time per output token is correspondingly higher. This result answers whether the narrowed pilots accelerate ordinary generation under the recorded procedure.

The pilot acceptance lengths remain approximately 1.1--1.25 committed tokens per verification round, despite proposing eight draft tokens. Most committed progress therefore comes from the retained target token, and accepted drafts add little progress per round. That limited acceptance helps explain why drafting offers little opportunity to recover its execution cost. The saved counts establish the acceptance result. Detailed pilot component timing remains unavailable.

Both narrowed pilots match their AR reference on 12 of 30 distinct prompts, or 36 of 90 repeated executions. The remaining output differences constrain an exact-preservation claim. The checkpoints' small training set and unresolved original training details also limit generalization. The present evidence supports improving and verifying the full-width drafter before selecting a narrower model for a larger study.

\section{Resource accounting}
Table~\ref{tab:allocations} reports the two evaluation allocations once each. The initial configurations and original 100\% pilot share job \texttt{14056875}. The narrowed pilots and active Transformers AR reference use job \texttt{14058516}. Their allocated GPU-hours and CPU-hours include model loading, warmups, qualification, idle periods, inference, diagnostics, and incomplete attempts within each job. The separate project ledger preserves preparation allocations. Output-token ratios for incomplete whole-job evidence remain empty.

Method-by-benchmark compute cost uses each benchmark's measured elapsed time and the GPU count from its matching final Slurm record. Allocated GPU time per million output tokens equals that timed interval multiplied by the GPU count, divided by delivered tokens, and multiplied by one million. Output throughput per GPU divides delivered tokens by allocated GPU-seconds. The interval includes client gaps within the block. Startup, warmups, and time outside the block remain in the whole-job accounting scope.

Per-prompt GPU-seconds multiply the allocated GPU count by the prompt's end-to-end completion time in seconds. With one allocated GPU, a response that takes two seconds consumes two GPU-seconds during that request. The request-latency tables report milliseconds, so dividing those values by 1,000 gives the corresponding per-prompt GPU-seconds. Peak memory provides a separate measure of the capacity occupied during the recorded interval.

Memory observations retain separate scopes. Whole-device samples describe memory used on the physical GPU. Framework-allocated memory describes live tensors, and framework-reserved memory describes the framework's retained memory pool. Diagnostic peaks require a reset scope covering the measured requests. Whole-device sampling at 100\,ms can miss shorter peaks. The tables identify the memory measure and its unit explicitly.

The continuation's passive sampler records 5,647 observations from 01:41:44.425 to 01:51:09.615 UTC on November 23, 2026. Sampling starts after the 75\% and 50\% runs and during AR evaluation. Their two pilot memory rows therefore remain empty. AR memory cells contain two or three eligible repetitions, depending on benchmark coverage. The original 100\% run has its own matching sampler log. A sampled memory cell requires that its matching log span the entire timed benchmark interval.

\section{Limitations and next experiments}
The sample contains three fixed prompts per benchmark. Each run uses one GPU, one active request, greedy decoding, and a fixed output limit. Each released-method result reflects its recorded local implementation. Each team result identifies its checkpoint origin and available training history. Benchmark task accuracy remains unmeasured. Component times describe the saved instrumented executions, including their synchronization overhead. Partial diagnostic coverage limits comparisons for methods with fewer eligible repetitions.

Response differences remain a limitation of the speed comparison. An exact-preservation claim requires an explanation of the divergent outputs or a corrected implementation that passes the check. Saved streams and score probes preserve evidence for that future investigation.

The next research phase targets higher speed and lower compute cost than ordinary generation and the released methods. A stronger, verified full-width pilot provides the starting point. The 75\% width is the more promising narrowed candidate in the present observations because it leads 50\% on seven benchmarks. A controlled conditioning experiment can compare full width with 75\% while holding training data, updates, seeds, target revision, hardware, and evaluation settings fixed. Further architecture experiments can vary the drafter's conditioning, layer count, and proposal length. Each experiment should measure whether improved acceptance offsets drafting and verification costs. The same OTPS, GPU-time, and response-preservation checks support comparison with matched AR, EAGLE-3, DFlash, and P-EAGLE. Repeated training seeds and a larger held-out evaluation can then estimate variability and generalization.

\begin{thebibliography}{8}
\bibitem{leviathan2023}
Y. Leviathan, M. Kalman, and Y. Matias, ``Fast inference from transformers via speculative decoding,'' in \emph{Proc. 40th Int. Conf. Machine Learning}, PMLR, vol.~202, pp.~19274--19286, 2023. \url{https://proceedings.mlr.press/v202/leviathan23a.html}.
\bibitem{eagle3}
Y. Li, F. Wei, C. Zhang, and H. Zhang, ``EAGLE-3: Scaling up inference acceleration of large language models via training-time test,'' arXiv:2503.01840, 2025. \url{https://arxiv.org/abs/2503.01840}.
\bibitem{peagle}
M. Hui \emph{et al.}, ``P-EAGLE: Parallel-drafting EAGLE with scalable training,'' arXiv:2602.01469, 2026. \url{https://arxiv.org/abs/2602.01469}.
\bibitem{dflash}
J. Chen, Y. Liang, and Z. Liu, ``DFlash: Block diffusion for flash speculative decoding,'' arXiv:2602.06036, 2026. \url{https://arxiv.org/abs/2602.06036}.
\bibitem{hspec}
W. Jiang \emph{et al.}, ``H-Spec: Parallel speculative decoding without a drafter-side KV cache,'' arXiv:2609.24197, 2026. \url{https://arxiv.org/abs/2609.24197}.
\bibitem{sglangsource}
SGLang contributors, ``Speculative decoding result processing,'' source revision \texttt{0bcd822377da}, 2026. \url{https://github.com/sgl-project/sglang/blob/0bcd822377da7b5718e674eaf9c870d349424dd1/python/sglang/srt/managers/scheduler_components/batch_result_processor.py}.
\end{thebibliography}

\clearpage
\appendices
\section{Metric definitions}
Each table uses the following measurement definition. Repetition estimates are summarized by their mean and sample standard deviation.

\subsection{Output tokens per second (OTPS)}

\textit{Unit: output tokens/s.} Add the output tokens that all requests receive, including partial responses. Divide the total by the measured interval through the final completion or timeout. Count output tokens after the stopping and output-length rules apply.

\subsection{Allocated GPU time per million output tokens}

\textit{Unit: GPU-seconds/million output tokens.} Multiply the benchmark's elapsed seconds by the number of allocated GPUs. Divide by its delivered output tokens and multiply by one million. The comparison tables cover each timed benchmark block. Full-job accounting also includes model loading, warmups and idle time outside those blocks.

\subsection{Exact output agreement with AR}

\textit{Unit: \% of checked requests.} Under greedy decoding, compare each method with the target model generating alone in the same engine. Match the prompt and generation settings. Count requests with identical output token IDs, output lengths and stopping reasons. Divide by all checked requests and multiply by 100.

\subsection{Speedup over AR}

\textit{Unit: \(\times\).} Divide the method's OTPS by the OTPS of autoregressive decoding in the matching engine. Match the prompts, output rules, hardware and request schedule. A ratio of 1.2 means the method delivers 20\% more output per second.

\subsection{End-to-end request latency}

\textit{Unit: ms/request, median and p95.} Time each successful request from actual submission until completion. Sort those times and report the median and the 95th percentile (p95). The median describes the middle request. The p95 describes the time within which roughly 95\% of successful requests finish.

\subsection{Time to first token (TTFT)}

\textit{Unit: ms/request, median and p95.} Time each request from actual client submission until the first streamed output arrives. Use the first nonempty output chunk as the client observation. Exclude metadata-only events. Report the median and p95 among requests that return output, and count requests that return none.

\subsection{Maximum inter-chunk latency per request}

\textit{Unit: ms, p95 across eligible requests.} A streaming server sends an answer in pieces called chunks. For each request with at least two nonempty chunks, find the longest gap between successive chunks. Take the p95 of those longest gaps, giving each eligible request one observation.

\subsection{Peak GPU memory}

\textit{Unit: GiB/device.} Record the highest observed memory use on each GPU during the stated interval. Keep separate readings for memory used on the whole device, memory occupied by live tensors and memory reserved by the framework. A gibibyte (GiB) equals 1,073,741,824 bytes.

\subsection{Output throughput per GPU}

\textit{Unit: output tokens/s/GPU.} Divide a benchmark's delivered output tokens by its elapsed seconds multiplied by the allocated GPU count. The comparison tables use timed benchmark blocks. Complete-job resource totals retain their own allocation interval.

\subsection{End-to-end output rate per request}

\textit{Unit: output tokens/s/request.} For each completed request, divide its returned output tokens by the seconds from actual submission through completion. Average the request rates equally. The clock includes waiting, prompt processing and generation.

\subsection{Time per output token (TPOT)}

\textit{Unit: ms/output token.} For each completed response with N output tokens, require N greater than one. Subtract TTFT from end-to-end request latency and divide by N minus one. Average the eligible request values equally. Count excluded responses separately.

\subsection{Inter-chunk latency}

\textit{Unit: ms/gap, p95.} Within each response, measure the gaps between successive nonempty chunks. Pool all observed gaps and report their p95. A response with more chunks contributes more gaps.

\subsection{Draft-token acceptance rate}

\textit{Unit: \% of proposed draft tokens.} Divide accepted draft tokens by proposed draft tokens and multiply by 100. The numerator and denominator count draft tokens only. Record target correction and bonus tokens separately. For a branching draft, state which branches the proposal count covers.

\subsection{Accepted draft tokens per round}

\textit{Unit: accepted draft tokens/round.} Divide the total accepted draft tokens by the number of verification rounds. Include rounds that accept zero draft tokens. Record tokens supplied by the target separately.

\subsection{Mean acceptance length (\(\tau\))}

\textit{Unit: committed tokens/round.} Add the tokens retained across all verification rounds before final output trimming. Include accepted draft tokens and any target correction or bonus tokens that the decoder retains. Divide by the total round count, giving each round equal weight. Record initial tokens generated outside verification rounds separately.

\subsection{Mean acceptance length (\(\tau\)), averaged across requests}

\textit{Unit: committed tokens/round.} For each request with at least one verification round, divide its retained tokens by its round count. Include accepted drafts and retained target tokens before final output trimming. Average the request values equally.

\subsection{Delivered output tokens per round}

\textit{Unit: delivered tokens/round.} Count the output tokens from the recorded rounds that actually reach the client after stopping and output limits. Divide by all those rounds, including rounds that return zero tokens. Record any initial target token outside the rounds separately.

\subsection{Prefix acceptance probability by position}

\textit{Unit: probability at each draft position.} For each draft position, count the rounds that proposed that position. Among those rounds, count how often the target accepted every draft token through that position before final output clipping. Divide the second count by the first.

\subsection{Conditional acceptance rate by position}

\textit{Unit: \% at each draft position.} At each draft position, count rounds that proposed the position and accepted every earlier proposal. Among these eligible rounds, count how often the current proposal passes. Divide passes by eligible rounds and multiply by 100. Report a result only when the eligible count is positive.

\subsection{Draft latency, verification latency and round latency}

\textit{Unit: ms/round for each component.} In a separate timing trace, measure a full cycle from drafting through verification and the final state update. Measure the time spent drafting, checking and moving data within that cycle. Record which operations overlap and average times over the observed cycles.
\section{Evidence availability}
Each empty cell has a reason in the accompanying primary or team measurement-evidence JSON record. Those records also identify the source files, hashes, repetition values, and eligible observations.

\paragraph*{Exact output agreement with AR}

50 of 60 cells contain validated measurements.

Remaining cells require completed benchmark records that pass the metric's checks. Streaming metrics also require eligible output chunks.

\paragraph*{Speedup over AR}

50 of 60 cells contain validated measurements.

Remaining cells require completed benchmark records that pass the metric's checks. Streaming metrics also require eligible output chunks.

\paragraph*{Peak GPU memory; Sampled whole-device memory}

70 of 90 cells contain validated measurements.

Remaining cells require completed benchmark records that pass the metric's checks. Streaming metrics also require eligible output chunks.

\paragraph*{Peak GPU memory; Framework-allocated memory}

14 of 90 cells contain validated measurements.

Each value requires a synchronized framework memory peak with a reset scope covering all three diagnostic requests.

\paragraph*{Peak GPU memory; Framework-reserved memory}

14 of 90 cells contain validated measurements.

Each value requires a synchronized framework memory peak with a reset scope covering all three diagnostic requests.

\paragraph*{Mean acceptance length (\(\tau\))}

50 of 60 cells contain validated measurements.

Each value requires validated native counts or a matching complete diagnostic trace.

\paragraph*{Mean acceptance length (\(\tau\)), averaged across requests}

50 of 60 cells contain validated measurements.

Each value requires validated native counts or a matching complete diagnostic trace.

\paragraph*{Delivered output tokens per round}

50 of 60 cells contain validated measurements.

Each value requires validated native counts or a matching complete diagnostic trace.

\paragraph*{Prefix acceptance probability by position}

330 of 540 cells contain validated measurements.

Each position requires a validated proposal depth and a positive eligible-round count. Methods with shorter proposals leave later positions empty. Tree proposals require an explicit path interpretation.

\paragraph*{Prefix acceptance probability by position; Tree-prefix diagnostic}

30 of 90 cells contain validated measurements.

Each position requires a validated proposal depth and a positive eligible-round count. Methods with shorter proposals leave later positions empty. Tree proposals require an explicit path interpretation.

\paragraph*{Conditional acceptance rate by position}

166 of 540 cells contain validated measurements.

Each position requires a validated proposal depth and a positive eligible-round count. Methods with shorter proposals leave later positions empty. Tree proposals require an explicit path interpretation.

\paragraph*{Conditional acceptance rate by position; Tree-prefix diagnostic}

30 of 90 cells contain validated measurements.

Each position requires a validated proposal depth and a positive eligible-round count. Methods with shorter proposals leave later positions empty. Tree proposals require an explicit path interpretation.

\paragraph*{Draft latency, verification latency and round latency; Round latency}

14 of 60 cells contain validated measurements.

The available raw worker intervals require validation against complete verification rounds, including the final state update.

\paragraph*{Draft latency, verification latency and round latency; Draft latency}

14 of 60 cells contain validated measurements.

The available raw worker intervals require validation against complete verification rounds, including the final state update.

\paragraph*{Draft latency, verification latency and round latency; Verification latency}

14 of 60 cells contain validated measurements.

The available raw worker intervals require validation against complete verification rounds, including the final state update.
\section{Additional measured results}
The following tables preserve the method-by-benchmark structure. Each table reports one metric, variant, and draft position where applicable. Tables appear when at least one measured cell exists. An empty cell within a measured table retains its evidence reason in the export.
% Generated appendix tables contain measured observations only.
\begin{table*}[!htp]
\centering\small
\setlength{\tabcolsep}{3pt}
\caption{Exact output agreement with AR (\% of checked requests).}
\label{tab:appendix-metric-quality-none}
\textit{Mathematics}\par\smallskip
\begin{tabular}{@{}>{\raggedright\arraybackslash}p{3.7cm}rrr@{}}
\toprule
Method & GSM8K & MATH-500 & AIME 2025 \\
\midrule
EAGLE-3 & \(\mathbf{66.67} \pm 0\;[3]\) & \(0 \pm 0\;[3]\) & \(0 \pm 0\;[3]\) \\
P-EAGLE & \(\mathbf{66.67} \pm 0\;[3]\) & \(\mathbf{33.33} \pm 0\;[3]\) & \(\mathbf{33.33} \pm 0\;[3]\) \\
DFlash & \(\mathbf{66.67} \pm 0\;[3]\) & \(\mathbf{33.33} \pm 0\;[3]\) & \(0 \pm 0\;[3]\) \\
Pilot FFN 75\% & \(33.33 \pm 0\;[3]\) & \(\mathbf{33.33} \pm 0\;[3]\) & \(0 \pm 0\;[3]\) \\
Pilot FFN 50\% & \(33.33 \pm 0\;[3]\) & \(\mathbf{33.33} \pm 0\;[3]\) & \(0 \pm 0\;[3]\) \\
\bottomrule
\end{tabular}\par\medskip
\textit{Code}\par\smallskip
\begin{tabular}{@{}>{\raggedright\arraybackslash}p{3.7cm}rrr@{}}
\toprule
Method & HumanEval & MBPP & LiveCodeBench \\
\midrule
EAGLE-3 & \(66.67 \pm 0\;[3]\) & \(66.67 \pm 0\;[3]\) & \(\mathbf{33.33} \pm 0\;[3]\) \\
P-EAGLE & \(66.67 \pm 0\;[3]\) & \(66.67 \pm 0\;[3]\) & \(\mathbf{33.33} \pm 0\;[3]\) \\
DFlash & \(66.67 \pm 0\;[3]\) & \(66.67 \pm 0\;[3]\) & \(\mathbf{33.33} \pm 0\;[3]\) \\
Pilot FFN 75\% & \(\mathbf{100} \pm 0\;[3]\) & \(\mathbf{100} \pm 0\;[3]\) & \(\mathbf{33.33} \pm 0\;[3]\) \\
Pilot FFN 50\% & \(\mathbf{100} \pm 0\;[3]\) & \(\mathbf{100} \pm 0\;[3]\) & \(\mathbf{33.33} \pm 0\;[3]\) \\
\bottomrule
\end{tabular}\par\medskip
\textit{Conversation, instruction following, and summarization}\par\smallskip
\begin{tabular}{@{}>{\raggedright\arraybackslash}p{3.7cm}rrrr@{}}
\toprule
Method & MT-Bench & Alpaca & Arena-Hard & CNN/DailyMail \\
\midrule
EAGLE-3 & \(33.33 \pm 0\;[3]\) & \(33.33 \pm 0\;[3]\) & \(\mathbf{66.67} \pm 0\;[3]\) & \(0 \pm 0\;[3]\) \\
P-EAGLE & \(\mathbf{66.67} \pm 0\;[3]\) & \(33.33 \pm 0\;[3]\) & \(33.33 \pm 0\;[3]\) & \(33.33 \pm 0\;[3]\) \\
DFlash & \(33.33 \pm 0\;[3]\) & \(33.33 \pm 0\;[3]\) & \(33.33 \pm 0\;[3]\) & \(\mathbf{66.67} \pm 0\;[3]\) \\
Pilot FFN 75\% & \(0 \pm 0\;[3]\) & \(\mathbf{66.67} \pm 0\;[3]\) & \(33.33 \pm 0\;[3]\) & \(0 \pm 0\;[3]\) \\
Pilot FFN 50\% & \(0 \pm 0\;[3]\) & \(\mathbf{66.67} \pm 0\;[3]\) & \(33.33 \pm 0\;[3]\) & \(0 \pm 0\;[3]\) \\
\bottomrule
\end{tabular}\par\medskip
\begin{minipage}{\linewidth}\footnotesize
Under greedy decoding, compare each method with the target model generating alone in the same engine. Match the prompt and generation settings. Count requests with identical output token IDs, output lengths and stopping reasons. Divide by all checked requests and multiply by 100.
The bold mean is the best observed value in each benchmark column. The comparison covers the recorded implementations under their stated configurations. At least two eligible method entries are required. Ties at the displayed precision share bold formatting. Bold formatting conveys an observed comparison and makes no statistical-significance claim. Brackets show the available repetition count. Response differences remain part of the interpretation. Higher values are preferred.
Brackets give eligible repetition counts. Standard deviation appears when at least two repetitions are available. Each repetition uses three prompts. The evidence records preserve observation counts and measurement scopes. Empty cells lack an eligible measured value.
\end{minipage}
\end{table*}

\begin{table*}[!htp]
\centering\small
\setlength{\tabcolsep}{3pt}
\caption{End-to-end request latency; Median (ms/request).}
\label{tab:appendix-metric-completion-median-none}
\textit{Mathematics}\par\smallskip
\begin{tabular}{@{}>{\raggedright\arraybackslash}p{3.7cm}rrr@{}}
\toprule
Method & GSM8K & MATH-500 & AIME 2025 \\
\midrule
EAGLE-3 & \(1375 \pm 32\;[3]\) & \(4110 \pm 6.9\;[3]\) & \(\mathbf{8338} \pm 40\;[3]\) \\
P-EAGLE & \(1328 \pm 0.18\;[3]\) & \(\mathbf{2833} \pm 0.16\;[3]\) & \(8932 \pm 0.37\;[3]\) \\
DFlash & \(\mathbf{1097} \pm 12\;[3]\) & \(3010 \pm 8.9\;[3]\) & \(11320 \pm 11\;[3]\) \\
AR \textperiodcentered{} SGLang & \(3321 \pm 36\;[3]\) & \(10040 \pm 9.8\;[3]\) & \(21160 \pm 11\;[3]\) \\
AR \textperiodcentered{} vLLM ModelRunnerV1 & \(4511 \pm 0.21\;[3]\) & \(11120 \pm 0.23\;[3]\) & \(18350 \pm 0.97\;[3]\) \\
Pilot FFN 100\% & \(3992 \pm 9.4\;[3]\) & \(12900 \pm 7.8\;[3]\) & \(62870 \pm 70\;[3]\) \\
Pilot FFN 75\% & \(4100 \pm 12\;[3]\) & \(12970 \pm 14\;[3]\) & \(62110 \pm 230\;[3]\) \\
Pilot FFN 50\% & \(4101 \pm 28\;[3]\) & \(13080 \pm 180\;[3]\) & \(64800 \pm 160\;[3]\) \\
AR \textperiodcentered{} Transformers & \(3596 \pm 0.82\;[3]\) & \(11070 \pm 5.2\;[3]\) & \(19430 \pm 42\;[3]\) \\
\bottomrule
\end{tabular}\par\medskip
\textit{Code}\par\smallskip
\begin{tabular}{@{}>{\raggedright\arraybackslash}p{3.7cm}rrr@{}}
\toprule
Method & HumanEval & MBPP & LiveCodeBench \\
\midrule
EAGLE-3 & \(833.6 \pm 1.1\;[3]\) & \(1018 \pm 12\;[3]\) & \(2522 \pm 9.1\;[3]\) \\
P-EAGLE & \(835 \pm 0.27\;[3]\) & \(782.9 \pm 0.24\;[3]\) & \(2128 \pm 0.072\;[3]\) \\
DFlash & \(\mathbf{703.9} \pm 7.4\;[3]\) & \(\mathbf{699.8} \pm 7.1\;[3]\) & \(\mathbf{1740} \pm 4.2\;[3]\) \\
AR \textperiodcentered{} SGLang & \(2422 \pm 3.2\;[3]\) & \(2556 \pm 2.4\;[3]\) & \(7439 \pm 11\;[3]\) \\
AR \textperiodcentered{} vLLM ModelRunnerV1 & \(2277 \pm 0.075\;[3]\) & \(2383 \pm 7.1\;[3]\) & \(6347 \pm 6.9\;[3]\) \\
Pilot FFN 100\% & \(2848 \pm 3.6\;[3]\) & \(3089 \pm 6.3\;[3]\) & \(7912 \pm 70\;[3]\) \\
Pilot FFN 75\% & \(3012 \pm 11\;[3]\) & \(3224 \pm 9.6\;[3]\) & \(8040 \pm 14\;[3]\) \\
Pilot FFN 50\% & \(2929 \pm 18\;[3]\) & \(3129 \pm 9.2\;[3]\) & \(8045 \pm 19\;[3]\) \\
AR \textperiodcentered{} Transformers & \(2614 \pm 3.1\;[3]\) & \(2751 \pm 7.4\;[3]\) & \(4959 \pm 9.9\;[3]\) \\
\bottomrule
\end{tabular}\par\medskip
\textit{Conversation, instruction following, and summarization}\par\smallskip
\begin{tabular}{@{}>{\raggedright\arraybackslash}p{3.7cm}rrrr@{}}
\toprule
Method & MT-Bench & Alpaca & Arena-Hard & CNN/DailyMail \\
\midrule
EAGLE-3 & \(1921 \pm 38\;[3]\) & \(2035 \pm 30\;[3]\) & \(4457 \pm 40\;[3]\) & \(2582 \pm 36\;[3]\) \\
P-EAGLE & \(2384 \pm 0.28\;[3]\) & \(1816 \pm 0.32\;[3]\) & \(4770 \pm 15\;[3]\) & \(2810 \pm 1.2\;[3]\) \\
DFlash & \(\mathbf{1851} \pm 9.1\;[3]\) & \(\mathbf{1679} \pm 5\;[3]\) & \(\mathbf{3927} \pm 13\;[3]\) & \(\mathbf{2491} \pm 4\;[3]\) \\
AR \textperiodcentered{} SGLang & \(4871 \pm 3.3\;[3]\) & \(5061 \pm 17\;[3]\) & \(8800 \pm 6\;[3]\) & \(5773 \pm 9.3\;[3]\) \\
AR \textperiodcentered{} vLLM ModelRunnerV1 & \(4455 \pm 0.48\;[3]\) & \(5032 \pm 0.25\;[3]\) & \(7315 \pm 2.4\;[3]\) & \(5064 \pm 0.17\;[3]\) \\
Pilot FFN 100\% & \(5780 \pm 5.1\;[3]\) & \(5874 \pm 5\;[3]\) & \(11500 \pm 21\;[3]\) & \(7236 \pm 25\;[3]\) \\
Pilot FFN 75\% & \(5770 \pm 17\;[3]\) & \(5854 \pm 19\;[3]\) & \(11790 \pm 5.5\;[3]\) & \(7042 \pm 12\;[3]\) \\
Pilot FFN 50\% & \(5909 \pm 56\;[3]\) & \(6016 \pm 64\;[3]\) & \(11740 \pm 15\;[3]\) & \(7329 \pm 11\;[3]\) \\
AR \textperiodcentered{} Transformers & \(5203 \pm 16\;[3]\) & \(5490 \pm 7.9\;[3]\) & \(8545 \pm 7\;[3]\) & \(5903 \pm 8\;[3]\) \\
\bottomrule
\end{tabular}\par\medskip
\begin{minipage}{\linewidth}\footnotesize
Time each successful request from actual submission until completion. Sort those times and report the median and the 95th percentile (p95). The median describes the middle request. The p95 describes the time within which roughly 95\% of successful requests finish.
The bold mean is the best observed value in each benchmark column. The comparison covers the recorded implementations under their stated configurations. At least two eligible method entries are required. Ties at the displayed precision share bold formatting. Bold formatting conveys an observed comparison and makes no statistical-significance claim. Brackets show the available repetition count. Response differences remain part of the interpretation. Lower values are preferred.
Brackets give eligible repetition counts. Standard deviation appears when at least two repetitions are available. Each repetition uses three prompts. The evidence records preserve observation counts and measurement scopes. Empty cells lack an eligible measured value.
\end{minipage}
\end{table*}

\begin{table*}[!htp]
\centering\small
\setlength{\tabcolsep}{3pt}
\caption{End-to-end request latency; 95th percentile (ms/request).}
\label{tab:appendix-metric-completion-95th-percentile-none}
\textit{Mathematics}\par\smallskip
\begin{tabular}{@{}>{\raggedright\arraybackslash}p{3.7cm}rrr@{}}
\toprule
Method & GSM8K & MATH-500 & AIME 2025 \\
\midrule
EAGLE-3 & \(2919 \pm 3.6\;[3]\) & \(5010 \pm 4\;[3]\) & \(17600 \pm 100\;[3]\) \\
P-EAGLE & \(\mathbf{2323} \pm 31\;[3]\) & \(4016 \pm 0.5\;[3]\) & \(12050 \pm 0.71\;[3]\) \\
DFlash & \(8009 \pm 2.1\;[3]\) & \(\mathbf{3811} \pm 14\;[3]\) & \(\mathbf{12000} \pm 4.5\;[3]\) \\
AR \textperiodcentered{} SGLang & \(42930 \pm 11\;[3]\) & \(14430 \pm 8\;[3]\) & \(44660 \pm 46\;[3]\) \\
AR \textperiodcentered{} vLLM ModelRunnerV1 & \(6273 \pm 0.036\;[3]\) & \(13670 \pm 0.86\;[3]\) & \(41650 \pm 0.48\;[3]\) \\
Pilot FFN 100\% & \(6866 \pm 7.4\;[3]\) & \(16180 \pm 44\;[3]\) & \(63710 \pm 54\;[3]\) \\
Pilot FFN 75\% & \(6804 \pm 8.7\;[3]\) & \(18640 \pm 3500\;[3]\) & \(62520 \pm 100\;[3]\) \\
Pilot FFN 50\% & \(7148 \pm 20\;[3]\) & \(16950 \pm 110\;[3]\) & \(64930 \pm 66\;[3]\) \\
AR \textperiodcentered{} Transformers & \(5390 \pm 39\;[3]\) & \(14650 \pm 2.9\;[3]\) & \(49570 \pm 73\;[3]\) \\
\bottomrule
\end{tabular}\par\medskip
\textit{Code}\par\smallskip
\begin{tabular}{@{}>{\raggedright\arraybackslash}p{3.7cm}rrr@{}}
\toprule
Method & HumanEval & MBPP & LiveCodeBench \\
\midrule
EAGLE-3 & \(1009 \pm 7.7\;[3]\) & \(2078 \pm 25\;[3]\) & \(3802 \pm 12\;[3]\) \\
P-EAGLE & \(839.5 \pm 0.17\;[3]\) & \(1922 \pm 0.15\;[3]\) & \(\mathbf{3233} \pm 0.55\;[3]\) \\
DFlash & \(\mathbf{758.9} \pm 9.8\;[3]\) & \(\mathbf{1368} \pm 2.5\;[3]\) & \(3555 \pm 11\;[3]\) \\
AR \textperiodcentered{} SGLang & \(2672 \pm 9.2\;[3]\) & \(5496 \pm 5\;[3]\) & \(7719 \pm 17\;[3]\) \\
AR \textperiodcentered{} vLLM ModelRunnerV1 & \(2493 \pm 0.33\;[3]\) & \(4997 \pm 0.3\;[3]\) & \(7694 \pm 1.1\;[3]\) \\
Pilot FFN 100\% & \(3139 \pm 2.6\;[3]\) & \(6370 \pm 11\;[3]\) & \(9136 \pm 33\;[3]\) \\
Pilot FFN 75\% & \(3136 \pm 12\;[3]\) & \(6777 \pm 8.4\;[3]\) & \(9235 \pm 34\;[3]\) \\
Pilot FFN 50\% & \(3175 \pm 15\;[3]\) & \(6447 \pm 18\;[3]\) & \(9495 \pm 19\;[3]\) \\
AR \textperiodcentered{} Transformers & \(2885 \pm 10\;[3]\) & \(5955 \pm 8.3\;[3]\) & \(7683 \pm 23\;[3]\) \\
\bottomrule
\end{tabular}\par\medskip
\textit{Conversation, instruction following, and summarization}\par\smallskip
\begin{tabular}{@{}>{\raggedright\arraybackslash}p{3.7cm}rrrr@{}}
\toprule
Method & MT-Bench & Alpaca & Arena-Hard & CNN/DailyMail \\
\midrule
EAGLE-3 & \(\mathbf{7080} \pm 44\;[3]\) & \(\mathbf{3570} \pm 9.6\;[3]\) & \(4809 \pm 44\;[3]\) & \(2816 \pm 28\;[3]\) \\
P-EAGLE & \(11000 \pm 0.39\;[3]\) & \(3821 \pm 0.32\;[3]\) & \(5057 \pm 1.6\;[3]\) & \(3210 \pm 0.32\;[3]\) \\
DFlash & \(7335 \pm 17\;[3]\) & \(3731 \pm 14\;[3]\) & \(\mathbf{4296} \pm 6.1\;[3]\) & \(\mathbf{2550} \pm 12\;[3]\) \\
AR \textperiodcentered{} SGLang & \(16800 \pm 32\;[3]\) & \(9296 \pm 10\;[3]\) & \(11290 \pm 10\;[3]\) & \(5931 \pm 19\;[3]\) \\
AR \textperiodcentered{} vLLM ModelRunnerV1 & \(15720 \pm 4.1\;[3]\) & \(7841 \pm 0.29\;[3]\) & \(10970 \pm 1.1\;[3]\) & \(5410 \pm 0.72\;[3]\) \\
Pilot FFN 100\% & \(20550 \pm 19\;[3]\) & \(11020 \pm 23\;[3]\) & \(13920 \pm 61\;[3]\) & \(7630 \pm 14\;[3]\) \\
Pilot FFN 75\% & \(20530 \pm 24\;[3]\) & \(10960 \pm 18\;[3]\) & \(14120 \pm 9.9\;[3]\) & \(7687 \pm 5\;[3]\) \\
Pilot FFN 50\% & \(20770 \pm 71\;[3]\) & \(11200 \pm 12\;[3]\) & \(14160 \pm 52\;[3]\) & \(8031 \pm 8.9\;[3]\) \\
AR \textperiodcentered{} Transformers & \(14890 \pm 31\;[3]\) & \(10210 \pm 8.3\;[3]\) & \(12080 \pm 2.6\;[3]\) & \(6296 \pm 3.5\;[3]\) \\
\bottomrule
\end{tabular}\par\medskip
\begin{minipage}{\linewidth}\footnotesize
Time each successful request from actual submission until completion. Sort those times and report the median and the 95th percentile (p95). The median describes the middle request. The p95 describes the time within which roughly 95\% of successful requests finish.
The bold mean is the best observed value in each benchmark column. The comparison covers the recorded implementations under their stated configurations. At least two eligible method entries are required. Ties at the displayed precision share bold formatting. Bold formatting conveys an observed comparison and makes no statistical-significance claim. Brackets show the available repetition count. Response differences remain part of the interpretation. Lower values are preferred.
Brackets give eligible repetition counts. Standard deviation appears when at least two repetitions are available. Each repetition uses three prompts. The evidence records preserve observation counts and measurement scopes. Empty cells lack an eligible measured value.
\end{minipage}
\end{table*}

\begin{table*}[!htp]
\centering\small
\setlength{\tabcolsep}{3pt}
\caption{Time to first token (TTFT); Median (ms/request).}
\label{tab:appendix-metric-ttft-median-none}
\textit{Mathematics}\par\smallskip
\begin{tabular}{@{}>{\raggedright\arraybackslash}p{3.7cm}rrr@{}}
\toprule
Method & GSM8K & MATH-500 & AIME 2025 \\
\midrule
EAGLE-3 & \(33.05 \pm 0.73\;[3]\) & \(34.53 \pm 0.16\;[3]\) & \(35.76 \pm 0.14\;[3]\) \\
P-EAGLE & \(56.65 \pm 0.088\;[3]\) & \(57.44 \pm 0.48\;[3]\) & \(59.73 \pm 0.18\;[3]\) \\
DFlash & \(31.74 \pm 0.81\;[3]\) & \(32.68 \pm 0.59\;[3]\) & \(33.96 \pm 0.57\;[3]\) \\
AR \textperiodcentered{} SGLang & \(\mathbf{30.82} \pm 0.067\;[3]\) & \(\mathbf{31.6} \pm 0.072\;[3]\) & \(\mathbf{32.82} \pm 0.48\;[3]\) \\
AR \textperiodcentered{} vLLM ModelRunnerV1 & \(48.11 \pm 0.28\;[3]\) & \(48.56 \pm 0.41\;[3]\) & \(50.38 \pm 0.15\;[3]\) \\
Pilot FFN 100\% & \(31.51 \pm 0.075\;[3]\) & \(31.86 \pm 0.54\;[3]\) & \(34.51 \pm 0.099\;[3]\) \\
Pilot FFN 75\% & \(31.51 \pm 0.16\;[3]\) & \(32.01 \pm 0.19\;[3]\) & \(34.4 \pm 0.34\;[3]\) \\
Pilot FFN 50\% & \(31.75 \pm 0.091\;[3]\) & \(32.29 \pm 0.41\;[3]\) & \(34.62 \pm 0.15\;[3]\) \\
AR \textperiodcentered{} Transformers & \(31.62 \pm 0.36\;[3]\) & \(32.17 \pm 0.064\;[3]\) & \(34.65 \pm 0.16\;[3]\) \\
\bottomrule
\end{tabular}\par\medskip
\textit{Code}\par\smallskip
\begin{tabular}{@{}>{\raggedright\arraybackslash}p{3.7cm}rrr@{}}
\toprule
Method & HumanEval & MBPP & LiveCodeBench \\
\midrule
EAGLE-3 & \(36.35 \pm 0.47\;[3]\) & \(33.21 \pm 0.56\;[3]\) & \(42.47 \pm 0.36\;[3]\) \\
P-EAGLE & \(58.23 \pm 0.18\;[3]\) & \(56.2 \pm 0.19\;[3]\) & \(59.96 \pm 0.55\;[3]\) \\
DFlash & \(34.29 \pm 0.26\;[3]\) & \(31.58 \pm 0.43\;[3]\) & \(40.84 \pm 0.28\;[3]\) \\
AR \textperiodcentered{} SGLang & \(\mathbf{33.32} \pm 0.2\;[3]\) & \(\mathbf{30.21} \pm 0.065\;[3]\) & \(\mathbf{39.13} \pm 0.62\;[3]\) \\
AR \textperiodcentered{} vLLM ModelRunnerV1 & \(49.92 \pm 0.23\;[3]\) & \(47.82 \pm 0.38\;[3]\) & \(50.83 \pm 0.15\;[3]\) \\
Pilot FFN 100\% & \(34.39 \pm 0.16\;[3]\) & \(30.86 \pm 0.096\;[3]\) & \(40.16 \pm 0.2\;[3]\) \\
Pilot FFN 75\% & \(34.38 \pm 0.1\;[3]\) & \(30.87 \pm 0.11\;[3]\) & \(39.9 \pm 0.079\;[3]\) \\
Pilot FFN 50\% & \(34.5 \pm 0.15\;[3]\) & \(31.02 \pm 0.22\;[3]\) & \(40.14 \pm 0.21\;[3]\) \\
AR \textperiodcentered{} Transformers & \(34.58 \pm 0.11\;[3]\) & \(30.99 \pm 0.15\;[3]\) & \(40.19 \pm 0.23\;[3]\) \\
\bottomrule
\end{tabular}\par\medskip
\textit{Conversation, instruction following, and summarization}\par\smallskip
\begin{tabular}{@{}>{\raggedright\arraybackslash}p{3.7cm}rrrr@{}}
\toprule
Method & MT-Bench & Alpaca & Arena-Hard & CNN/DailyMail \\
\midrule
EAGLE-3 & \(33.46 \pm 1.1\;[3]\) & \(32.39 \pm 0.42\;[3]\) & \(33.23 \pm 0.6\;[3]\) & \(62.16 \pm 0.32\;[3]\) \\
P-EAGLE & \(56.71 \pm 0.3\;[3]\) & \(55.46 \pm 0.49\;[3]\) & \(57.18 \pm 0.31\;[3]\) & \(79.96 \pm 1.1\;[3]\) \\
DFlash & \(32.01 \pm 0.16\;[3]\) & \(30.77 \pm 0.064\;[3]\) & \(31.64 \pm 0.25\;[3]\) & \(60.61 \pm 0.41\;[3]\) \\
AR \textperiodcentered{} SGLang & \(\mathbf{30.66} \pm 0.22\;[3]\) & \(29.59 \pm 0.0031\;[3]\) & \(\mathbf{30.38} \pm 0.33\;[3]\) & \(\mathbf{59.51} \pm 0.64\;[3]\) \\
AR \textperiodcentered{} vLLM ModelRunnerV1 & \(48.8 \pm 0.4\;[3]\) & \(47.1 \pm 0.19\;[3]\) & \(48.35 \pm 0.28\;[3]\) & \(69.97 \pm 0.61\;[3]\) \\
Pilot FFN 100\% & \(31.54 \pm 0.1\;[3]\) & \(29.31 \pm 0.12\;[3]\) & \(31.46 \pm 0.27\;[3]\) & \(65.81 \pm 1\;[3]\) \\
Pilot FFN 75\% & \(31.48 \pm 0.22\;[3]\) & \(\mathbf{29.29} \pm 0.24\;[3]\) & \(31.68 \pm 0.16\;[3]\) & \(64.71 \pm 0.91\;[3]\) \\
Pilot FFN 50\% & \(31.36 \pm 0.28\;[3]\) & \(29.51 \pm 0.14\;[3]\) & \(31.69 \pm 0.23\;[3]\) & \(65.26 \pm 0.48\;[3]\) \\
AR \textperiodcentered{} Transformers & \(31.59 \pm 0.17\;[3]\) & \(29.65 \pm 0.21\;[3]\) & \(31.82 \pm 0.15\;[3]\) & \(65.56 \pm 0.39\;[3]\) \\
\bottomrule
\end{tabular}\par\medskip
\begin{minipage}{\linewidth}\footnotesize
Time each request from actual client submission until the first streamed output arrives. Use the first nonempty output chunk as the client observation. Exclude metadata-only events. Report the median and p95 among requests that return output, and count requests that return none.
The bold mean is the best observed value in each benchmark column. The comparison covers the recorded implementations under their stated configurations. At least two eligible method entries are required. Ties at the displayed precision share bold formatting. Bold formatting conveys an observed comparison and makes no statistical-significance claim. Brackets show the available repetition count. Response differences remain part of the interpretation. Lower values are preferred.
Brackets give eligible repetition counts. Standard deviation appears when at least two repetitions are available. Each repetition uses three prompts. The evidence records preserve observation counts and measurement scopes. Empty cells lack an eligible measured value.
\end{minipage}
\end{table*}

\begin{table*}[!htp]
\centering\small
\setlength{\tabcolsep}{3pt}
\caption{Time to first token (TTFT); 95th percentile (ms/request).}
\label{tab:appendix-metric-ttft-95th-percentile-none}
\textit{Mathematics}\par\smallskip
\begin{tabular}{@{}>{\raggedright\arraybackslash}p{3.7cm}rrr@{}}
\toprule
Method & GSM8K & MATH-500 & AIME 2025 \\
\midrule
EAGLE-3 & \(34.02 \pm 0.44\;[3]\) & \(35.7 \pm 0.089\;[3]\) & \(59.57 \pm 1.2\;[3]\) \\
P-EAGLE & \(78.33 \pm 37\;[3]\) & \(59.09 \pm 0.34\;[3]\) & \(86.14 \pm 1.6\;[3]\) \\
DFlash & \(32.88 \pm 0.68\;[3]\) & \(33.84 \pm 0.41\;[3]\) & \(57.28 \pm 0.74\;[3]\) \\
AR \textperiodcentered{} SGLang & \(\mathbf{31.55} \pm 0.17\;[3]\) & \(\mathbf{33.14} \pm 0.24\;[3]\) & \(\mathbf{56.18} \pm 0.64\;[3]\) \\
AR \textperiodcentered{} vLLM ModelRunnerV1 & \(48.54 \pm 0.14\;[3]\) & \(50.35 \pm 0.34\;[3]\) & \(58.67 \pm 0.55\;[3]\) \\
Pilot FFN 100\% & \(32.33 \pm 0.18\;[3]\) & \(50.37 \pm 29\;[3]\) & \(63.01 \pm 0.93\;[3]\) \\
Pilot FFN 75\% & \(32.62 \pm 0.27\;[3]\) & \(2072 \pm 3500\;[3]\) & \(61.93 \pm 0.75\;[3]\) \\
Pilot FFN 50\% & \(32.55 \pm 0.3\;[3]\) & \(35.39 \pm 2.3\;[3]\) & \(62.48 \pm 0.38\;[3]\) \\
AR \textperiodcentered{} Transformers & \(32.36 \pm 0.1\;[3]\) & \(34.41 \pm 0.21\;[3]\) & \(62.42 \pm 0.099\;[3]\) \\
\bottomrule
\end{tabular}\par\medskip
\textit{Code}\par\smallskip
\begin{tabular}{@{}>{\raggedright\arraybackslash}p{3.7cm}rrr@{}}
\toprule
Method & HumanEval & MBPP & LiveCodeBench \\
\midrule
EAGLE-3 & \(37.18 \pm 0.56\;[3]\) & \(33.42 \pm 0.51\;[3]\) & \(42.77 \pm 0.43\;[3]\) \\
P-EAGLE & \(59.17 \pm 0.09\;[3]\) & \(56.39 \pm 0.11\;[3]\) & \(66.06 \pm 0.28\;[3]\) \\
DFlash & \(34.83 \pm 0.62\;[3]\) & \(32.05 \pm 0.47\;[3]\) & \(41.2 \pm 0.2\;[3]\) \\
AR \textperiodcentered{} SGLang & \(\mathbf{34.02} \pm 0.18\;[3]\) & \(\mathbf{30.69} \pm 0.12\;[3]\) & \(\mathbf{39.77} \pm 0.2\;[3]\) \\
AR \textperiodcentered{} vLLM ModelRunnerV1 & \(50.36 \pm 0.14\;[3]\) & \(48.02 \pm 0.56\;[3]\) & \(56.1 \pm 0.098\;[3]\) \\
Pilot FFN 100\% & \(34.95 \pm 0.24\;[3]\) & \(31.02 \pm 0.084\;[3]\) & \(44.33 \pm 0.087\;[3]\) \\
Pilot FFN 75\% & \(34.91 \pm 0.13\;[3]\) & \(31.13 \pm 0.22\;[3]\) & \(44.1 \pm 0.2\;[3]\) \\
Pilot FFN 50\% & \(35.21 \pm 0.091\;[3]\) & \(31.3 \pm 0.36\;[3]\) & \(44.21 \pm 0.13\;[3]\) \\
AR \textperiodcentered{} Transformers & \(34.69 \pm 0.17\;[3]\) & \(31.33 \pm 0.25\;[3]\) & \(44.21 \pm 0.4\;[3]\) \\
\bottomrule
\end{tabular}\par\medskip
\textit{Conversation, instruction following, and summarization}\par\smallskip
\begin{tabular}{@{}>{\raggedright\arraybackslash}p{3.7cm}rrrr@{}}
\toprule
Method & MT-Bench & Alpaca & Arena-Hard & CNN/DailyMail \\
\midrule
EAGLE-3 & \(37.79 \pm 6.8\;[3]\) & \(32.75 \pm 0.95\;[3]\) & \(79.17 \pm 0.4\;[3]\) & \(62.8 \pm 0.29\;[3]\) \\
P-EAGLE & \(97.28 \pm 70\;[3]\) & \(55.7 \pm 0.31\;[3]\) & \(95.92 \pm 18\;[3]\) & \(87.77 \pm 0.86\;[3]\) \\
DFlash & \(\mathbf{32.09} \pm 0.15\;[3]\) & \(31.36 \pm 0.28\;[3]\) & \(76.04 \pm 0.83\;[3]\) & \(61.58 \pm 0.71\;[3]\) \\
AR \textperiodcentered{} SGLang & \(34.53 \pm 6.5\;[3]\) & \(30.17 \pm 0.18\;[3]\) & \(\mathbf{76.03} \pm 0.35\;[3]\) & \(\mathbf{60.87} \pm 1.3\;[3]\) \\
AR \textperiodcentered{} vLLM ModelRunnerV1 & \(51.89 \pm 5.3\;[3]\) & \(47.39 \pm 0.089\;[3]\) & \(85.33 \pm 1.9\;[3]\) & \(77.24 \pm 0.66\;[3]\) \\
Pilot FFN 100\% & \(35.71 \pm 6.9\;[3]\) & \(29.67 \pm 0.32\;[3]\) & \(83.3 \pm 2.8\;[3]\) & \(67.52 \pm 0.64\;[3]\) \\
Pilot FFN 75\% & \(35.47 \pm 6.7\;[3]\) & \(\mathbf{29.62} \pm 0.14\;[3]\) & \(81.65 \pm 0.98\;[3]\) & \(66.07 \pm 0.32\;[3]\) \\
Pilot FFN 50\% & \(35.61 \pm 7.2\;[3]\) & \(29.81 \pm 0.51\;[3]\) & \(82.32 \pm 1.2\;[3]\) & \(66.86 \pm 0.071\;[3]\) \\
AR \textperiodcentered{} Transformers & \(36.27 \pm 8.2\;[3]\) & \(30.05 \pm 0.43\;[3]\) & \(80.98 \pm 1.8\;[3]\) & \(66.73 \pm 0.28\;[3]\) \\
\bottomrule
\end{tabular}\par\medskip
\begin{minipage}{\linewidth}\footnotesize
Time each request from actual client submission until the first streamed output arrives. Use the first nonempty output chunk as the client observation. Exclude metadata-only events. Report the median and p95 among requests that return output, and count requests that return none.
The bold mean is the best observed value in each benchmark column. The comparison covers the recorded implementations under their stated configurations. At least two eligible method entries are required. Ties at the displayed precision share bold formatting. Bold formatting conveys an observed comparison and makes no statistical-significance claim. Brackets show the available repetition count. Response differences remain part of the interpretation. Lower values are preferred.
Brackets give eligible repetition counts. Standard deviation appears when at least two repetitions are available. Each repetition uses three prompts. The evidence records preserve observation counts and measurement scopes. Empty cells lack an eligible measured value.
\end{minipage}
\end{table*}

\begin{table*}[!htp]
\centering\small
\setlength{\tabcolsep}{3pt}
\caption{Maximum inter-chunk latency per request (ms, p95 across eligible requests).}
\label{tab:appendix-metric-pause-none}
\textit{Mathematics}\par\smallskip
\begin{tabular}{@{}>{\raggedright\arraybackslash}p{3.7cm}rrr@{}}
\toprule
Method & GSM8K & MATH-500 & AIME 2025 \\
\midrule
EAGLE-3 & \(32.34 \pm 0.32\;[3]\) & \(33.17 \pm 0.6\;[3]\) & \(34.21 \pm 0.52\;[3]\) \\
P-EAGLE & \(27.3 \pm 0.48\;[3]\) & \(27.42 \pm 0.8\;[3]\) & \(28.13 \pm 0.37\;[3]\) \\
DFlash & \(31.8 \pm 0.071\;[3]\) & \(31.26 \pm 0.16\;[3]\) & \(32.17 \pm 0.47\;[3]\) \\
AR \textperiodcentered{} SGLang & \(24.33 \pm 0.43\;[3]\) & \(30.81 \pm 12\;[3]\) & \(\mathbf{24.41} \pm 0.092\;[3]\) \\
AR \textperiodcentered{} vLLM ModelRunnerV1 & \(\mathbf{22.09} \pm 0.48\;[3]\) & \(\mathbf{21.94} \pm 0.23\;[3]\) & \(28.44 \pm 10\;[3]\) \\
Pilot FFN 100\% & \(32.59 \pm 0.11\;[3]\) & \(33.59 \pm 0.16\;[3]\) & \(38.92 \pm 0.024\;[3]\) \\
Pilot FFN 75\% & \(32.65 \pm 0.2\;[3]\) & \(33.87 \pm 0.19\;[3]\) & \(38.99 \pm 0.23\;[3]\) \\
Pilot FFN 50\% & \(33.41 \pm 0.14\;[3]\) & \(34.03 \pm 0.36\;[3]\) & \(39.5 \pm 0.098\;[3]\) \\
AR \textperiodcentered{} Transformers & \(25.21 \pm 0.18\;[3]\) & \(25.56 \pm 0.15\;[3]\) & \(26.9 \pm 0.16\;[3]\) \\
\bottomrule
\end{tabular}\par\medskip
\textit{Code}\par\smallskip
\begin{tabular}{@{}>{\raggedright\arraybackslash}p{3.7cm}rrr@{}}
\toprule
Method & HumanEval & MBPP & LiveCodeBench \\
\midrule
EAGLE-3 & \(32.99 \pm 0.94\;[3]\) & \(32.74 \pm 0.6\;[3]\) & \(32.84 \pm 0.26\;[3]\) \\
P-EAGLE & \(26.2 \pm 0.098\;[3]\) & \(26.63 \pm 0.55\;[3]\) & \(27.4 \pm 0.25\;[3]\) \\
DFlash & \(31.07 \pm 0.083\;[3]\) & \(31.02 \pm 0.066\;[3]\) & \(31.38 \pm 0.16\;[3]\) \\
AR \textperiodcentered{} SGLang & \(23.91 \pm 0.83\;[3]\) & \(23.69 \pm 0.43\;[3]\) & \(24.49 \pm 1.2\;[3]\) \\
AR \textperiodcentered{} vLLM ModelRunnerV1 & \(\mathbf{22.11} \pm 0.63\;[3]\) & \(\mathbf{21.62} \pm 0.1\;[3]\) & \(\mathbf{21.95} \pm 0.24\;[3]\) \\
Pilot FFN 100\% & \(32.24 \pm 0.2\;[3]\) & \(32.41 \pm 0.31\;[3]\) & \(34.1 \pm 0.061\;[3]\) \\
Pilot FFN 75\% & \(32.34 \pm 0.067\;[3]\) & \(32.5 \pm 0.32\;[3]\) & \(33.78 \pm 0.44\;[3]\) \\
Pilot FFN 50\% & \(32.96 \pm 0.31\;[3]\) & \(33.19 \pm 0.58\;[3]\) & \(34.5 \pm 0.16\;[3]\) \\
AR \textperiodcentered{} Transformers & \(25.22 \pm 0.19\;[3]\) & \(25.43 \pm 0.43\;[3]\) & \(25.36 \pm 0.088\;[3]\) \\
\bottomrule
\end{tabular}\par\medskip
\textit{Conversation, instruction following, and summarization}\par\smallskip
\begin{tabular}{@{}>{\raggedright\arraybackslash}p{3.7cm}rrrr@{}}
\toprule
Method & MT-Bench & Alpaca & Arena-Hard & CNN/DailyMail \\
\midrule
EAGLE-3 & \(33.33 \pm 0.41\;[3]\) & \(32.75 \pm 0.5\;[3]\) & \(33.64 \pm 0.44\;[3]\) & \(32.24 \pm 0.96\;[3]\) \\
P-EAGLE & \(27.28 \pm 0.29\;[3]\) & \(26.81 \pm 0.35\;[3]\) & \(27.19 \pm 0.33\;[3]\) & \(27.05 \pm 0.24\;[3]\) \\
DFlash & \(31.66 \pm 0.4\;[3]\) & \(31.61 \pm 0.51\;[3]\) & \(31.82 \pm 0.43\;[3]\) & \(31.68 \pm 0.43\;[3]\) \\
AR \textperiodcentered{} SGLang & \(24.1 \pm 0.22\;[3]\) & \(23.71 \pm 0.35\;[3]\) & \(23.87 \pm 0.2\;[3]\) & \(23.91 \pm 0.36\;[3]\) \\
AR \textperiodcentered{} vLLM ModelRunnerV1 & \(\mathbf{24} \pm 2\;[3]\) & \(\mathbf{21.88} \pm 0.088\;[3]\) & \(\mathbf{22.21} \pm 0.43\;[3]\) & \(\mathbf{22.62} \pm 1.1\;[3]\) \\
Pilot FFN 100\% & \(36.08 \pm 3\;[3]\) & \(32.98 \pm 0.027\;[3]\) & \(43.36 \pm 1.1\;[3]\) & \(38.85 \pm 0.2\;[3]\) \\
Pilot FFN 75\% & \(34.54 \pm 0.11\;[3]\) & \(32.85 \pm 0.14\;[3]\) & \(40.18 \pm 0.3\;[3]\) & \(36.94 \pm 0.18\;[3]\) \\
Pilot FFN 50\% & \(34.95 \pm 0.43\;[3]\) & \(34.01 \pm 0.18\;[3]\) & \(38.92 \pm 0.34\;[3]\) & \(36.38 \pm 0.24\;[3]\) \\
AR \textperiodcentered{} Transformers & \(25.7 \pm 0.31\;[3]\) & \(25.4 \pm 0.24\;[3]\) & \(25.81 \pm 0.0024\;[3]\) & \(25.67 \pm 0.044\;[3]\) \\
\bottomrule
\end{tabular}\par\medskip
\begin{minipage}{\linewidth}\footnotesize
A streaming server sends an answer in pieces called chunks. For each request with at least two nonempty chunks, find the longest gap between successive chunks. Take the p95 of those longest gaps, giving each eligible request one observation.
The bold mean is the best observed value in each benchmark column. The comparison covers the recorded implementations under their stated configurations. At least two eligible method entries are required. Ties at the displayed precision share bold formatting. Bold formatting conveys an observed comparison and makes no statistical-significance claim. Brackets show the available repetition count. Response differences remain part of the interpretation. Lower values are preferred.
Brackets give eligible repetition counts. Standard deviation appears when at least two repetitions are available. Each repetition uses three prompts. The evidence records preserve observation counts and measurement scopes. Empty cells lack an eligible measured value.
\end{minipage}
\end{table*}

\begin{table*}[!htp]
\centering\small
\setlength{\tabcolsep}{3pt}
\caption{Peak GPU memory; Sampled whole-device memory (GiB/device).}
\label{tab:appendix-metric-memory-sampled-whole-device-memory-none}
\textit{Mathematics}\par\smallskip
\begin{tabular}{@{}>{\raggedright\arraybackslash}p{3.7cm}rrr@{}}
\toprule
Method & GSM8K & MATH-500 & AIME 2025 \\
\midrule
EAGLE-3 & \(31.15 \pm 0.0011\;[3]\) & \(31.15 \pm 0.0011\;[3]\) & \(31.15 \pm 0.0011\;[3]\) \\
P-EAGLE & \(29.26 \pm 0\;[3]\) & \(29.26 \pm 0\;[3]\) & \(29.26 \pm 0\;[3]\) \\
DFlash & \(32.2 \pm 0.016\;[3]\) & \(32.2 \pm 0.016\;[3]\) & \(32.2 \pm 0.016\;[3]\) \\
AR \textperiodcentered{} SGLang & \(32.01 \pm 0\;[3]\) & \(32.01 \pm 0\;[3]\) & \(32.01 \pm 0\;[3]\) \\
AR \textperiodcentered{} vLLM ModelRunnerV1 & \(29.2 \pm 0\;[3]\) & \(29.2 \pm 0\;[3]\) & \(29.2 \pm 0\;[3]\) \\
Pilot FFN 100\% & \(18.39 \pm 0.015\;[3]\) & \(18.39 \pm 0.015\;[3]\) & \(18.39 \pm 0.015\;[3]\) \\
AR \textperiodcentered{} Transformers & \(\mathbf{16.57} \pm 0\;[2]\) & \(\mathbf{16.57} \pm 0\;[2]\) & \(\mathbf{16.57} \pm 0\;[2]\) \\
\bottomrule
\end{tabular}\par\medskip
\textit{Code}\par\smallskip
\begin{tabular}{@{}>{\raggedright\arraybackslash}p{3.7cm}rrr@{}}
\toprule
Method & HumanEval & MBPP & LiveCodeBench \\
\midrule
EAGLE-3 & \(31.15 \pm 0.0011\;[3]\) & \(31.15 \pm 0.0011\;[3]\) & \(31.15 \pm 0.0011\;[3]\) \\
P-EAGLE & \(29.26 \pm 0\;[3]\) & \(29.26 \pm 0\;[3]\) & \(29.26 \pm 0\;[3]\) \\
DFlash & \(32.2 \pm 0.016\;[3]\) & \(32.2 \pm 0.016\;[3]\) & \(32.2 \pm 0.016\;[3]\) \\
AR \textperiodcentered{} SGLang & \(32.01 \pm 0\;[3]\) & \(32.01 \pm 0\;[3]\) & \(32.01 \pm 0\;[3]\) \\
AR \textperiodcentered{} vLLM ModelRunnerV1 & \(29.2 \pm 0\;[3]\) & \(29.2 \pm 0\;[3]\) & \(29.2 \pm 0\;[3]\) \\
Pilot FFN 100\% & \(18.39 \pm 0.015\;[3]\) & \(18.39 \pm 0.015\;[3]\) & \(18.39 \pm 0\;[3]\) \\
AR \textperiodcentered{} Transformers & \(\mathbf{16.57} \pm 0\;[3]\) & \(\mathbf{16.57} \pm 0\;[3]\) & \(\mathbf{16.57} \pm 0\;[3]\) \\
\bottomrule
\end{tabular}\par\medskip
\textit{Conversation, instruction following, and summarization}\par\smallskip
\begin{tabular}{@{}>{\raggedright\arraybackslash}p{3.7cm}rrrr@{}}
\toprule
Method & MT-Bench & Alpaca & Arena-Hard & CNN/DailyMail \\
\midrule
EAGLE-3 & \(31.13 \pm 0.047\;[3]\) & \(31.13 \pm 0.047\;[3]\) & \(31.15 \pm 0\;[3]\) & \(31.15 \pm 0\;[3]\) \\
P-EAGLE & \(29.17 \pm 0.15\;[3]\) & \(29.17 \pm 0.15\;[3]\) & \(29.26 \pm 0.0011\;[3]\) & \(29.26 \pm 0.0011\;[3]\) \\
DFlash & \(32.13 \pm 0.11\;[3]\) & \(32.13 \pm 0.11\;[3]\) & \(32.2 \pm 0.016\;[3]\) & \(32.2 \pm 0.016\;[3]\) \\
AR \textperiodcentered{} SGLang & \(31.97 \pm 0.07\;[3]\) & \(31.97 \pm 0.07\;[3]\) & \(32.01 \pm 0\;[3]\) & \(32.01 \pm 0\;[3]\) \\
AR \textperiodcentered{} vLLM ModelRunnerV1 & \(29.16 \pm 0.07\;[3]\) & \(29.16 \pm 0.07\;[3]\) & \(29.2 \pm 0\;[3]\) & \(29.2 \pm 0\;[3]\) \\
Pilot FFN 100\% & \(18.13 \pm 0.45\;[3]\) & \(18.13 \pm 0.45\;[3]\) & \(18.39 \pm 0.015\;[3]\) & \(18.39 \pm 0.015\;[3]\) \\
AR \textperiodcentered{} Transformers & \(\mathbf{16.57} \pm 0\;[2]\) & \(\mathbf{16.57} \pm 0\;[2]\) & \(\mathbf{16.57} \pm 0\;[2]\) & \(\mathbf{16.57} \pm 0\;[2]\) \\
\bottomrule
\end{tabular}\par\medskip
\begin{minipage}{\linewidth}\footnotesize
Record the highest observed memory use on each GPU during the stated interval. Keep separate readings for memory used on the whole device, memory occupied by live tensors and memory reserved by the framework. A gibibyte (GiB) equals 1,073,741,824 bytes.
The bold mean is the best observed value in each benchmark column. The comparison covers the recorded implementations under their stated configurations. At least two eligible method entries are required. Ties at the displayed precision share bold formatting. Bold formatting conveys an observed comparison and makes no statistical-significance claim. Brackets show the available repetition count. Response differences remain part of the interpretation. Lower values are preferred.
Brackets give eligible repetition counts. Standard deviation appears when at least two repetitions are available. Each repetition uses three prompts. The evidence records preserve observation counts and measurement scopes. Empty cells lack an eligible measured value.
\end{minipage}
\end{table*}

\begin{table*}[!htp]
\centering\small
\setlength{\tabcolsep}{3pt}
\caption{Peak GPU memory; Framework-allocated memory (GiB/device).}
\label{tab:appendix-metric-memory-framework-allocated-memory-none}
\textit{Mathematics}\par\smallskip
\begin{tabular}{@{}>{\raggedright\arraybackslash}p{3.7cm}rrr@{}}
\toprule
Method & GSM8K & MATH-500 & AIME 2025 \\
\midrule
EAGLE-3 & \(30.38 \pm 0\;[3]\) & \(30.39 \pm 0\;[3]\) & \(30.46 \pm 0\;[3]\) \\
DFlash &  &  &  \\
\bottomrule
\end{tabular}\par\medskip
\textit{Code}\par\smallskip
\begin{tabular}{@{}>{\raggedright\arraybackslash}p{3.7cm}rrr@{}}
\toprule
Method & HumanEval & MBPP & LiveCodeBench \\
\midrule
EAGLE-3 & \(30.39 \pm 0\;[3]\) & \(30.38 \pm 0\;[3]\) & \(30.42 \pm 0\;[3]\) \\
DFlash &  &  &  \\
\bottomrule
\end{tabular}\par\medskip
\textit{Conversation, instruction following, and summarization}\par\smallskip
\begin{tabular}{@{}>{\raggedright\arraybackslash}p{3.7cm}rrrr@{}}
\toprule
Method & MT-Bench & Alpaca & Arena-Hard & CNN/DailyMail \\
\midrule
EAGLE-3 & \(\mathbf{30.38} \pm 0\;[3]\) & \(\mathbf{30.38} \pm 7\times 10^{-5}\;[3]\) & \(\mathbf{30.51} \pm 7.5\times 10^{-5}\;[3]\) & \(\mathbf{30.47} \pm 0\;[3]\) \\
DFlash & \(31.29\;[1]\) & \(31.28\;[1]\) & \(31.42\;[1]\) & \(31.41\;[1]\) \\
\bottomrule
\end{tabular}\par\medskip
\begin{minipage}{\linewidth}\footnotesize
Record the highest observed memory use on each GPU during the stated interval. Keep separate readings for memory used on the whole device, memory occupied by live tensors and memory reserved by the framework. A gibibyte (GiB) equals 1,073,741,824 bytes.
The bold mean is the best observed value in each benchmark column. The comparison covers the recorded implementations under their stated configurations. At least two eligible method entries are required. Ties at the displayed precision share bold formatting. Bold formatting conveys an observed comparison and makes no statistical-significance claim. Brackets show the available repetition count. Response differences remain part of the interpretation. Lower values are preferred.
Brackets give eligible repetition counts. Standard deviation appears when at least two repetitions are available. Each repetition uses three prompts. The evidence records preserve observation counts and measurement scopes. Empty cells lack an eligible measured value.
\end{minipage}
\end{table*}

\begin{table*}[!htp]
\centering\small
\setlength{\tabcolsep}{3pt}
\caption{Peak GPU memory; Framework-reserved memory (GiB/device).}
\label{tab:appendix-metric-memory-framework-reserved-memory-none}
\textit{Mathematics}\par\smallskip
\begin{tabular}{@{}>{\raggedright\arraybackslash}p{3.7cm}rrr@{}}
\toprule
Method & GSM8K & MATH-500 & AIME 2025 \\
\midrule
EAGLE-3 & \(30.61 \pm 0.0011\;[3]\) & \(30.61 \pm 0.0011\;[3]\) & \(30.61 \pm 0.0011\;[3]\) \\
DFlash &  &  &  \\
\bottomrule
\end{tabular}\par\medskip
\textit{Code}\par\smallskip
\begin{tabular}{@{}>{\raggedright\arraybackslash}p{3.7cm}rrr@{}}
\toprule
Method & HumanEval & MBPP & LiveCodeBench \\
\midrule
EAGLE-3 & \(30.61 \pm 0.0011\;[3]\) & \(30.61 \pm 0.0011\;[3]\) & \(30.61 \pm 0.0011\;[3]\) \\
DFlash &  &  &  \\
\bottomrule
\end{tabular}\par\medskip
\textit{Conversation, instruction following, and summarization}\par\smallskip
\begin{tabular}{@{}>{\raggedright\arraybackslash}p{3.7cm}rrrr@{}}
\toprule
Method & MT-Bench & Alpaca & Arena-Hard & CNN/DailyMail \\
\midrule
EAGLE-3 & \(\mathbf{30.56} \pm 0.081\;[3]\) & \(\mathbf{30.56} \pm 0.081\;[3]\) & \(\mathbf{30.61} \pm 0\;[3]\) & \(\mathbf{30.61} \pm 0\;[3]\) \\
DFlash & \(31.47\;[1]\) & \(31.47\;[1]\) & \(31.65\;[1]\) & \(31.65\;[1]\) \\
\bottomrule
\end{tabular}\par\medskip
\begin{minipage}{\linewidth}\footnotesize
Record the highest observed memory use on each GPU during the stated interval. Keep separate readings for memory used on the whole device, memory occupied by live tensors and memory reserved by the framework. A gibibyte (GiB) equals 1,073,741,824 bytes.
The bold mean is the best observed value in each benchmark column. The comparison covers the recorded implementations under their stated configurations. At least two eligible method entries are required. Ties at the displayed precision share bold formatting. Bold formatting conveys an observed comparison and makes no statistical-significance claim. Brackets show the available repetition count. Response differences remain part of the interpretation. Lower values are preferred.
Brackets give eligible repetition counts. Standard deviation appears when at least two repetitions are available. Each repetition uses three prompts. The evidence records preserve observation counts and measurement scopes. Empty cells lack an eligible measured value.
\end{minipage}
\end{table*}

\begin{table*}[!htp]
\centering\small
\setlength{\tabcolsep}{3pt}
\caption{Output throughput per GPU (output tokens/s/GPU).}
\label{tab:appendix-metric-per-gpu-throughput-none}
\textit{Mathematics}\par\smallskip
\begin{tabular}{@{}>{\raggedright\arraybackslash}p{3.7cm}rrr@{}}
\toprule
Method & GSM8K & MATH-500 & AIME 2025 \\
\midrule
EAGLE-3 & \(106.5 \pm 1.1\;[3]\) & \(103.2 \pm 0.086\;[3]\) & \(104.2 \pm 0.26\;[3]\) \\
P-EAGLE & \(117.5 \pm 0.82\;[3]\) & \(\mathbf{153.5} \pm 0.012\;[3]\) & \(\mathbf{157.9} \pm 0.0071\;[3]\) \\
DFlash & \(\mathbf{213.4} \pm 0.29\;[3]\) & \(138.8 \pm 0.31\;[3]\) & \(154.9 \pm 0.062\;[3]\) \\
AR \textperiodcentered{} SGLang & \(43.27 \pm 0.046\;[3]\) & \(43.44 \pm 0.0078\;[3]\) & \(43.34 \pm 0.058\;[3]\) \\
AR \textperiodcentered{} vLLM ModelRunnerV1 & \(46.28 \pm 0.0013\;[3]\) & \(46.47 \pm 0.002\;[3]\) & \(46.3 \pm 0.00078\;[3]\) \\
Pilot FFN 100\% & \(39.02 \pm 0.036\;[3]\) & \(33.96 \pm 0.022\;[3]\) & \(32.42 \pm 0.026\;[3]\) \\
Pilot FFN 75\% & \(39.02 \pm 0.045\;[3]\) & \(31.92 \pm 2.8\;[3]\) & \(32.74 \pm 0.054\;[3]\) \\
Pilot FFN 50\% & \(37.59 \pm 0.082\;[3]\) & \(32.85 \pm 0.22\;[3]\) & \(31.56 \pm 0.035\;[3]\) \\
AR \textperiodcentered{} Transformers & \(40.17 \pm 0.15\;[3]\) & \(40.01 \pm 0.027\;[3]\) & \(39.12 \pm 0.05\;[3]\) \\
\bottomrule
\end{tabular}\par\medskip
\textit{Code}\par\smallskip
\begin{tabular}{@{}>{\raggedright\arraybackslash}p{3.7cm}rrr@{}}
\toprule
Method & HumanEval & MBPP & LiveCodeBench \\
\midrule
EAGLE-3 & \(100.3 \pm 0.28\;[3]\) & \(110.7 \pm 1.1\;[3]\) & \(99.48 \pm 0.39\;[3]\) \\
P-EAGLE & \(114.1 \pm 0.011\;[3]\) & \(124.9 \pm 0.0019\;[3]\) & \(107.9 \pm 0.01\;[3]\) \\
DFlash & \(\mathbf{137.2} \pm 1.4\;[3]\) & \(\mathbf{155.1} \pm 0.35\;[3]\) & \(\mathbf{117.2} \pm 0.21\;[3]\) \\
AR \textperiodcentered{} SGLang & \(42.88 \pm 0.081\;[3]\) & \(43.27 \pm 0.056\;[3]\) & \(43.27 \pm 0.043\;[3]\) \\
AR \textperiodcentered{} vLLM ModelRunnerV1 & \(45.72 \pm 0.0035\;[3]\) & \(46.17 \pm 0.029\;[3]\) & \(46.21 \pm 0.021\;[3]\) \\
Pilot FFN 100\% & \(36.12 \pm 0.052\;[3]\) & \(37.1 \pm 0.057\;[3]\) & \(35.49 \pm 0.18\;[3]\) \\
Pilot FFN 75\% & \(35.42 \pm 0.064\;[3]\) & \(35.39 \pm 0.044\;[3]\) & \(35.1 \pm 0.082\;[3]\) \\
Pilot FFN 50\% & \(35.66 \pm 0.13\;[3]\) & \(36.68 \pm 0.06\;[3]\) & \(34.45 \pm 0.03\;[3]\) \\
AR \textperiodcentered{} Transformers & \(39.73 \pm 0.077\;[3]\) & \(40.03 \pm 0.042\;[3]\) & \(39.87 \pm 0.13\;[3]\) \\
\bottomrule
\end{tabular}\par\medskip
\textit{Conversation, instruction following, and summarization}\par\smallskip
\begin{tabular}{@{}>{\raggedright\arraybackslash}p{3.7cm}rrrr@{}}
\toprule
Method & MT-Bench & Alpaca & Arena-Hard & CNN/DailyMail \\
\midrule
EAGLE-3 & \(\mathbf{110.3} \pm 0.77\;[3]\) & \(116.7 \pm 0.37\;[3]\) & \(94.66 \pm 0.49\;[3]\) & \(\mathbf{97.98} \pm 1.2\;[3]\) \\
P-EAGLE & \(82.75 \pm 0.37\;[3]\) & \(105.3 \pm 0.0069\;[3]\) & \(93.02 \pm 0.12\;[3]\) & \(83.56 \pm 0.012\;[3]\) \\
DFlash & \(109.7 \pm 0.13\;[3]\) & \(\mathbf{117.6} \pm 0.39\;[3]\) & \(\mathbf{112.3} \pm 0.14\;[3]\) & \(97.67 \pm 0.067\;[3]\) \\
AR \textperiodcentered{} SGLang & \(43.37 \pm 0.067\;[3]\) & \(43.38 \pm 0.011\;[3]\) & \(43.25 \pm 0.018\;[3]\) & \(43 \pm 0.044\;[3]\) \\
AR \textperiodcentered{} vLLM ModelRunnerV1 & \(46.43 \pm 0.0034\;[3]\) & \(46.37 \pm 0.0016\;[3]\) & \(46.19 \pm 0.0067\;[3]\) & \(45.8 \pm 0.0037\;[3]\) \\
Pilot FFN 100\% & \(36.3 \pm 0.038\;[3]\) & \(36.34 \pm 0.055\;[3]\) & \(34.89 \pm 0.1\;[3]\) & \(34.28 \pm 0.019\;[3]\) \\
Pilot FFN 75\% & \(36.33 \pm 0.0059\;[3]\) & \(36.42 \pm 0.045\;[3]\) & \(34.36 \pm 0.012\;[3]\) & \(34.48 \pm 0.044\;[3]\) \\
Pilot FFN 50\% & \(35.71 \pm 0.11\;[3]\) & \(35.54 \pm 0.13\;[3]\) & \(33.96 \pm 0.095\;[3]\) & \(33.25 \pm 0.078\;[3]\) \\
AR \textperiodcentered{} Transformers & \(39.97 \pm 0.093\;[3]\) & \(40.02 \pm 0.029\;[3]\) & \(39.76 \pm 0.039\;[3]\) & \(39.39 \pm 0.013\;[3]\) \\
\bottomrule
\end{tabular}\par\medskip
\begin{minipage}{\linewidth}\footnotesize
Divide a benchmark's delivered output tokens by its elapsed seconds multiplied by the allocated GPU count. The comparison tables use timed benchmark blocks. Complete-job resource totals retain their own allocation interval. Compute cost uses each timed benchmark interval and the GPU count from its final Slurm allocation. The interval includes client gaps between requests. Startup, warmups, and time outside the benchmark have separate whole-job accounting. Per-prompt GPU-seconds multiply the allocated GPU count by that prompt's completion time in seconds.
The bold mean is the best observed value in each benchmark column. The comparison covers the recorded implementations under their stated configurations. At least two eligible method entries are required. Ties at the displayed precision share bold formatting. Bold formatting conveys an observed comparison and makes no statistical-significance claim. Brackets show the available repetition count. Response differences remain part of the interpretation. Higher values are preferred.
Brackets give eligible repetition counts. Standard deviation appears when at least two repetitions are available. Each repetition uses three prompts. The evidence records preserve observation counts and measurement scopes. Empty cells lack an eligible measured value.
\end{minipage}
\end{table*}

\begin{table*}[!htp]
\centering\small
\setlength{\tabcolsep}{3pt}
\caption{End-to-end output rate per request (output tokens/s/request).}
\label{tab:appendix-metric-request-speed-none}
\textit{Mathematics}\par\smallskip
\begin{tabular}{@{}>{\raggedright\arraybackslash}p{3.7cm}rrr@{}}
\toprule
Method & GSM8K & MATH-500 & AIME 2025 \\
\midrule
EAGLE-3 & \(106.6 \pm 1.7\;[3]\) & \(103.1 \pm 0.085\;[3]\) & \(100.6 \pm 0.15\;[3]\) \\
P-EAGLE & \(115.8 \pm 0.56\;[3]\) & \(\mathbf{152.5} \pm 0.0097\;[3]\) & \(\mathbf{151} \pm 0.017\;[3]\) \\
DFlash & \(\mathbf{163.6} \pm 0.63\;[3]\) & \(137.3 \pm 0.27\;[3]\) & \(\mathbf{151} \pm 0.11\;[3]\) \\
AR \textperiodcentered{} SGLang & \(43.3 \pm 0.28\;[3]\) & \(43.44 \pm 0.0061\;[3]\) & \(43.33 \pm 0.078\;[3]\) \\
AR \textperiodcentered{} vLLM ModelRunnerV1 & \(46.23 \pm 0.0015\;[3]\) & \(46.47 \pm 0.0019\;[3]\) & \(46.28 \pm 0.0024\;[3]\) \\
Pilot FFN 100\% & \(38.54 \pm 0.044\;[3]\) & \(34.07 \pm 0.031\;[3]\) & \(32.6 \pm 0.025\;[3]\) \\
Pilot FFN 75\% & \(38.48 \pm 0.061\;[3]\) & \(32.75 \pm 1.7\;[3]\) & \(32.5 \pm 0.043\;[3]\) \\
Pilot FFN 50\% & \(37.15 \pm 0.066\;[3]\) & \(33.01 \pm 0.21\;[3]\) & \(31.53 \pm 0.017\;[3]\) \\
AR \textperiodcentered{} Transformers & \(40.15 \pm 0.12\;[3]\) & \(40.03 \pm 0.038\;[3]\) & \(39.54 \pm 0.032\;[3]\) \\
\bottomrule
\end{tabular}\par\medskip
\textit{Code}\par\smallskip
\begin{tabular}{@{}>{\raggedright\arraybackslash}p{3.7cm}rrr@{}}
\toprule
Method & HumanEval & MBPP & LiveCodeBench \\
\midrule
EAGLE-3 & \(99.18 \pm 0.24\;[3]\) & \(108.5 \pm 1\;[3]\) & \(105.4 \pm 0.86\;[3]\) \\
P-EAGLE & \(114.9 \pm 0.004\;[3]\) & \(125.7 \pm 0.0013\;[3]\) & \(111.2 \pm 0.015\;[3]\) \\
DFlash & \(\mathbf{143.7} \pm 1.4\;[3]\) & \(\mathbf{153.1} \pm 0.41\;[3]\) & \(\mathbf{131.9} \pm 0.56\;[3]\) \\
AR \textperiodcentered{} SGLang & \(42.86 \pm 0.073\;[3]\) & \(43.18 \pm 0.11\;[3]\) & \(43.2 \pm 0.032\;[3]\) \\
AR \textperiodcentered{} vLLM ModelRunnerV1 & \(45.66 \pm 0.0037\;[3]\) & \(46.07 \pm 0.042\;[3]\) & \(46.09 \pm 0.019\;[3]\) \\
Pilot FFN 100\% & \(35.99 \pm 0.057\;[3]\) & \(36.98 \pm 0.058\;[3]\) & \(35.62 \pm 0.16\;[3]\) \\
Pilot FFN 75\% & \(35.32 \pm 0.075\;[3]\) & \(35.62 \pm 0.046\;[3]\) & \(35.32 \pm 0.088\;[3]\) \\
Pilot FFN 50\% & \(35.58 \pm 0.14\;[3]\) & \(36.57 \pm 0.041\;[3]\) & \(34.62 \pm 0.036\;[3]\) \\
AR \textperiodcentered{} Transformers & \(39.71 \pm 0.067\;[3]\) & \(40.01 \pm 0.055\;[3]\) & \(39.89 \pm 0.17\;[3]\) \\
\bottomrule
\end{tabular}\par\medskip
\textit{Conversation, instruction following, and summarization}\par\smallskip
\begin{tabular}{@{}>{\raggedright\arraybackslash}p{3.7cm}rrrr@{}}
\toprule
Method & MT-Bench & Alpaca & Arena-Hard & CNN/DailyMail \\
\midrule
EAGLE-3 & \(111.6 \pm 1\;[3]\) & \(113.9 \pm 0.53\;[3]\) & \(98.26 \pm 0.66\;[3]\) & \(\mathbf{98.35} \pm 1.2\;[3]\) \\
P-EAGLE & \(88.29 \pm 1.5\;[3]\) & \(103.4 \pm 0.02\;[3]\) & \(98.25 \pm 0.074\;[3]\) & \(83.52 \pm 0.012\;[3]\) \\
DFlash & \(\mathbf{117.1} \pm 0.17\;[3]\) & \(\mathbf{121.2} \pm 0.31\;[3]\) & \(\mathbf{122.2} \pm 0.37\;[3]\) & \(97.6 \pm 0.067\;[3]\) \\
AR \textperiodcentered{} SGLang & \(43.3 \pm 0.054\;[3]\) & \(43.21 \pm 0.048\;[3]\) & \(43.26 \pm 0.035\;[3]\) & \(43 \pm 0.041\;[3]\) \\
AR \textperiodcentered{} vLLM ModelRunnerV1 & \(46.27 \pm 0.026\;[3]\) & \(45.94 \pm 0.0027\;[3]\) & \(46.14 \pm 0.0062\;[3]\) & \(45.8 \pm 0.0036\;[3]\) \\
Pilot FFN 100\% & \(36.21 \pm 0.049\;[3]\) & \(36.46 \pm 0.028\;[3]\) & \(35.38 \pm 0.094\;[3]\) & \(34.3 \pm 0.019\;[3]\) \\
Pilot FFN 75\% & \(36.22 \pm 0.039\;[3]\) & \(36.14 \pm 0.089\;[3]\) & \(34.99 \pm 0.015\;[3]\) & \(34.5 \pm 0.049\;[3]\) \\
Pilot FFN 50\% & \(35.35 \pm 0.12\;[3]\) & \(35.16 \pm 0.17\;[3]\) & \(34.14 \pm 0.097\;[3]\) & \(33.29 \pm 0.086\;[3]\) \\
AR \textperiodcentered{} Transformers & \(40.01 \pm 0.11\;[3]\) & \(39.83 \pm 0.023\;[3]\) & \(39.76 \pm 0.08\;[3]\) & \(39.4 \pm 0.012\;[3]\) \\
\bottomrule
\end{tabular}\par\medskip
\begin{minipage}{\linewidth}\footnotesize
For each completed request, divide its returned output tokens by the seconds from actual submission through completion. Average the request rates equally. The clock includes waiting, prompt processing and generation.
The bold mean is the best observed value in each benchmark column. The comparison covers the recorded implementations under their stated configurations. At least two eligible method entries are required. Ties at the displayed precision share bold formatting. Bold formatting conveys an observed comparison and makes no statistical-significance claim. Brackets show the available repetition count. Response differences remain part of the interpretation. Higher values are preferred.
Brackets give eligible repetition counts. Standard deviation appears when at least two repetitions are available. Each repetition uses three prompts. The evidence records preserve observation counts and measurement scopes. Empty cells lack an eligible measured value.
\end{minipage}
\end{table*}

\begin{table*}[!htp]
\centering\small
\setlength{\tabcolsep}{3pt}
\caption{Time per output token (TPOT) (ms/output token).}
\label{tab:appendix-metric-tpot-none}
\textit{Mathematics}\par\smallskip
\begin{tabular}{@{}>{\raggedright\arraybackslash}p{3.7cm}rrr@{}}
\toprule
Method & GSM8K & MATH-500 & AIME 2025 \\
\midrule
EAGLE-3 & \(9.238 \pm 0.15\;[3]\) & \(9.653 \pm 0.0081\;[3]\) & \(9.967 \pm 0.015\;[3]\) \\
P-EAGLE & \(8.386 \pm 0.0086\;[3]\) & \(\mathbf{6.474} \pm 0.00062\;[3]\) & \(\mathbf{6.646} \pm 2.5\times 10^{-5}\;[3]\) \\
DFlash & \(\mathbf{6.485} \pm 0.034\;[3]\) & \(7.396 \pm 0.016\;[3]\) & \(6.698 \pm 0.0046\;[3]\) \\
AR \textperiodcentered{} SGLang & \(23.06 \pm 0.15\;[3]\) & \(23 \pm 0.0031\;[3]\) & \(23.05 \pm 0.042\;[3]\) \\
AR \textperiodcentered{} vLLM ModelRunnerV1 & \(21.5 \pm 0.00063\;[3]\) & \(21.47 \pm 0.0008\;[3]\) & \(21.57 \pm 0.00076\;[3]\) \\
Pilot FFN 100\% & \(25.98 \pm 0.028\;[3]\) & \(29.34 \pm 0.03\;[3]\) & \(30.65 \pm 0.023\;[3]\) \\
Pilot FFN 75\% & \(26.07 \pm 0.045\;[3]\) & \(29.61 \pm 0.039\;[3]\) & \(30.75 \pm 0.04\;[3]\) \\
Pilot FFN 50\% & \(26.95 \pm 0.051\;[3]\) & \(30.3 \pm 0.19\;[3]\) & \(31.69 \pm 0.017\;[3]\) \\
AR \textperiodcentered{} Transformers & \(24.86 \pm 0.072\;[3]\) & \(24.97 \pm 0.024\;[3]\) & \(25.28 \pm 0.021\;[3]\) \\
\bottomrule
\end{tabular}\par\medskip
\textit{Code}\par\smallskip
\begin{tabular}{@{}>{\raggedright\arraybackslash}p{3.7cm}rrr@{}}
\toprule
Method & HumanEval & MBPP & LiveCodeBench \\
\midrule
EAGLE-3 & \(9.864 \pm 0.014\;[3]\) & \(9.042 \pm 0.082\;[3]\) & \(9.52 \pm 0.067\;[3]\) \\
P-EAGLE & \(8.452 \pm 0.00073\;[3]\) & \(7.691 \pm 0.00077\;[3]\) & \(9 \pm 0.0014\;[3]\) \\
DFlash & \(\mathbf{6.985} \pm 0.073\;[3]\) & \(\mathbf{6.447} \pm 0.017\;[3]\) & \(\mathbf{7.956} \pm 0.023\;[3]\) \\
AR \textperiodcentered{} SGLang & \(23.23 \pm 0.04\;[3]\) & \(23.1 \pm 0.062\;[3]\) & \(23.08 \pm 0.018\;[3]\) \\
AR \textperiodcentered{} vLLM ModelRunnerV1 & \(21.65 \pm 0.00083\;[3]\) & \(21.53 \pm 0.00065\;[3]\) & \(21.55 \pm 0.0017\;[3]\) \\
Pilot FFN 100\% & \(27.73 \pm 0.044\;[3]\) & \(27.03 \pm 0.041\;[3]\) & \(28.04 \pm 0.12\;[3]\) \\
Pilot FFN 75\% & \(28.28 \pm 0.058\;[3]\) & \(28.09 \pm 0.037\;[3]\) & \(28.28 \pm 0.071\;[3]\) \\
Pilot FFN 50\% & \(28.05 \pm 0.11\;[3]\) & \(27.34 \pm 0.03\;[3]\) & \(28.88 \pm 0.028\;[3]\) \\
AR \textperiodcentered{} Transformers & \(25.08 \pm 0.042\;[3]\) & \(24.94 \pm 0.033\;[3]\) & \(24.99 \pm 0.11\;[3]\) \\
\bottomrule
\end{tabular}\par\medskip
\textit{Conversation, instruction following, and summarization}\par\smallskip
\begin{tabular}{@{}>{\raggedright\arraybackslash}p{3.7cm}rrrr@{}}
\toprule
Method & MT-Bench & Alpaca & Arena-Hard & CNN/DailyMail \\
\midrule
EAGLE-3 & \(8.857 \pm 0.067\;[3]\) & \(8.61 \pm 0.032\;[3]\) & \(10.41 \pm 0.054\;[3]\) & \(\mathbf{9.987} \pm 0.13\;[3]\) \\
P-EAGLE & \(11.12 \pm 0.018\;[3]\) & \(9.429 \pm 0.0021\;[3]\) & \(10.61 \pm 0.0046\;[3]\) & \(11.7 \pm 0.0013\;[3]\) \\
DFlash & \(\mathbf{8.524} \pm 0.011\;[3]\) & \(\mathbf{8.129} \pm 0.021\;[3]\) & \(\mathbf{8.522} \pm 0.017\;[3]\) & \(10.06 \pm 0.0078\;[3]\) \\
AR \textperiodcentered{} SGLang & \(23.05 \pm 0.02\;[3]\) & \(23.08 \pm 0.027\;[3]\) & \(23.05 \pm 0.018\;[3]\) & \(23.12 \pm 0.023\;[3]\) \\
AR \textperiodcentered{} vLLM ModelRunnerV1 & \(21.49 \pm 0.00035\;[3]\) & \(21.55 \pm 0.0027\;[3]\) & \(21.55 \pm 0.0015\;[3]\) & \(21.62 \pm 0.00099\;[3]\) \\
Pilot FFN 100\% & \(27.6 \pm 0.022\;[3]\) & \(27.42 \pm 0.019\;[3]\) & \(28.3 \pm 0.071\;[3]\) & \(29.05 \pm 0.02\;[3]\) \\
Pilot FFN 75\% & \(27.59 \pm 0.034\;[3]\) & \(27.69 \pm 0.069\;[3]\) & \(28.66 \pm 0.0098\;[3]\) & \(28.87 \pm 0.037\;[3]\) \\
Pilot FFN 50\% & \(28.29 \pm 0.082\;[3]\) & \(28.47 \pm 0.13\;[3]\) & \(29.29 \pm 0.081\;[3]\) & \(29.93 \pm 0.072\;[3]\) \\
AR \textperiodcentered{} Transformers & \(24.95 \pm 0.047\;[3]\) & \(25.07 \pm 0.014\;[3]\) & \(25.08 \pm 0.051\;[3]\) & \(25.22 \pm 0.008\;[3]\) \\
\bottomrule
\end{tabular}\par\medskip
\begin{minipage}{\linewidth}\footnotesize
For each completed response with N output tokens, require N greater than one. Subtract TTFT from end-to-end request latency and divide by N minus one. Average the eligible request values equally. Count excluded responses separately.
The bold mean is the best observed value in each benchmark column. The comparison covers the recorded implementations under their stated configurations. At least two eligible method entries are required. Ties at the displayed precision share bold formatting. Bold formatting conveys an observed comparison and makes no statistical-significance claim. Brackets show the available repetition count. Response differences remain part of the interpretation. Lower values are preferred.
Brackets give eligible repetition counts. Standard deviation appears when at least two repetitions are available. Each repetition uses three prompts. The evidence records preserve observation counts and measurement scopes. Empty cells lack an eligible measured value.
\end{minipage}
\end{table*}

\begin{table*}[!htp]
\centering\small
\setlength{\tabcolsep}{3pt}
\caption{Inter-chunk latency (ms/gap, p95).}
\label{tab:appendix-metric-pooled-gap-none}
\textit{Mathematics}\par\smallskip
\begin{tabular}{@{}>{\raggedright\arraybackslash}p{3.7cm}rrr@{}}
\toprule
Method & GSM8K & MATH-500 & AIME 2025 \\
\midrule
EAGLE-3 & \(32.21 \pm 0.29\;[3]\) & \(32.16 \pm 0.041\;[3]\) & \(32.52 \pm 0.13\;[3]\) \\
P-EAGLE & \(26.27 \pm 0.26\;[3]\) & \(26.38 \pm 0.15\;[3]\) & \(27.07 \pm 0.043\;[3]\) \\
DFlash & \(31.41 \pm 0.13\;[3]\) & \(30.99 \pm 0.039\;[3]\) & \(31.27 \pm 0.074\;[3]\) \\
AR \textperiodcentered{} SGLang & \(23.39 \pm 0.013\;[3]\) & \(23.18 \pm 0.018\;[3]\) & \(23.37 \pm 0.027\;[3]\) \\
AR \textperiodcentered{} vLLM ModelRunnerV1 & \(\mathbf{21.54} \pm 0.047\;[3]\) & \(\mathbf{21.54} \pm 0.0061\;[3]\) & \(\mathbf{21.81} \pm 0.01\;[3]\) \\
Pilot FFN 100\% & \(32.22 \pm 0.024\;[3]\) & \(33.2 \pm 0.048\;[3]\) & \(37.5 \pm 0.031\;[3]\) \\
Pilot FFN 75\% & \(32.37 \pm 0.089\;[3]\) & \(33.31 \pm 0.092\;[3]\) & \(37.51 \pm 0.082\;[3]\) \\
Pilot FFN 50\% & \(32.81 \pm 0.084\;[3]\) & \(33.48 \pm 0.3\;[3]\) & \(37.84 \pm 0.027\;[3]\) \\
AR \textperiodcentered{} Transformers & \(24.92 \pm 0.057\;[3]\) & \(25.26 \pm 0.016\;[3]\) & \(26.52 \pm 0.13\;[3]\) \\
\bottomrule
\end{tabular}\par\medskip
\textit{Code}\par\smallskip
\begin{tabular}{@{}>{\raggedright\arraybackslash}p{3.7cm}rrr@{}}
\toprule
Method & HumanEval & MBPP & LiveCodeBench \\
\midrule
EAGLE-3 & \(32.62 \pm 0.76\;[3]\) & \(32.31 \pm 0.38\;[3]\) & \(32.3 \pm 0.23\;[3]\) \\
P-EAGLE & \(26.01 \pm 0.043\;[3]\) & \(26.15 \pm 0.17\;[3]\) & \(26.32 \pm 0.034\;[3]\) \\
DFlash & \(30.95 \pm 0.041\;[3]\) & \(30.88 \pm 0.064\;[3]\) & \(31.06 \pm 0.15\;[3]\) \\
AR \textperiodcentered{} SGLang & \(23.2 \pm 0.014\;[3]\) & \(23.1 \pm 0.052\;[3]\) & \(23.14 \pm 0.018\;[3]\) \\
AR \textperiodcentered{} vLLM ModelRunnerV1 & \(\mathbf{21.51} \pm 0.047\;[3]\) & \(\mathbf{21.49} \pm 0.051\;[3]\) & \(\mathbf{21.53} \pm 0.0063\;[3]\) \\
Pilot FFN 100\% & \(32.01 \pm 0.051\;[3]\) & \(32.11 \pm 0.099\;[3]\) & \(33.18 \pm 0.025\;[3]\) \\
Pilot FFN 75\% & \(32.07 \pm 0.06\;[3]\) & \(32.23 \pm 0.036\;[3]\) & \(33.18 \pm 0.068\;[3]\) \\
Pilot FFN 50\% & \(32.65 \pm 0.16\;[3]\) & \(32.67 \pm 0.11\;[3]\) & \(33.88 \pm 0.065\;[3]\) \\
AR \textperiodcentered{} Transformers & \(24.96 \pm 0.071\;[3]\) & \(24.98 \pm 0.032\;[3]\) & \(25.24 \pm 0.042\;[3]\) \\
\bottomrule
\end{tabular}\par\medskip
\textit{Conversation, instruction following, and summarization}\par\smallskip
\begin{tabular}{@{}>{\raggedright\arraybackslash}p{3.7cm}rrrr@{}}
\toprule
Method & MT-Bench & Alpaca & Arena-Hard & CNN/DailyMail \\
\midrule
EAGLE-3 & \(32.65 \pm 0.18\;[3]\) & \(32.17 \pm 0.19\;[3]\) & \(32.62 \pm 0.29\;[3]\) & \(31.58 \pm 0.62\;[3]\) \\
P-EAGLE & \(26.49 \pm 0.068\;[3]\) & \(26.27 \pm 0.23\;[3]\) & \(26.73 \pm 0.065\;[3]\) & \(26.57 \pm 0.11\;[3]\) \\
DFlash & \(31.04 \pm 0.08\;[3]\) & \(30.94 \pm 0.11\;[3]\) & \(31.12 \pm 0.12\;[3]\) & \(31.19 \pm 0.1\;[3]\) \\
AR \textperiodcentered{} SGLang & \(23.22 \pm 0.056\;[3]\) & \(23.15 \pm 0.039\;[3]\) & \(23.25 \pm 0.031\;[3]\) & \(23.27 \pm 0.019\;[3]\) \\
AR \textperiodcentered{} vLLM ModelRunnerV1 & \(\mathbf{21.57} \pm 0.01\;[3]\) & \(\mathbf{21.49} \pm 0.026\;[3]\) & \(\mathbf{21.68} \pm 0.0078\;[3]\) & \(\mathbf{21.62} \pm 0.0038\;[3]\) \\
Pilot FFN 100\% & \(33.87 \pm 0.04\;[3]\) & \(32.53 \pm 0.16\;[3]\) & \(35.13 \pm 0.026\;[3]\) & \(33.95 \pm 0.038\;[3]\) \\
Pilot FFN 75\% & \(33.96 \pm 0.052\;[3]\) & \(32.53 \pm 0.086\;[3]\) & \(35.1 \pm 0.058\;[3]\) & \(34.04 \pm 0.039\;[3]\) \\
Pilot FFN 50\% & \(34.39 \pm 0.14\;[3]\) & \(33.21 \pm 0.037\;[3]\) & \(35.66 \pm 0.039\;[3]\) & \(34.52 \pm 0.071\;[3]\) \\
AR \textperiodcentered{} Transformers & \(25.37 \pm 0.058\;[3]\) & \(25.09 \pm 0.04\;[3]\) & \(25.49 \pm 0.019\;[3]\) & \(25.28 \pm 0.0051\;[3]\) \\
\bottomrule
\end{tabular}\par\medskip
\begin{minipage}{\linewidth}\footnotesize
Within each response, measure the gaps between successive nonempty chunks. Pool all observed gaps and report their p95. A response with more chunks contributes more gaps.
The bold mean is the best observed value in each benchmark column. The comparison covers the recorded implementations under their stated configurations. At least two eligible method entries are required. Ties at the displayed precision share bold formatting. Bold formatting conveys an observed comparison and makes no statistical-significance claim. Brackets show the available repetition count. Response differences remain part of the interpretation. Lower values are preferred.
Brackets give eligible repetition counts. Standard deviation appears when at least two repetitions are available. Each repetition uses three prompts. The evidence records preserve observation counts and measurement scopes. Empty cells lack an eligible measured value.
\end{minipage}
\end{table*}

\begin{table*}[!htp]
\centering\small
\setlength{\tabcolsep}{3pt}
\caption{Draft-token acceptance rate (\% of proposed draft tokens).}
\label{tab:appendix-metric-draft-fraction-none}
\textit{Mathematics}\par\smallskip
\begin{tabular}{@{}>{\raggedright\arraybackslash}p{3.7cm}rrr@{}}
\toprule
Method & GSM8K & MATH-500 & AIME 2025 \\
\midrule
EAGLE-3 & \(16.28 \pm 0\;[3]\) & \(15.37 \pm 0\;[3]\) & \(15.52 \pm 0\;[3]\) \\
P-EAGLE & \(30.71 \pm 0\;[3]\) & \(\mathbf{43.63} \pm 0\;[3]\) & \(\mathbf{45.7} \pm 0\;[3]\) \\
DFlash & \(\mathbf{62.85} \pm 0\;[3]\) & \(36.63 \pm 0\;[3]\) & \(42.25 \pm 0\;[3]\) \\
Pilot FFN 100\% & \(3.063 \pm 0\;[3]\) & \(1.191 \pm 0\;[3]\) & \(1.515 \pm 0\;[3]\) \\
Pilot FFN 75\% & \(3.097 \pm 0\;[3]\) & \(1.067 \pm 0\;[3]\) & \(1.657 \pm 0\;[3]\) \\
Pilot FFN 50\% & \(2.694 \pm 0\;[3]\) & \(0.9023 \pm 0\;[3]\) & \(1.322 \pm 0\;[3]\) \\
\bottomrule
\end{tabular}\par\medskip
\textit{Code}\par\smallskip
\begin{tabular}{@{}>{\raggedright\arraybackslash}p{3.7cm}rrr@{}}
\toprule
Method & HumanEval & MBPP & LiveCodeBench \\
\midrule
EAGLE-3 & \(15.77 \pm 0\;[3]\) & \(17.49 \pm 0\;[3]\) & \(14.79 \pm 0\;[3]\) \\
P-EAGLE & \(31.14 \pm 0\;[3]\) & \(33.8 \pm 0\;[3]\) & \(26.98 \pm 0\;[3]\) \\
DFlash & \(\mathbf{38.1} \pm 0\;[3]\) & \(\mathbf{43.43} \pm 0\;[3]\) & \(\mathbf{29.67} \pm 0\;[3]\) \\
Pilot FFN 100\% & \(1.984 \pm 0\;[3]\) & \(2.258 \pm 0\;[3]\) & \(1.962 \pm 0\;[3]\) \\
Pilot FFN 75\% & \(1.702 \pm 0\;[3]\) & \(1.629 \pm 0\;[3]\) & \(1.846 \pm 0\;[3]\) \\
Pilot FFN 50\% & \(2.042 \pm 0\;[3]\) & \(2.297 \pm 0\;[3]\) & \(1.8 \pm 0\;[3]\) \\
\bottomrule
\end{tabular}\par\medskip
\textit{Conversation, instruction following, and summarization}\par\smallskip
\begin{tabular}{@{}>{\raggedright\arraybackslash}p{3.7cm}rrrr@{}}
\toprule
Method & MT-Bench & Alpaca & Arena-Hard & CNN/DailyMail \\
\midrule
EAGLE-3 & \(17.05 \pm 0\;[3]\) & \(18.35 \pm 0\;[3]\) & \(13.63 \pm 0\;[3]\) & \(14.3 \pm 0\;[3]\) \\
P-EAGLE & \(16.86 \pm 0\;[3]\) & \(25.55 \pm 0\;[3]\) & \(21.09 \pm 0\;[3]\) & \(17.97 \pm 0\;[3]\) \\
DFlash & \(\mathbf{26.65} \pm 0\;[3]\) & \(\mathbf{29.55} \pm 0\;[3]\) & \(\mathbf{27.81} \pm 0\;[3]\) & \(\mathbf{23.16} \pm 0\;[3]\) \\
Pilot FFN 100\% & \(2.227 \pm 0\;[3]\) & \(1.977 \pm 0\;[3]\) & \(1.914 \pm 0\;[3]\) & \(1.922 \pm 0\;[3]\) \\
Pilot FFN 75\% & \(2.257 \pm 0\;[3]\) & \(2.051 \pm 0\;[3]\) & \(1.733 \pm 0\;[3]\) & \(2.034 \pm 0\;[3]\) \\
Pilot FFN 50\% & \(2.182 \pm 0\;[3]\) & \(1.879 \pm 0\;[3]\) & \(1.733 \pm 0\;[3]\) & \(1.701 \pm 0\;[3]\) \\
\bottomrule
\end{tabular}\par\medskip
\begin{minipage}{\linewidth}\footnotesize
Divide accepted draft tokens by proposed draft tokens and multiply by 100. The numerator and denominator count draft tokens only. Record target correction and bonus tokens separately. For a branching draft, state which branches the proposal count covers. The current EAGLE-3 setting proposes 15 candidate tree nodes per round. DFlash proposes nine linear draft tokens, and P-EAGLE proposes seven linear draft tokens. EAGLE-3's denominator includes candidates on alternative branches; its numerator counts the accepted path. Each linear method's denominator contains its proposed sequence. These different proposal structures limit comparisons of the acceptance percentages. The team drafter proposes eight linear tokens per round.
The bold mean is the best observed value in each benchmark column. The comparison covers the recorded implementations under their stated configurations. At least two eligible method entries are required. Ties at the displayed precision share bold formatting. Bold formatting conveys an observed comparison and makes no statistical-significance claim. Brackets show the available repetition count. Response differences remain part of the interpretation. Higher values are preferred. Tree and linear proposal denominators form separate comparison groups.
Brackets give eligible repetition counts. Standard deviation appears when at least two repetitions are available. Each repetition uses three prompts. The evidence records preserve observation counts and measurement scopes. Empty cells lack an eligible measured value.
\end{minipage}
\end{table*}

\begin{table*}[!htp]
\centering\small
\setlength{\tabcolsep}{3pt}
\caption{Accepted draft tokens per round (accepted draft tokens/round).}
\label{tab:appendix-metric-accepted-round-none}
\textit{Mathematics}\par\smallskip
\begin{tabular}{@{}>{\raggedright\arraybackslash}p{3.7cm}rrr@{}}
\toprule
Method & GSM8K & MATH-500 & AIME 2025 \\
\midrule
EAGLE-3 & \(2.443 \pm 0\;[3]\) & \(2.306 \pm 0\;[3]\) & \(2.328 \pm 0\;[3]\) \\
P-EAGLE & \(2.15 \pm 0\;[3]\) & \(3.054 \pm 0\;[3]\) & \(3.199 \pm 0\;[3]\) \\
DFlash & \(\mathbf{5.656} \pm 0\;[3]\) & \(\mathbf{3.297} \pm 0\;[3]\) & \(\mathbf{3.803} \pm 0\;[3]\) \\
Pilot FFN 100\% & \(0.245 \pm 0\;[3]\) & \(0.09528 \pm 0\;[3]\) & \(0.1212 \pm 0\;[3]\) \\
Pilot FFN 75\% & \(0.2478 \pm 0\;[3]\) & \(0.08539 \pm 0\;[3]\) & \(0.1326 \pm 0\;[3]\) \\
Pilot FFN 50\% & \(0.2155 \pm 0\;[3]\) & \(0.07218 \pm 0\;[3]\) & \(0.1057 \pm 0\;[3]\) \\
\bottomrule
\end{tabular}\par\medskip
\textit{Code}\par\smallskip
\begin{tabular}{@{}>{\raggedright\arraybackslash}p{3.7cm}rrr@{}}
\toprule
Method & HumanEval & MBPP & LiveCodeBench \\
\midrule
EAGLE-3 & \(2.365 \pm 0\;[3]\) & \(2.624 \pm 0\;[3]\) & \(2.218 \pm 0\;[3]\) \\
P-EAGLE & \(2.179 \pm 0\;[3]\) & \(2.366 \pm 0\;[3]\) & \(1.888 \pm 0\;[3]\) \\
DFlash & \(\mathbf{3.429} \pm 0\;[3]\) & \(\mathbf{3.909} \pm 0\;[3]\) & \(\mathbf{2.67} \pm 0\;[3]\) \\
Pilot FFN 100\% & \(0.1587 \pm 0\;[3]\) & \(0.1806 \pm 0\;[3]\) & \(0.157 \pm 0\;[3]\) \\
Pilot FFN 75\% & \(0.1362 \pm 0\;[3]\) & \(0.1303 \pm 0\;[3]\) & \(0.1477 \pm 0\;[3]\) \\
Pilot FFN 50\% & \(0.1633 \pm 0\;[3]\) & \(0.1837 \pm 0\;[3]\) & \(0.144 \pm 0\;[3]\) \\
\bottomrule
\end{tabular}\par\medskip
\textit{Conversation, instruction following, and summarization}\par\smallskip
\begin{tabular}{@{}>{\raggedright\arraybackslash}p{3.7cm}rrrr@{}}
\toprule
Method & MT-Bench & Alpaca & Arena-Hard & CNN/DailyMail \\
\midrule
EAGLE-3 & \(\mathbf{2.558} \pm 0\;[3]\) & \(\mathbf{2.753} \pm 0\;[3]\) & \(2.045 \pm 0\;[3]\) & \(\mathbf{2.145} \pm 0\;[3]\) \\
P-EAGLE & \(1.18 \pm 0\;[3]\) & \(1.788 \pm 0\;[3]\) & \(1.476 \pm 0\;[3]\) & \(1.258 \pm 0\;[3]\) \\
DFlash & \(2.399 \pm 0\;[3]\) & \(2.66 \pm 0\;[3]\) & \(\mathbf{2.503} \pm 0\;[3]\) & \(2.085 \pm 0\;[3]\) \\
Pilot FFN 100\% & \(0.1781 \pm 0\;[3]\) & \(0.1582 \pm 0\;[3]\) & \(0.1531 \pm 0\;[3]\) & \(0.1537 \pm 0\;[3]\) \\
Pilot FFN 75\% & \(0.1805 \pm 0\;[3]\) & \(0.1641 \pm 0\;[3]\) & \(0.1386 \pm 0\;[3]\) & \(0.1628 \pm 0\;[3]\) \\
Pilot FFN 50\% & \(0.1746 \pm 0\;[3]\) & \(0.1503 \pm 0\;[3]\) & \(0.1386 \pm 0\;[3]\) & \(0.1361 \pm 0\;[3]\) \\
\bottomrule
\end{tabular}\par\medskip
\begin{minipage}{\linewidth}\footnotesize
Divide the total accepted draft tokens by the number of verification rounds. Include rounds that accept zero draft tokens. Record tokens supplied by the target separately.
The bold mean is the best observed value in each benchmark column. The comparison covers the recorded implementations under their stated configurations. At least two eligible method entries are required. Ties at the displayed precision share bold formatting. Bold formatting conveys an observed comparison and makes no statistical-significance claim. Brackets show the available repetition count. Response differences remain part of the interpretation. Higher values are preferred.
Brackets give eligible repetition counts. Standard deviation appears when at least two repetitions are available. Each repetition uses three prompts. The evidence records preserve observation counts and measurement scopes. Empty cells lack an eligible measured value.
\end{minipage}
\end{table*}

\begin{table*}[!htp]
\centering\small
\setlength{\tabcolsep}{3pt}
\caption{Mean acceptance length (\(\tau\)), averaged across requests (committed tokens/round).}
\label{tab:appendix-metric-acceptance-length-request-none}
\textit{Mathematics}\par\smallskip
\begin{tabular}{@{}>{\raggedright\arraybackslash}p{3.7cm}rrr@{}}
\toprule
Method & GSM8K & MATH-500 & AIME 2025 \\
\midrule
EAGLE-3 & \(3.451 \pm 0\;[3]\) & \(3.301 \pm 0\;[3]\) & \(3.222 \pm 0\;[3]\) \\
DFlash & \(\mathbf{5.116} \pm 0\;[3]\) & \(\mathbf{4.25} \pm 0\;[3]\) & \(\mathbf{4.688} \pm 0\;[3]\) \\
Pilot FFN 100\% & \(1.229 \pm 0\;[3]\) & \(1.097 \pm 0\;[3]\) & \(1.121 \pm 0\;[3]\) \\
Pilot FFN 75\% & \(1.23 \pm 0\;[3]\) & \(1.091 \pm 0\;[3]\) & \(1.119 \pm 0\;[3]\) \\
Pilot FFN 50\% & \(1.2 \pm 0\;[3]\) & \(1.076 \pm 0\;[3]\) & \(1.1 \pm 0\;[3]\) \\
\bottomrule
\end{tabular}\par\medskip
\textit{Code}\par\smallskip
\begin{tabular}{@{}>{\raggedright\arraybackslash}p{3.7cm}rrr@{}}
\toprule
Method & HumanEval & MBPP & LiveCodeBench \\
\midrule
EAGLE-3 & \(3.337 \pm 0\;[3]\) & \(3.575 \pm 0\;[3]\) & \(3.436 \pm 0\;[3]\) \\
DFlash & \(\mathbf{4.702} \pm 0\;[3]\) & \(\mathbf{4.862} \pm 0\;[3]\) & \(\mathbf{4.164} \pm 0\;[3]\) \\
Pilot FFN 100\% & \(1.155 \pm 0\;[3]\) & \(1.178 \pm 0\;[3]\) & \(1.158 \pm 0\;[3]\) \\
Pilot FFN 75\% & \(1.133 \pm 0\;[3]\) & \(1.136 \pm 0\;[3]\) & \(1.15 \pm 0\;[3]\) \\
Pilot FFN 50\% & \(1.161 \pm 0\;[3]\) & \(1.18 \pm 0\;[3]\) & \(1.146 \pm 0\;[3]\) \\
\bottomrule
\end{tabular}\par\medskip
\textit{Conversation, instruction following, and summarization}\par\smallskip
\begin{tabular}{@{}>{\raggedright\arraybackslash}p{3.7cm}rrrr@{}}
\toprule
Method & MT-Bench & Alpaca & Arena-Hard & CNN/DailyMail \\
\midrule
EAGLE-3 & \(3.625 \pm 0\;[3]\) & \(3.74 \pm 0\;[3]\) & \(3.145 \pm 0\;[3]\) & \(\mathbf{3.158} \pm 0\;[3]\) \\
DFlash & \(\mathbf{3.665} \pm 0\;[3]\) & \(\mathbf{3.842} \pm 0\;[3]\) & \(\mathbf{3.808} \pm 0\;[3]\) & \(3.082 \pm 0\;[3]\) \\
Pilot FFN 100\% & \(1.166 \pm 0\;[3]\) & \(1.161 \pm 0\;[3]\) & \(1.163 \pm 0\;[3]\) & \(1.154 \pm 0\;[3]\) \\
Pilot FFN 75\% & \(1.167 \pm 0\;[3]\) & \(1.157 \pm 0\;[3]\) & \(1.152 \pm 0\;[3]\) & \(1.163 \pm 0\;[3]\) \\
Pilot FFN 50\% & \(1.152 \pm 0\;[3]\) & \(1.138 \pm 0\;[3]\) & \(1.138 \pm 0\;[3]\) & \(1.137 \pm 0\;[3]\) \\
\bottomrule
\end{tabular}\par\medskip
\begin{minipage}{\linewidth}\footnotesize
For each request with at least one verification round, divide its retained tokens by its round count. Include accepted drafts and retained target tokens before final output trimming. Average the request values equally.
The bold mean is the best observed value in each benchmark column. The comparison covers the recorded implementations under their stated configurations. At least two eligible method entries are required. Ties at the displayed precision share bold formatting. Bold formatting conveys an observed comparison and makes no statistical-significance claim. Brackets show the available repetition count. Response differences remain part of the interpretation. Higher values are preferred.
Brackets give eligible repetition counts. Standard deviation appears when at least two repetitions are available. Each repetition uses three prompts. The evidence records preserve observation counts and measurement scopes. Empty cells lack an eligible measured value.
\end{minipage}
\end{table*}

\begin{table*}[!htp]
\centering\small
\setlength{\tabcolsep}{3pt}
\caption{Delivered output tokens per round (delivered tokens/round).}
\label{tab:appendix-metric-returned-round-none}
\textit{Mathematics}\par\smallskip
\begin{tabular}{@{}>{\raggedright\arraybackslash}p{3.7cm}rrr@{}}
\toprule
Method & GSM8K & MATH-500 & AIME 2025 \\
\midrule
EAGLE-3 & \(3.408 \pm 0\;[3]\) & \(3.293 \pm 0\;[3]\) & \(3.324 \pm 0\;[3]\) \\
DFlash & \(\mathbf{6.625} \pm 0\;[3]\) & \(\mathbf{4.287} \pm 0\;[3]\) & \(\mathbf{4.797} \pm 0\;[3]\) \\
Pilot FFN 100\% & \(1.238 \pm 0\;[3]\) & \(1.093 \pm 0\;[3]\) & \(1.121 \pm 0\;[3]\) \\
Pilot FFN 75\% & \(1.241 \pm 0\;[3]\) & \(1.083 \pm 0\;[3]\) & \(1.132 \pm 0\;[3]\) \\
Pilot FFN 50\% & \(1.209 \pm 0\;[3]\) & \(1.07 \pm 0\;[3]\) & \(1.105 \pm 0\;[3]\) \\
\bottomrule
\end{tabular}\par\medskip
\textit{Code}\par\smallskip
\begin{tabular}{@{}>{\raggedright\arraybackslash}p{3.7cm}rrr@{}}
\toprule
Method & HumanEval & MBPP & LiveCodeBench \\
\midrule
EAGLE-3 & \(3.311 \pm 0\;[3]\) & \(3.584 \pm 0\;[3]\) & \(3.197 \pm 0\;[3]\) \\
DFlash & \(\mathbf{4.375} \pm 0\;[3]\) & \(\mathbf{4.875} \pm 0\;[3]\) & \(\mathbf{3.655} \pm 0\;[3]\) \\
Pilot FFN 100\% & \(1.147 \pm 0\;[3]\) & \(1.173 \pm 0\;[3]\) & \(1.152 \pm 0\;[3]\) \\
Pilot FFN 75\% & \(1.125 \pm 0\;[3]\) & \(1.123 \pm 0\;[3]\) & \(1.143 \pm 0\;[3]\) \\
Pilot FFN 50\% & \(1.151 \pm 0\;[3]\) & \(1.176 \pm 0\;[3]\) & \(1.139 \pm 0\;[3]\) \\
\bottomrule
\end{tabular}\par\medskip
\textit{Conversation, instruction following, and summarization}\par\smallskip
\begin{tabular}{@{}>{\raggedright\arraybackslash}p{3.7cm}rrrr@{}}
\toprule
Method & MT-Bench & Alpaca & Arena-Hard & CNN/DailyMail \\
\midrule
EAGLE-3 & \(\mathbf{3.546} \pm 0\;[3]\) & \(\mathbf{3.727} \pm 0\;[3]\) & \(3.028 \pm 0\;[3]\) & \(\mathbf{3.12} \pm 0\;[3]\) \\
DFlash & \(3.39 \pm 0\;[3]\) & \(3.644 \pm 0\;[3]\) & \(\mathbf{3.493} \pm 0\;[3]\) & \(3.072 \pm 0\;[3]\) \\
Pilot FFN 100\% & \(1.175 \pm 0\;[3]\) & \(1.153 \pm 0\;[3]\) & \(1.149 \pm 0\;[3]\) & \(1.149 \pm 0\;[3]\) \\
Pilot FFN 75\% & \(1.177 \pm 0\;[3]\) & \(1.159 \pm 0\;[3]\) & \(1.135 \pm 0\;[3]\) & \(1.158 \pm 0\;[3]\) \\
Pilot FFN 50\% & \(1.172 \pm 0\;[3]\) & \(1.145 \pm 0\;[3]\) & \(1.135 \pm 0\;[3]\) & \(1.131 \pm 0\;[3]\) \\
\bottomrule
\end{tabular}\par\medskip
\begin{minipage}{\linewidth}\footnotesize
Count the output tokens from the recorded rounds that actually reach the client after stopping and output limits. Divide by all those rounds, including rounds that return zero tokens. Record any initial target token outside the rounds separately.
The bold mean is the best observed value in each benchmark column. The comparison covers the recorded implementations under their stated configurations. At least two eligible method entries are required. Ties at the displayed precision share bold formatting. Bold formatting conveys an observed comparison and makes no statistical-significance claim. Brackets show the available repetition count. Response differences remain part of the interpretation. Higher values are preferred.
Brackets give eligible repetition counts. Standard deviation appears when at least two repetitions are available. Each repetition uses three prompts. The evidence records preserve observation counts and measurement scopes. Empty cells lack an eligible measured value.
\end{minipage}
\end{table*}

\begin{table*}[!htp]
\centering\small
\setlength{\tabcolsep}{3pt}
\caption{Prefix acceptance probability by position; Position 1 (probability at each draft position).}
\label{tab:appendix-metric-prefix-survival-1}
\textit{Mathematics}\par\smallskip
\begin{tabular}{@{}>{\raggedright\arraybackslash}p{3.7cm}rrr@{}}
\toprule
Method & GSM8K & MATH-500 & AIME 2025 \\
\midrule
DFlash & \(\mathbf{0.9198} \pm 0\;[3]\) & \(\mathbf{0.8594} \pm 0\;[3]\) & \(\mathbf{0.844} \pm 0\;[3]\) \\
Pilot FFN 100\% & \(0.245 \pm 0\;[3]\) & \(0.09528 \pm 0\;[3]\) & \(0.121 \pm 0\;[3]\) \\
Pilot FFN 75\% & \(0.2434 \pm 0\;[3]\) & \(0.08539 \pm 0\;[3]\) & \(0.1301 \pm 0\;[3]\) \\
Pilot FFN 50\% & \(0.2091 \pm 0\;[3]\) & \(0.07218 \pm 0\;[3]\) & \(0.105 \pm 0\;[3]\) \\
\bottomrule
\end{tabular}\par\medskip
\textit{Code}\par\smallskip
\begin{tabular}{@{}>{\raggedright\arraybackslash}p{3.7cm}rrr@{}}
\toprule
Method & HumanEval & MBPP & LiveCodeBench \\
\midrule
DFlash & \(\mathbf{0.8036} \pm 0\;[3]\) & \(\mathbf{0.8068} \pm 0\;[3]\) & \(\mathbf{0.79} \pm 0\;[3]\) \\
Pilot FFN 100\% & \(0.1587 \pm 0\;[3]\) & \(0.178 \pm 0\;[3]\) & \(0.1489 \pm 0\;[3]\) \\
Pilot FFN 75\% & \(0.1362 \pm 0\;[3]\) & \(0.1303 \pm 0\;[3]\) & \(0.1413 \pm 0\;[3]\) \\
Pilot FFN 50\% & \(0.1633 \pm 0\;[3]\) & \(0.1811 \pm 0\;[3]\) & \(0.136 \pm 0\;[3]\) \\
\bottomrule
\end{tabular}\par\medskip
\textit{Conversation, instruction following, and summarization}\par\smallskip
\begin{tabular}{@{}>{\raggedright\arraybackslash}p{3.7cm}rrrr@{}}
\toprule
Method & MT-Bench & Alpaca & Arena-Hard & CNN/DailyMail \\
\midrule
DFlash & \(\mathbf{0.7977} \pm 0\;[3]\) & \(\mathbf{0.8144} \pm 0\;[3]\) & \(\mathbf{0.7222} \pm 0\;[3]\) & \(\mathbf{0.7712} \pm 0\;[3]\) \\
Pilot FFN 100\% & \(0.1781 \pm 0\;[3]\) & \(0.1582 \pm 0\;[3]\) & \(0.1456 \pm 0\;[3]\) & \(0.1537 \pm 0\;[3]\) \\
Pilot FFN 75\% & \(0.1805 \pm 0\;[3]\) & \(0.1573 \pm 0\;[3]\) & \(0.1302 \pm 0\;[3]\) & \(0.1612 \pm 0\;[3]\) \\
Pilot FFN 50\% & \(0.1746 \pm 0\;[3]\) & \(0.1503 \pm 0\;[3]\) & \(0.1333 \pm 0\;[3]\) & \(0.1346 \pm 0\;[3]\) \\
\bottomrule
\end{tabular}\par\medskip
\begin{minipage}{\linewidth}\footnotesize
For each draft position, count the rounds that proposed that position. Among those rounds, count how often the target accepted every draft token through that position before final output clipping. Divide the second count by the first.
The bold mean is the best observed value in each benchmark column. The comparison covers the recorded implementations under their stated configurations. At least two eligible method entries are required. Ties at the displayed precision share bold formatting. Bold formatting conveys an observed comparison and makes no statistical-significance claim. Brackets show the available repetition count. Response differences remain part of the interpretation. Higher values are preferred. Tree and linear proposal denominators form separate comparison groups.
Brackets give eligible repetition counts. Standard deviation appears when at least two repetitions are available. Each repetition uses three prompts. The evidence records preserve observation counts and measurement scopes. Empty cells lack an eligible measured value.
\end{minipage}
\end{table*}

\begin{table*}[!htp]
\centering\small
\setlength{\tabcolsep}{3pt}
\caption{Prefix acceptance probability by position; Position 2 (probability at each draft position).}
\label{tab:appendix-metric-prefix-survival-2}
\textit{Mathematics}\par\smallskip
\begin{tabular}{@{}>{\raggedright\arraybackslash}p{3.7cm}rrr@{}}
\toprule
Method & GSM8K & MATH-500 & AIME 2025 \\
\midrule
DFlash & \(\mathbf{0.8252} \pm 0\;[3]\) & \(\mathbf{0.6813} \pm 0\;[3]\) & \(\mathbf{0.6663} \pm 0\;[3]\) \\
Pilot FFN 100\% & \(0 \pm 0\;[3]\) & \(0 \pm 0\;[3]\) & \(0.0002453 \pm 0\;[3]\) \\
Pilot FFN 75\% & \(0.004425 \pm 0\;[3]\) & \(0 \pm 0\;[3]\) & \(0.002478 \pm 0\;[3]\) \\
Pilot FFN 50\% & \(0.006466 \pm 0\;[3]\) & \(0 \pm 0\;[3]\) & \(0.0007259 \pm 0\;[3]\) \\
\bottomrule
\end{tabular}\par\medskip
\textit{Code}\par\smallskip
\begin{tabular}{@{}>{\raggedright\arraybackslash}p{3.7cm}rrr@{}}
\toprule
Method & HumanEval & MBPP & LiveCodeBench \\
\midrule
DFlash & \(\mathbf{0.6964} \pm 0\;[3]\) & \(\mathbf{0.6818} \pm 0\;[3]\) & \(\mathbf{0.605} \pm 0\;[3]\) \\
Pilot FFN 100\% & \(0 \pm 0\;[3]\) & \(0.002618 \pm 0\;[3]\) & \(0.008091 \pm 0\;[3]\) \\
Pilot FFN 75\% & \(0 \pm 0\;[3]\) & \(0 \pm 0\;[3]\) & \(0.006421 \pm 0\;[3]\) \\
Pilot FFN 50\% & \(0 \pm 0\;[3]\) & \(0.002625 \pm 0\;[3]\) & \(0.008 \pm 0\;[3]\) \\
\bottomrule
\end{tabular}\par\medskip
\textit{Conversation, instruction following, and summarization}\par\smallskip
\begin{tabular}{@{}>{\raggedright\arraybackslash}p{3.7cm}rrrr@{}}
\toprule
Method & MT-Bench & Alpaca & Arena-Hard & CNN/DailyMail \\
\midrule
DFlash & \(\mathbf{0.5897} \pm 0\;[3]\) & \(\mathbf{0.634} \pm 0\;[3]\) & \(\mathbf{0.5588} \pm 0\;[3]\) & \(\mathbf{0.5339} \pm 0\;[3]\) \\
Pilot FFN 100\% & \(0 \pm 0\;[3]\) & \(0 \pm 0\;[3]\) & \(0.007495 \pm 0\;[3]\) & \(0 \pm 0\;[3]\) \\
Pilot FFN 75\% & \(0 \pm 0\;[3]\) & \(0.006838 \pm 0\;[3]\) & \(0.008466 \pm 0\;[3]\) & \(0.001565 \pm 0\;[3]\) \\
Pilot FFN 50\% & \(0 \pm 0\;[3]\) & \(0 \pm 0\;[3]\) & \(0.005291 \pm 0\;[3]\) & \(0.001529 \pm 0\;[3]\) \\
\bottomrule
\end{tabular}\par\medskip
\begin{minipage}{\linewidth}\footnotesize
For each draft position, count the rounds that proposed that position. Among those rounds, count how often the target accepted every draft token through that position before final output clipping. Divide the second count by the first.
The bold mean is the best observed value in each benchmark column. The comparison covers the recorded implementations under their stated configurations. At least two eligible method entries are required. Ties at the displayed precision share bold formatting. Bold formatting conveys an observed comparison and makes no statistical-significance claim. Brackets show the available repetition count. Response differences remain part of the interpretation. Higher values are preferred. Tree and linear proposal denominators form separate comparison groups.
Brackets give eligible repetition counts. Standard deviation appears when at least two repetitions are available. Each repetition uses three prompts. The evidence records preserve observation counts and measurement scopes. Empty cells lack an eligible measured value.
\end{minipage}
\end{table*}

\begin{table*}[!htp]
\centering\small
\setlength{\tabcolsep}{3pt}
\caption{Prefix acceptance probability by position; Position 3 (probability at each draft position).}
\label{tab:appendix-metric-prefix-survival-3}
\textit{Mathematics}\par\smallskip
\begin{tabular}{@{}>{\raggedright\arraybackslash}p{3.7cm}rrr@{}}
\toprule
Method & GSM8K & MATH-500 & AIME 2025 \\
\midrule
DFlash & \(\mathbf{0.745} \pm 0\;[3]\) & \(\mathbf{0.4906} \pm 0\;[3]\) & \(\mathbf{0.5333} \pm 0\;[3]\) \\
Pilot FFN 100\% & \(0 \pm 0\;[3]\) & \(0 \pm 0\;[3]\) & \(0 \pm 0\;[3]\) \\
Pilot FFN 75\% & \(0 \pm 0\;[3]\) & \(0 \pm 0\;[3]\) & \(0 \pm 0\;[3]\) \\
Pilot FFN 50\% & \(0 \pm 0\;[3]\) & \(0 \pm 0\;[3]\) & \(0 \pm 0\;[3]\) \\
\bottomrule
\end{tabular}\par\medskip
\textit{Code}\par\smallskip
\begin{tabular}{@{}>{\raggedright\arraybackslash}p{3.7cm}rrr@{}}
\toprule
Method & HumanEval & MBPP & LiveCodeBench \\
\midrule
DFlash & \(\mathbf{0.5357} \pm 0\;[3]\) & \(\mathbf{0.5227} \pm 0\;[3]\) & \(\mathbf{0.415} \pm 0\;[3]\) \\
Pilot FFN 100\% & \(0 \pm 0\;[3]\) & \(0 \pm 0\;[3]\) & \(0 \pm 0\;[3]\) \\
Pilot FFN 75\% & \(0 \pm 0\;[3]\) & \(0 \pm 0\;[3]\) & \(0 \pm 0\;[3]\) \\
Pilot FFN 50\% & \(0 \pm 0\;[3]\) & \(0 \pm 0\;[3]\) & \(0 \pm 0\;[3]\) \\
\bottomrule
\end{tabular}\par\medskip
\textit{Conversation, instruction following, and summarization}\par\smallskip
\begin{tabular}{@{}>{\raggedright\arraybackslash}p{3.7cm}rrrr@{}}
\toprule
Method & MT-Bench & Alpaca & Arena-Hard & CNN/DailyMail \\
\midrule
DFlash & \(\mathbf{0.4046} \pm 0\;[3]\) & \(\mathbf{0.4433} \pm 0\;[3]\) & \(\mathbf{0.4183} \pm 0\;[3]\) & \(\mathbf{0.3263} \pm 0\;[3]\) \\
Pilot FFN 100\% & \(0 \pm 0\;[3]\) & \(0 \pm 0\;[3]\) & \(0 \pm 0\;[3]\) & \(0 \pm 0\;[3]\) \\
Pilot FFN 75\% & \(0 \pm 0\;[3]\) & \(0 \pm 0\;[3]\) & \(0 \pm 0\;[3]\) & \(0 \pm 0\;[3]\) \\
Pilot FFN 50\% & \(0 \pm 0\;[3]\) & \(0 \pm 0\;[3]\) & \(0 \pm 0\;[3]\) & \(0 \pm 0\;[3]\) \\
\bottomrule
\end{tabular}\par\medskip
\begin{minipage}{\linewidth}\footnotesize
For each draft position, count the rounds that proposed that position. Among those rounds, count how often the target accepted every draft token through that position before final output clipping. Divide the second count by the first.
The bold mean is the best observed value in each benchmark column. The comparison covers the recorded implementations under their stated configurations. At least two eligible method entries are required. Ties at the displayed precision share bold formatting. Bold formatting conveys an observed comparison and makes no statistical-significance claim. Brackets show the available repetition count. Response differences remain part of the interpretation. Higher values are preferred. Tree and linear proposal denominators form separate comparison groups.
Brackets give eligible repetition counts. Standard deviation appears when at least two repetitions are available. Each repetition uses three prompts. The evidence records preserve observation counts and measurement scopes. Empty cells lack an eligible measured value.
\end{minipage}
\end{table*}

\begin{table*}[!htp]
\centering\small
\setlength{\tabcolsep}{3pt}
\caption{Prefix acceptance probability by position; Position 4 (probability at each draft position).}
\label{tab:appendix-metric-prefix-survival-4}
\textit{Mathematics}\par\smallskip
\begin{tabular}{@{}>{\raggedright\arraybackslash}p{3.7cm}rrr@{}}
\toprule
Method & GSM8K & MATH-500 & AIME 2025 \\
\midrule
DFlash & \(\mathbf{0.6447} \pm 0\;[3]\) & \(\mathbf{0.3812} \pm 0\;[3]\) & \(\mathbf{0.4472} \pm 0\;[3]\) \\
Pilot FFN 100\% & \(0 \pm 0\;[3]\) & \(0 \pm 0\;[3]\) & \(0 \pm 0\;[3]\) \\
Pilot FFN 75\% & \(0 \pm 0\;[3]\) & \(0 \pm 0\;[3]\) & \(0 \pm 0\;[3]\) \\
Pilot FFN 50\% & \(0 \pm 0\;[3]\) & \(0 \pm 0\;[3]\) & \(0 \pm 0\;[3]\) \\
\bottomrule
\end{tabular}\par\medskip
\textit{Code}\par\smallskip
\begin{tabular}{@{}>{\raggedright\arraybackslash}p{3.7cm}rrr@{}}
\toprule
Method & HumanEval & MBPP & LiveCodeBench \\
\midrule
DFlash & \(\mathbf{0.4107} \pm 0\;[3]\) & \(\mathbf{0.4659} \pm 0\;[3]\) & \(\mathbf{0.285} \pm 0\;[3]\) \\
Pilot FFN 100\% & \(0 \pm 0\;[3]\) & \(0 \pm 0\;[3]\) & \(0 \pm 0\;[3]\) \\
Pilot FFN 75\% & \(0 \pm 0\;[3]\) & \(0 \pm 0\;[3]\) & \(0 \pm 0\;[3]\) \\
Pilot FFN 50\% & \(0 \pm 0\;[3]\) & \(0 \pm 0\;[3]\) & \(0 \pm 0\;[3]\) \\
\bottomrule
\end{tabular}\par\medskip
\textit{Conversation, instruction following, and summarization}\par\smallskip
\begin{tabular}{@{}>{\raggedright\arraybackslash}p{3.7cm}rrrr@{}}
\toprule
Method & MT-Bench & Alpaca & Arena-Hard & CNN/DailyMail \\
\midrule
DFlash & \(\mathbf{0.2621} \pm 0\;[3]\) & \(\mathbf{0.299} \pm 0\;[3]\) & \(\mathbf{0.2941} \pm 0\;[3]\) & \(\mathbf{0.1992} \pm 0\;[3]\) \\
Pilot FFN 100\% & \(0 \pm 0\;[3]\) & \(0 \pm 0\;[3]\) & \(0 \pm 0\;[3]\) & \(0 \pm 0\;[3]\) \\
Pilot FFN 75\% & \(0 \pm 0\;[3]\) & \(0 \pm 0\;[3]\) & \(0 \pm 0\;[3]\) & \(0 \pm 0\;[3]\) \\
Pilot FFN 50\% & \(0 \pm 0\;[3]\) & \(0 \pm 0\;[3]\) & \(0 \pm 0\;[3]\) & \(0 \pm 0\;[3]\) \\
\bottomrule
\end{tabular}\par\medskip
\begin{minipage}{\linewidth}\footnotesize
For each draft position, count the rounds that proposed that position. Among those rounds, count how often the target accepted every draft token through that position before final output clipping. Divide the second count by the first.
The bold mean is the best observed value in each benchmark column. The comparison covers the recorded implementations under their stated configurations. At least two eligible method entries are required. Ties at the displayed precision share bold formatting. Bold formatting conveys an observed comparison and makes no statistical-significance claim. Brackets show the available repetition count. Response differences remain part of the interpretation. Higher values are preferred. Tree and linear proposal denominators form separate comparison groups.
Brackets give eligible repetition counts. Standard deviation appears when at least two repetitions are available. Each repetition uses three prompts. The evidence records preserve observation counts and measurement scopes. Empty cells lack an eligible measured value.
\end{minipage}
\end{table*}

\begin{table*}[!htp]
\centering\small
\setlength{\tabcolsep}{3pt}
\caption{Prefix acceptance probability by position; Position 5 (probability at each draft position).}
\label{tab:appendix-metric-prefix-survival-5}
\textit{Mathematics}\par\smallskip
\begin{tabular}{@{}>{\raggedright\arraybackslash}p{3.7cm}rrr@{}}
\toprule
Method & GSM8K & MATH-500 & AIME 2025 \\
\midrule
DFlash & \(\mathbf{0.5874} \pm 0\;[3]\) & \(\mathbf{0.3031} \pm 0\;[3]\) & \(\mathbf{0.3555} \pm 0\;[3]\) \\
Pilot FFN 100\% & \(0 \pm 0\;[3]\) & \(0 \pm 0\;[3]\) & \(0 \pm 0\;[3]\) \\
Pilot FFN 75\% & \(0 \pm 0\;[3]\) & \(0 \pm 0\;[3]\) & \(0 \pm 0\;[3]\) \\
Pilot FFN 50\% & \(0 \pm 0\;[3]\) & \(0 \pm 0\;[3]\) & \(0 \pm 0\;[3]\) \\
\bottomrule
\end{tabular}\par\medskip
\textit{Code}\par\smallskip
\begin{tabular}{@{}>{\raggedright\arraybackslash}p{3.7cm}rrr@{}}
\toprule
Method & HumanEval & MBPP & LiveCodeBench \\
\midrule
DFlash & \(\mathbf{0.3393} \pm 0\;[3]\) & \(\mathbf{0.4091} \pm 0\;[3]\) & \(\mathbf{0.21} \pm 0\;[3]\) \\
Pilot FFN 100\% & \(0 \pm 0\;[3]\) & \(0 \pm 0\;[3]\) & \(0 \pm 0\;[3]\) \\
Pilot FFN 75\% & \(0 \pm 0\;[3]\) & \(0 \pm 0\;[3]\) & \(0 \pm 0\;[3]\) \\
Pilot FFN 50\% & \(0 \pm 0\;[3]\) & \(0 \pm 0\;[3]\) & \(0 \pm 0\;[3]\) \\
\bottomrule
\end{tabular}\par\medskip
\textit{Conversation, instruction following, and summarization}\par\smallskip
\begin{tabular}{@{}>{\raggedright\arraybackslash}p{3.7cm}rrrr@{}}
\toprule
Method & MT-Bench & Alpaca & Arena-Hard & CNN/DailyMail \\
\midrule
DFlash & \(\mathbf{0.1709} \pm 0\;[3]\) & \(\mathbf{0.1804} \pm 0\;[3]\) & \(\mathbf{0.1928} \pm 0\;[3]\) & \(\mathbf{0.1059} \pm 0\;[3]\) \\
Pilot FFN 100\% & \(0 \pm 0\;[3]\) & \(0 \pm 0\;[3]\) & \(0 \pm 0\;[3]\) & \(0 \pm 0\;[3]\) \\
Pilot FFN 75\% & \(0 \pm 0\;[3]\) & \(0 \pm 0\;[3]\) & \(0 \pm 0\;[3]\) & \(0 \pm 0\;[3]\) \\
Pilot FFN 50\% & \(0 \pm 0\;[3]\) & \(0 \pm 0\;[3]\) & \(0 \pm 0\;[3]\) & \(0 \pm 0\;[3]\) \\
\bottomrule
\end{tabular}\par\medskip
\begin{minipage}{\linewidth}\footnotesize
For each draft position, count the rounds that proposed that position. Among those rounds, count how often the target accepted every draft token through that position before final output clipping. Divide the second count by the first.
The bold mean is the best observed value in each benchmark column. The comparison covers the recorded implementations under their stated configurations. At least two eligible method entries are required. Ties at the displayed precision share bold formatting. Bold formatting conveys an observed comparison and makes no statistical-significance claim. Brackets show the available repetition count. Response differences remain part of the interpretation. Higher values are preferred. Tree and linear proposal denominators form separate comparison groups.
Brackets give eligible repetition counts. Standard deviation appears when at least two repetitions are available. Each repetition uses three prompts. The evidence records preserve observation counts and measurement scopes. Empty cells lack an eligible measured value.
\end{minipage}
\end{table*}

\begin{table*}[!htp]
\centering\small
\setlength{\tabcolsep}{3pt}
\caption{Prefix acceptance probability by position; Position 6 (probability at each draft position).}
\label{tab:appendix-metric-prefix-survival-6}
\textit{Mathematics}\par\smallskip
\begin{tabular}{@{}>{\raggedright\arraybackslash}p{3.7cm}rrr@{}}
\toprule
Method & GSM8K & MATH-500 & AIME 2025 \\
\midrule
DFlash & \(\mathbf{0.5301} \pm 0\;[3]\) & \(\mathbf{0.2281} \pm 0\;[3]\) & \(\mathbf{0.2959} \pm 0\;[3]\) \\
Pilot FFN 100\% & \(0 \pm 0\;[3]\) & \(0 \pm 0\;[3]\) & \(0 \pm 0\;[3]\) \\
Pilot FFN 75\% & \(0 \pm 0\;[3]\) & \(0 \pm 0\;[3]\) & \(0 \pm 0\;[3]\) \\
Pilot FFN 50\% & \(0 \pm 0\;[3]\) & \(0 \pm 0\;[3]\) & \(0 \pm 0\;[3]\) \\
\bottomrule
\end{tabular}\par\medskip
\textit{Code}\par\smallskip
\begin{tabular}{@{}>{\raggedright\arraybackslash}p{3.7cm}rrr@{}}
\toprule
Method & HumanEval & MBPP & LiveCodeBench \\
\midrule
DFlash & \(\mathbf{0.2679} \pm 0\;[3]\) & \(\mathbf{0.3182} \pm 0\;[3]\) & \(\mathbf{0.145} \pm 0\;[3]\) \\
Pilot FFN 100\% & \(0 \pm 0\;[3]\) & \(0 \pm 0\;[3]\) & \(0 \pm 0\;[3]\) \\
Pilot FFN 75\% & \(0 \pm 0\;[3]\) & \(0 \pm 0\;[3]\) & \(0 \pm 0\;[3]\) \\
Pilot FFN 50\% & \(0 \pm 0\;[3]\) & \(0 \pm 0\;[3]\) & \(0 \pm 0\;[3]\) \\
\bottomrule
\end{tabular}\par\medskip
\textit{Conversation, instruction following, and summarization}\par\smallskip
\begin{tabular}{@{}>{\raggedright\arraybackslash}p{3.7cm}rrrr@{}}
\toprule
Method & MT-Bench & Alpaca & Arena-Hard & CNN/DailyMail \\
\midrule
DFlash & \(\mathbf{0.09117} \pm 0\;[3]\) & \(\mathbf{0.1082} \pm 0\;[3]\) & \(\mathbf{0.1209} \pm 0\;[3]\) & \(\mathbf{0.0678} \pm 0\;[3]\) \\
Pilot FFN 100\% & \(0 \pm 0\;[3]\) & \(0 \pm 0\;[3]\) & \(0 \pm 0\;[3]\) & \(0 \pm 0\;[3]\) \\
Pilot FFN 75\% & \(0 \pm 0\;[3]\) & \(0 \pm 0\;[3]\) & \(0 \pm 0\;[3]\) & \(0 \pm 0\;[3]\) \\
Pilot FFN 50\% & \(0 \pm 0\;[3]\) & \(0 \pm 0\;[3]\) & \(0 \pm 0\;[3]\) & \(0 \pm 0\;[3]\) \\
\bottomrule
\end{tabular}\par\medskip
\begin{minipage}{\linewidth}\footnotesize
For each draft position, count the rounds that proposed that position. Among those rounds, count how often the target accepted every draft token through that position before final output clipping. Divide the second count by the first.
The bold mean is the best observed value in each benchmark column. The comparison covers the recorded implementations under their stated configurations. At least two eligible method entries are required. Ties at the displayed precision share bold formatting. Bold formatting conveys an observed comparison and makes no statistical-significance claim. Brackets show the available repetition count. Response differences remain part of the interpretation. Higher values are preferred. Tree and linear proposal denominators form separate comparison groups.
Brackets give eligible repetition counts. Standard deviation appears when at least two repetitions are available. Each repetition uses three prompts. The evidence records preserve observation counts and measurement scopes. Empty cells lack an eligible measured value.
\end{minipage}
\end{table*}

\begin{table*}[!htp]
\centering\small
\setlength{\tabcolsep}{3pt}
\caption{Prefix acceptance probability by position; Position 7 (probability at each draft position).}
\label{tab:appendix-metric-prefix-survival-7}
\textit{Mathematics}\par\smallskip
\begin{tabular}{@{}>{\raggedright\arraybackslash}p{3.7cm}rrr@{}}
\toprule
Method & GSM8K & MATH-500 & AIME 2025 \\
\midrule
DFlash & \(\mathbf{0.5014} \pm 0\;[3]\) & \(\mathbf{0.1688} \pm 0\;[3]\) & \(\mathbf{0.2603} \pm 0\;[3]\) \\
Pilot FFN 100\% & \(0 \pm 0\;[3]\) & \(0 \pm 0\;[3]\) & \(0 \pm 0\;[3]\) \\
Pilot FFN 75\% & \(0 \pm 0\;[3]\) & \(0 \pm 0\;[3]\) & \(0 \pm 0\;[3]\) \\
Pilot FFN 50\% & \(0 \pm 0\;[3]\) & \(0 \pm 0\;[3]\) & \(0 \pm 0\;[3]\) \\
\bottomrule
\end{tabular}\par\medskip
\textit{Code}\par\smallskip
\begin{tabular}{@{}>{\raggedright\arraybackslash}p{3.7cm}rrr@{}}
\toprule
Method & HumanEval & MBPP & LiveCodeBench \\
\midrule
DFlash & \(\mathbf{0.1786} \pm 0\;[3]\) & \(\mathbf{0.2841} \pm 0\;[3]\) & \(\mathbf{0.09} \pm 0\;[3]\) \\
Pilot FFN 100\% & \(0 \pm 0\;[3]\) & \(0 \pm 0\;[3]\) & \(0 \pm 0\;[3]\) \\
Pilot FFN 75\% & \(0 \pm 0\;[3]\) & \(0 \pm 0\;[3]\) & \(0 \pm 0\;[3]\) \\
Pilot FFN 50\% & \(0 \pm 0\;[3]\) & \(0 \pm 0\;[3]\) & \(0 \pm 0\;[3]\) \\
\bottomrule
\end{tabular}\par\medskip
\textit{Conversation, instruction following, and summarization}\par\smallskip
\begin{tabular}{@{}>{\raggedright\arraybackslash}p{3.7cm}rrrr@{}}
\toprule
Method & MT-Bench & Alpaca & Arena-Hard & CNN/DailyMail \\
\midrule
DFlash & \(\mathbf{0.04274} \pm 0\;[3]\) & \(\mathbf{0.08247} \pm 0\;[3]\) & \(\mathbf{0.08824} \pm 0\;[3]\) & \(\mathbf{0.03814} \pm 0\;[3]\) \\
Pilot FFN 100\% & \(0 \pm 0\;[3]\) & \(0 \pm 0\;[3]\) & \(0 \pm 0\;[3]\) & \(0 \pm 0\;[3]\) \\
Pilot FFN 75\% & \(0 \pm 0\;[3]\) & \(0 \pm 0\;[3]\) & \(0 \pm 0\;[3]\) & \(0 \pm 0\;[3]\) \\
Pilot FFN 50\% & \(0 \pm 0\;[3]\) & \(0 \pm 0\;[3]\) & \(0 \pm 0\;[3]\) & \(0 \pm 0\;[3]\) \\
\bottomrule
\end{tabular}\par\medskip
\begin{minipage}{\linewidth}\footnotesize
For each draft position, count the rounds that proposed that position. Among those rounds, count how often the target accepted every draft token through that position before final output clipping. Divide the second count by the first.
The bold mean is the best observed value in each benchmark column. The comparison covers the recorded implementations under their stated configurations. At least two eligible method entries are required. Ties at the displayed precision share bold formatting. Bold formatting conveys an observed comparison and makes no statistical-significance claim. Brackets show the available repetition count. Response differences remain part of the interpretation. Higher values are preferred. Tree and linear proposal denominators form separate comparison groups.
Brackets give eligible repetition counts. Standard deviation appears when at least two repetitions are available. Each repetition uses three prompts. The evidence records preserve observation counts and measurement scopes. Empty cells lack an eligible measured value.
\end{minipage}
\end{table*}

\begin{table*}[!htp]
\centering\small
\setlength{\tabcolsep}{3pt}
\caption{Prefix acceptance probability by position; Position 8 (probability at each draft position).}
\label{tab:appendix-metric-prefix-survival-8}
\textit{Mathematics}\par\smallskip
\begin{tabular}{@{}>{\raggedright\arraybackslash}p{3.7cm}rrr@{}}
\toprule
Method & GSM8K & MATH-500 & AIME 2025 \\
\midrule
DFlash & \(\mathbf{0.4728} \pm 0\;[3]\) & \(\mathbf{0.1125} \pm 0\;[3]\) & \(\mathbf{0.219} \pm 0\;[3]\) \\
Pilot FFN 100\% & \(0 \pm 0\;[3]\) & \(0 \pm 0\;[3]\) & \(0 \pm 0\;[3]\) \\
Pilot FFN 75\% & \(0 \pm 0\;[3]\) & \(0 \pm 0\;[3]\) & \(0 \pm 0\;[3]\) \\
Pilot FFN 50\% & \(0 \pm 0\;[3]\) & \(0 \pm 0\;[3]\) & \(0 \pm 0\;[3]\) \\
\bottomrule
\end{tabular}\par\medskip
\textit{Code}\par\smallskip
\begin{tabular}{@{}>{\raggedright\arraybackslash}p{3.7cm}rrr@{}}
\toprule
Method & HumanEval & MBPP & LiveCodeBench \\
\midrule
DFlash & \(\mathbf{0.1071} \pm 0\;[3]\) & \(\mathbf{0.25} \pm 0\;[3]\) & \(\mathbf{0.065} \pm 0\;[3]\) \\
Pilot FFN 100\% & \(0 \pm 0\;[3]\) & \(0 \pm 0\;[3]\) & \(0 \pm 0\;[3]\) \\
Pilot FFN 75\% & \(0 \pm 0\;[3]\) & \(0 \pm 0\;[3]\) & \(0 \pm 0\;[3]\) \\
Pilot FFN 50\% & \(0 \pm 0\;[3]\) & \(0 \pm 0\;[3]\) & \(0 \pm 0\;[3]\) \\
\bottomrule
\end{tabular}\par\medskip
\textit{Conversation, instruction following, and summarization}\par\smallskip
\begin{tabular}{@{}>{\raggedright\arraybackslash}p{3.7cm}rrrr@{}}
\toprule
Method & MT-Bench & Alpaca & Arena-Hard & CNN/DailyMail \\
\midrule
DFlash & \(\mathbf{0.01994} \pm 0\;[3]\) & \(\mathbf{0.06186} \pm 0\;[3]\) & \(\mathbf{0.06209} \pm 0\;[3]\) & \(\mathbf{0.02542} \pm 0\;[3]\) \\
Pilot FFN 100\% & \(0 \pm 0\;[3]\) & \(0 \pm 0\;[3]\) & \(0 \pm 0\;[3]\) & \(0 \pm 0\;[3]\) \\
Pilot FFN 75\% & \(0 \pm 0\;[3]\) & \(0 \pm 0\;[3]\) & \(0 \pm 0\;[3]\) & \(0 \pm 0\;[3]\) \\
Pilot FFN 50\% & \(0 \pm 0\;[3]\) & \(0 \pm 0\;[3]\) & \(0 \pm 0\;[3]\) & \(0 \pm 0\;[3]\) \\
\bottomrule
\end{tabular}\par\medskip
\begin{minipage}{\linewidth}\footnotesize
For each draft position, count the rounds that proposed that position. Among those rounds, count how often the target accepted every draft token through that position before final output clipping. Divide the second count by the first.
The bold mean is the best observed value in each benchmark column. The comparison covers the recorded implementations under their stated configurations. At least two eligible method entries are required. Ties at the displayed precision share bold formatting. Bold formatting conveys an observed comparison and makes no statistical-significance claim. Brackets show the available repetition count. Response differences remain part of the interpretation. Higher values are preferred. Tree and linear proposal denominators form separate comparison groups.
Brackets give eligible repetition counts. Standard deviation appears when at least two repetitions are available. Each repetition uses three prompts. The evidence records preserve observation counts and measurement scopes. Empty cells lack an eligible measured value.
\end{minipage}
\end{table*}

\begin{table*}[!htp]
\centering\small
\setlength{\tabcolsep}{3pt}
\caption{Prefix acceptance probability by position; Position 9 (probability at each draft position).}
\label{tab:appendix-metric-prefix-survival-9}
\textit{Mathematics}\par\smallskip
\begin{tabular}{@{}>{\raggedright\arraybackslash}p{3.7cm}rrr@{}}
\toprule
Method & GSM8K & MATH-500 & AIME 2025 \\
\midrule
DFlash & \(0.4298 \pm 0\;[3]\) & \(0.07187 \pm 0\;[3]\) & \(0.1812 \pm 0\;[3]\) \\
\bottomrule
\end{tabular}\par\medskip
\textit{Code}\par\smallskip
\begin{tabular}{@{}>{\raggedright\arraybackslash}p{3.7cm}rrr@{}}
\toprule
Method & HumanEval & MBPP & LiveCodeBench \\
\midrule
DFlash & \(0.08929 \pm 0\;[3]\) & \(0.1705 \pm 0\;[3]\) & \(0.065 \pm 0\;[3]\) \\
\bottomrule
\end{tabular}\par\medskip
\textit{Conversation, instruction following, and summarization}\par\smallskip
\begin{tabular}{@{}>{\raggedright\arraybackslash}p{3.7cm}rrrr@{}}
\toprule
Method & MT-Bench & Alpaca & Arena-Hard & CNN/DailyMail \\
\midrule
DFlash & \(0.01994 \pm 0\;[3]\) & \(0.03608 \pm 0\;[3]\) & \(0.04575 \pm 0\;[3]\) & \(0.01695 \pm 0\;[3]\) \\
\bottomrule
\end{tabular}\par\medskip
\begin{minipage}{\linewidth}\footnotesize
For each draft position, count the rounds that proposed that position. Among those rounds, count how often the target accepted every draft token through that position before final output clipping. Divide the second count by the first.
The bold mean is the best observed value in each benchmark column. The comparison covers the recorded implementations under their stated configurations. At least two eligible method entries are required. Ties at the displayed precision share bold formatting. Bold formatting conveys an observed comparison and makes no statistical-significance claim. Brackets show the available repetition count. Response differences remain part of the interpretation. Higher values are preferred. Tree and linear proposal denominators form separate comparison groups.
Brackets give eligible repetition counts. Standard deviation appears when at least two repetitions are available. Each repetition uses three prompts. The evidence records preserve observation counts and measurement scopes. Empty cells lack an eligible measured value.
\end{minipage}
\end{table*}

\begin{table*}[!htp]
\centering\small
\setlength{\tabcolsep}{3pt}
\caption{Prefix acceptance probability by position; Tree-prefix diagnostic; Position 1 (probability at each draft position).}
\label{tab:appendix-metric-prefix-survival-tree-prefix-diagnostic-1}
\textit{Mathematics}\par\smallskip
\begin{tabular}{@{}>{\raggedright\arraybackslash}p{3.7cm}rrr@{}}
\toprule
Method & GSM8K & MATH-500 & AIME 2025 \\
\midrule
EAGLE-3 & \(0.9253 \pm 0\;[3]\) & \(0.9158 \pm 0\;[3]\) & \(0.8703 \pm 0\;[3]\) \\
\bottomrule
\end{tabular}\par\medskip
\textit{Code}\par\smallskip
\begin{tabular}{@{}>{\raggedright\arraybackslash}p{3.7cm}rrr@{}}
\toprule
Method & HumanEval & MBPP & LiveCodeBench \\
\midrule
EAGLE-3 & \(0.8514 \pm 0\;[3]\) & \(0.944 \pm 0\;[3]\) & \(0.8559 \pm 0\;[3]\) \\
\bottomrule
\end{tabular}\par\medskip
\textit{Conversation, instruction following, and summarization}\par\smallskip
\begin{tabular}{@{}>{\raggedright\arraybackslash}p{3.7cm}rrrr@{}}
\toprule
Method & MT-Bench & Alpaca & Arena-Hard & CNN/DailyMail \\
\midrule
EAGLE-3 & \(0.9493 \pm 0\;[3]\) & \(0.9948 \pm 0\;[3]\) & \(0.7994 \pm 0\;[3]\) & \(0.8755 \pm 0\;[3]\) \\
\bottomrule
\end{tabular}\par\medskip
\begin{minipage}{\linewidth}\footnotesize
For EAGLE-3, count verification rounds whose proposal tree contains a node at the stated depth. Among those rounds, count rounds whose accepted path reaches that depth. Divide accepted rounds by eligible rounds. Depth one contains the root's children. The root is a previously generated target token. Counts precede final output clipping. These tree-prefix diagnostics use the recorded branching tree. The EAGLE-3 paper measures n-alpha with a chain and defines n by the draft feature history. The tree values preserve their separate interpretation.
The bold mean is the best observed value in each benchmark column. The comparison covers the recorded implementations under their stated configurations. At least two eligible method entries are required. Ties at the displayed precision share bold formatting. Bold formatting conveys an observed comparison and makes no statistical-significance claim. Brackets show the available repetition count. Response differences remain part of the interpretation. Higher values are preferred. Tree and linear proposal denominators form separate comparison groups.
Brackets give eligible repetition counts. Standard deviation appears when at least two repetitions are available. Each repetition uses three prompts. The evidence records preserve observation counts and measurement scopes. Empty cells lack an eligible measured value.
\end{minipage}
\end{table*}

\begin{table*}[!htp]
\centering\small
\setlength{\tabcolsep}{3pt}
\caption{Prefix acceptance probability by position; Tree-prefix diagnostic; Position 2 (probability at each draft position).}
\label{tab:appendix-metric-prefix-survival-tree-prefix-diagnostic-2}
\textit{Mathematics}\par\smallskip
\begin{tabular}{@{}>{\raggedright\arraybackslash}p{3.7cm}rrr@{}}
\toprule
Method & GSM8K & MATH-500 & AIME 2025 \\
\midrule
EAGLE-3 & \(0.8333 \pm 0\;[3]\) & \(0.7806 \pm 0\;[3]\) & \(0.7781 \pm 0\;[3]\) \\
\bottomrule
\end{tabular}\par\medskip
\textit{Code}\par\smallskip
\begin{tabular}{@{}>{\raggedright\arraybackslash}p{3.7cm}rrr@{}}
\toprule
Method & HumanEval & MBPP & LiveCodeBench \\
\midrule
EAGLE-3 & \(0.8243 \pm 0\;[3]\) & \(0.88 \pm 0\;[3]\) & \(0.7336 \pm 0\;[3]\) \\
\bottomrule
\end{tabular}\par\medskip
\textit{Conversation, instruction following, and summarization}\par\smallskip
\begin{tabular}{@{}>{\raggedright\arraybackslash}p{3.7cm}rrrr@{}}
\toprule
Method & MT-Bench & Alpaca & Arena-Hard & CNN/DailyMail \\
\midrule
EAGLE-3 & \(0.8478 \pm 0\;[3]\) & \(0.9175 \pm 0\;[3]\) & \(0.678 \pm 0\;[3]\) & \(0.7427 \pm 0\;[3]\) \\
\bottomrule
\end{tabular}\par\medskip
\begin{minipage}{\linewidth}\footnotesize
For EAGLE-3, count verification rounds whose proposal tree contains a node at the stated depth. Among those rounds, count rounds whose accepted path reaches that depth. Divide accepted rounds by eligible rounds. Depth one contains the root's children. The root is a previously generated target token. Counts precede final output clipping. These tree-prefix diagnostics use the recorded branching tree. The EAGLE-3 paper measures n-alpha with a chain and defines n by the draft feature history. The tree values preserve their separate interpretation.
The bold mean is the best observed value in each benchmark column. The comparison covers the recorded implementations under their stated configurations. At least two eligible method entries are required. Ties at the displayed precision share bold formatting. Bold formatting conveys an observed comparison and makes no statistical-significance claim. Brackets show the available repetition count. Response differences remain part of the interpretation. Higher values are preferred. Tree and linear proposal denominators form separate comparison groups.
Brackets give eligible repetition counts. Standard deviation appears when at least two repetitions are available. Each repetition uses three prompts. The evidence records preserve observation counts and measurement scopes. Empty cells lack an eligible measured value.
\end{minipage}
\end{table*}

\begin{table*}[!htp]
\centering\small
\setlength{\tabcolsep}{3pt}
\caption{Prefix acceptance probability by position; Tree-prefix diagnostic; Position 3 (probability at each draft position).}
\label{tab:appendix-metric-prefix-survival-tree-prefix-diagnostic-3}
\textit{Mathematics}\par\smallskip
\begin{tabular}{@{}>{\raggedright\arraybackslash}p{3.7cm}rrr@{}}
\toprule
Method & GSM8K & MATH-500 & AIME 2025 \\
\midrule
EAGLE-3 & \(0.6839 \pm 0\;[3]\) & \(0.6097 \pm 0\;[3]\) & \(0.6805 \pm 0\;[3]\) \\
\bottomrule
\end{tabular}\par\medskip
\textit{Code}\par\smallskip
\begin{tabular}{@{}>{\raggedright\arraybackslash}p{3.7cm}rrr@{}}
\toprule
Method & HumanEval & MBPP & LiveCodeBench \\
\midrule
EAGLE-3 & \(0.6892 \pm 0\;[3]\) & \(0.8 \pm 0\;[3]\) & \(0.6288 \pm 0\;[3]\) \\
\bottomrule
\end{tabular}\par\medskip
\textit{Conversation, instruction following, and summarization}\par\smallskip
\begin{tabular}{@{}>{\raggedright\arraybackslash}p{3.7cm}rrrr@{}}
\toprule
Method & MT-Bench & Alpaca & Arena-Hard & CNN/DailyMail \\
\midrule
EAGLE-3 & \(0.7612 \pm 0\;[3]\) & \(0.8402 \pm 0\;[3]\) & \(0.5694 \pm 0\;[3]\) & \(0.527 \pm 0\;[3]\) \\
\bottomrule
\end{tabular}\par\medskip
\begin{minipage}{\linewidth}\footnotesize
For EAGLE-3, count verification rounds whose proposal tree contains a node at the stated depth. Among those rounds, count rounds whose accepted path reaches that depth. Divide accepted rounds by eligible rounds. Depth one contains the root's children. The root is a previously generated target token. Counts precede final output clipping. These tree-prefix diagnostics use the recorded branching tree. The EAGLE-3 paper measures n-alpha with a chain and defines n by the draft feature history. The tree values preserve their separate interpretation.
The bold mean is the best observed value in each benchmark column. The comparison covers the recorded implementations under their stated configurations. At least two eligible method entries are required. Ties at the displayed precision share bold formatting. Bold formatting conveys an observed comparison and makes no statistical-significance claim. Brackets show the available repetition count. Response differences remain part of the interpretation. Higher values are preferred. Tree and linear proposal denominators form separate comparison groups.
Brackets give eligible repetition counts. Standard deviation appears when at least two repetitions are available. Each repetition uses three prompts. The evidence records preserve observation counts and measurement scopes. Empty cells lack an eligible measured value.
\end{minipage}
\end{table*}

\begin{table*}[!htp]
\centering\small
\setlength{\tabcolsep}{3pt}
\caption{Conditional acceptance rate by position; Position 1 (\% at each draft position).}
\label{tab:appendix-metric-conditional-acceptance-1}
\textit{Mathematics}\par\smallskip
\begin{tabular}{@{}>{\raggedright\arraybackslash}p{3.7cm}rrr@{}}
\toprule
Method & GSM8K & MATH-500 & AIME 2025 \\
\midrule
DFlash & \(\mathbf{91.98} \pm 0\;[3]\) & \(\mathbf{85.94} \pm 0\;[3]\) & \(\mathbf{84.4} \pm 0\;[3]\) \\
Pilot FFN 100\% & \(24.5 \pm 0\;[3]\) & \(9.528 \pm 0\;[3]\) & \(12.1 \pm 0\;[3]\) \\
Pilot FFN 75\% & \(24.34 \pm 0\;[3]\) & \(8.539 \pm 0\;[3]\) & \(13.01 \pm 0\;[3]\) \\
Pilot FFN 50\% & \(20.91 \pm 0\;[3]\) & \(7.218 \pm 0\;[3]\) & \(10.5 \pm 0\;[3]\) \\
\bottomrule
\end{tabular}\par\medskip
\textit{Code}\par\smallskip
\begin{tabular}{@{}>{\raggedright\arraybackslash}p{3.7cm}rrr@{}}
\toprule
Method & HumanEval & MBPP & LiveCodeBench \\
\midrule
DFlash & \(\mathbf{80.36} \pm 0\;[3]\) & \(\mathbf{80.68} \pm 0\;[3]\) & \(\mathbf{79} \pm 0\;[3]\) \\
Pilot FFN 100\% & \(15.87 \pm 0\;[3]\) & \(17.8 \pm 0\;[3]\) & \(14.89 \pm 0\;[3]\) \\
Pilot FFN 75\% & \(13.62 \pm 0\;[3]\) & \(13.03 \pm 0\;[3]\) & \(14.13 \pm 0\;[3]\) \\
Pilot FFN 50\% & \(16.33 \pm 0\;[3]\) & \(18.11 \pm 0\;[3]\) & \(13.6 \pm 0\;[3]\) \\
\bottomrule
\end{tabular}\par\medskip
\textit{Conversation, instruction following, and summarization}\par\smallskip
\begin{tabular}{@{}>{\raggedright\arraybackslash}p{3.7cm}rrrr@{}}
\toprule
Method & MT-Bench & Alpaca & Arena-Hard & CNN/DailyMail \\
\midrule
DFlash & \(\mathbf{79.77} \pm 0\;[3]\) & \(\mathbf{81.44} \pm 0\;[3]\) & \(\mathbf{72.22} \pm 0\;[3]\) & \(\mathbf{77.12} \pm 0\;[3]\) \\
Pilot FFN 100\% & \(17.81 \pm 0\;[3]\) & \(15.82 \pm 0\;[3]\) & \(14.56 \pm 0\;[3]\) & \(15.37 \pm 0\;[3]\) \\
Pilot FFN 75\% & \(18.05 \pm 0\;[3]\) & \(15.73 \pm 0\;[3]\) & \(13.02 \pm 0\;[3]\) & \(16.12 \pm 0\;[3]\) \\
Pilot FFN 50\% & \(17.46 \pm 0\;[3]\) & \(15.03 \pm 0\;[3]\) & \(13.33 \pm 0\;[3]\) & \(13.46 \pm 0\;[3]\) \\
\bottomrule
\end{tabular}\par\medskip
\begin{minipage}{\linewidth}\footnotesize
At each draft position, count rounds that proposed the position and accepted every earlier proposal. Among these eligible rounds, count how often the current proposal passes. Divide passes by eligible rounds and multiply by 100. Report a result only when the eligible count is positive.
The bold mean is the best observed value in each benchmark column. The comparison covers the recorded implementations under their stated configurations. At least two eligible method entries are required. Ties at the displayed precision share bold formatting. Bold formatting conveys an observed comparison and makes no statistical-significance claim. Brackets show the available repetition count. Response differences remain part of the interpretation. Higher values are preferred. Tree and linear proposal denominators form separate comparison groups.
Brackets give eligible repetition counts. Standard deviation appears when at least two repetitions are available. Each repetition uses three prompts. The evidence records preserve observation counts and measurement scopes. Empty cells lack an eligible measured value.
\end{minipage}
\end{table*}

\begin{table*}[!htp]
\centering\small
\setlength{\tabcolsep}{3pt}
\caption{Conditional acceptance rate by position; Position 2 (\% at each draft position).}
\label{tab:appendix-metric-conditional-acceptance-2}
\textit{Mathematics}\par\smallskip
\begin{tabular}{@{}>{\raggedright\arraybackslash}p{3.7cm}rrr@{}}
\toprule
Method & GSM8K & MATH-500 & AIME 2025 \\
\midrule
DFlash & \(\mathbf{89.72} \pm 0\;[3]\) & \(\mathbf{79.27} \pm 0\;[3]\) & \(\mathbf{78.94} \pm 0\;[3]\) \\
Pilot FFN 100\% & \(0 \pm 0\;[3]\) & \(0 \pm 0\;[3]\) & \(0.2028 \pm 0\;[3]\) \\
Pilot FFN 75\% & \(1.818 \pm 0\;[3]\) & \(0 \pm 0\;[3]\) & \(1.905 \pm 0\;[3]\) \\
Pilot FFN 50\% & \(3.093 \pm 0\;[3]\) & \(0 \pm 0\;[3]\) & \(0.6912 \pm 0\;[3]\) \\
\bottomrule
\end{tabular}\par\medskip
\textit{Code}\par\smallskip
\begin{tabular}{@{}>{\raggedright\arraybackslash}p{3.7cm}rrr@{}}
\toprule
Method & HumanEval & MBPP & LiveCodeBench \\
\midrule
DFlash & \(\mathbf{86.67} \pm 0\;[3]\) & \(\mathbf{84.51} \pm 0\;[3]\) & \(\mathbf{76.58} \pm 0\;[3]\) \\
Pilot FFN 100\% & \(0 \pm 0\;[3]\) & \(1.471 \pm 0\;[3]\) & \(5.435 \pm 0\;[3]\) \\
Pilot FFN 75\% & \(0 \pm 0\;[3]\) & \(0 \pm 0\;[3]\) & \(4.545 \pm 0\;[3]\) \\
Pilot FFN 50\% & \(0 \pm 0\;[3]\) & \(1.449 \pm 0\;[3]\) & \(5.882 \pm 0\;[3]\) \\
\bottomrule
\end{tabular}\par\medskip
\textit{Conversation, instruction following, and summarization}\par\smallskip
\begin{tabular}{@{}>{\raggedright\arraybackslash}p{3.7cm}rrrr@{}}
\toprule
Method & MT-Bench & Alpaca & Arena-Hard & CNN/DailyMail \\
\midrule
DFlash & \(\mathbf{73.93} \pm 0\;[3]\) & \(\mathbf{77.85} \pm 0\;[3]\) & \(\mathbf{77.38} \pm 0\;[3]\) & \(\mathbf{69.23} \pm 0\;[3]\) \\
Pilot FFN 100\% & \(0 \pm 0\;[3]\) & \(0 \pm 0\;[3]\) & \(5.147 \pm 0\;[3]\) & \(0 \pm 0\;[3]\) \\
Pilot FFN 75\% & \(0 \pm 0\;[3]\) & \(4.348 \pm 0\;[3]\) & \(6.504 \pm 0\;[3]\) & \(0.9709 \pm 0\;[3]\) \\
Pilot FFN 50\% & \(0 \pm 0\;[3]\) & \(0 \pm 0\;[3]\) & \(3.968 \pm 0\;[3]\) & \(1.136 \pm 0\;[3]\) \\
\bottomrule
\end{tabular}\par\medskip
\begin{minipage}{\linewidth}\footnotesize
At each draft position, count rounds that proposed the position and accepted every earlier proposal. Among these eligible rounds, count how often the current proposal passes. Divide passes by eligible rounds and multiply by 100. Report a result only when the eligible count is positive.
The bold mean is the best observed value in each benchmark column. The comparison covers the recorded implementations under their stated configurations. At least two eligible method entries are required. Ties at the displayed precision share bold formatting. Bold formatting conveys an observed comparison and makes no statistical-significance claim. Brackets show the available repetition count. Response differences remain part of the interpretation. Higher values are preferred. Tree and linear proposal denominators form separate comparison groups.
Brackets give eligible repetition counts. Standard deviation appears when at least two repetitions are available. Each repetition uses three prompts. The evidence records preserve observation counts and measurement scopes. Empty cells lack an eligible measured value.
\end{minipage}
\end{table*}

\begin{table*}[!htp]
\centering\small
\setlength{\tabcolsep}{3pt}
\caption{Conditional acceptance rate by position; Position 3 (\% at each draft position).}
\label{tab:appendix-metric-conditional-acceptance-3}
\textit{Mathematics}\par\smallskip
\begin{tabular}{@{}>{\raggedright\arraybackslash}p{3.7cm}rrr@{}}
\toprule
Method & GSM8K & MATH-500 & AIME 2025 \\
\midrule
DFlash & \(\mathbf{90.28} \pm 0\;[3]\) & \(72.02 \pm 0\;[3]\) & \(\mathbf{80.03} \pm 0\;[3]\) \\
Pilot FFN 100\% &  &  & \(0 \pm 0\;[3]\) \\
Pilot FFN 75\% & \(0 \pm 0\;[3]\) &  & \(0 \pm 0\;[3]\) \\
Pilot FFN 50\% & \(0 \pm 0\;[3]\) &  & \(0 \pm 0\;[3]\) \\
\bottomrule
\end{tabular}\par\medskip
\textit{Code}\par\smallskip
\begin{tabular}{@{}>{\raggedright\arraybackslash}p{3.7cm}rrr@{}}
\toprule
Method & HumanEval & MBPP & LiveCodeBench \\
\midrule
DFlash & \(76.92 \pm 0\;[3]\) & \(\mathbf{76.67} \pm 0\;[3]\) & \(\mathbf{68.6} \pm 0\;[3]\) \\
Pilot FFN 100\% &  & \(0 \pm 0\;[3]\) & \(0 \pm 0\;[3]\) \\
Pilot FFN 75\% &  &  & \(0 \pm 0\;[3]\) \\
Pilot FFN 50\% &  & \(0 \pm 0\;[3]\) & \(0 \pm 0\;[3]\) \\
\bottomrule
\end{tabular}\par\medskip
\textit{Conversation, instruction following, and summarization}\par\smallskip
\begin{tabular}{@{}>{\raggedright\arraybackslash}p{3.7cm}rrrr@{}}
\toprule
Method & MT-Bench & Alpaca & Arena-Hard & CNN/DailyMail \\
\midrule
DFlash & \(68.6 \pm 0\;[3]\) & \(\mathbf{69.92} \pm 0\;[3]\) & \(\mathbf{74.85} \pm 0\;[3]\) & \(\mathbf{61.11} \pm 0\;[3]\) \\
Pilot FFN 100\% &  &  & \(0 \pm 0\;[3]\) &  \\
Pilot FFN 75\% &  & \(0 \pm 0\;[3]\) & \(0 \pm 0\;[3]\) & \(0 \pm 0\;[3]\) \\
Pilot FFN 50\% &  &  & \(0 \pm 0\;[3]\) & \(0 \pm 0\;[3]\) \\
\bottomrule
\end{tabular}\par\medskip
\begin{minipage}{\linewidth}\footnotesize
At each draft position, count rounds that proposed the position and accepted every earlier proposal. Among these eligible rounds, count how often the current proposal passes. Divide passes by eligible rounds and multiply by 100. Report a result only when the eligible count is positive.
The bold mean is the best observed value in each benchmark column. The comparison covers the recorded implementations under their stated configurations. At least two eligible method entries are required. Ties at the displayed precision share bold formatting. Bold formatting conveys an observed comparison and makes no statistical-significance claim. Brackets show the available repetition count. Response differences remain part of the interpretation. Higher values are preferred. Tree and linear proposal denominators form separate comparison groups.
Brackets give eligible repetition counts. Standard deviation appears when at least two repetitions are available. Each repetition uses three prompts. The evidence records preserve observation counts and measurement scopes. Empty cells lack an eligible measured value.
\end{minipage}
\end{table*}

\begin{table*}[!htp]
\centering\small
\setlength{\tabcolsep}{3pt}
\caption{Conditional acceptance rate by position; Position 4 (\% at each draft position).}
\label{tab:appendix-metric-conditional-acceptance-4}
\textit{Mathematics}\par\smallskip
\begin{tabular}{@{}>{\raggedright\arraybackslash}p{3.7cm}rrr@{}}
\toprule
Method & GSM8K & MATH-500 & AIME 2025 \\
\midrule
DFlash & \(86.54 \pm 0\;[3]\) & \(77.71 \pm 0\;[3]\) & \(83.87 \pm 0\;[3]\) \\
\bottomrule
\end{tabular}\par\medskip
\textit{Code}\par\smallskip
\begin{tabular}{@{}>{\raggedright\arraybackslash}p{3.7cm}rrr@{}}
\toprule
Method & HumanEval & MBPP & LiveCodeBench \\
\midrule
DFlash & \(76.67 \pm 0\;[3]\) & \(89.13 \pm 0\;[3]\) & \(68.67 \pm 0\;[3]\) \\
\bottomrule
\end{tabular}\par\medskip
\textit{Conversation, instruction following, and summarization}\par\smallskip
\begin{tabular}{@{}>{\raggedright\arraybackslash}p{3.7cm}rrrr@{}}
\toprule
Method & MT-Bench & Alpaca & Arena-Hard & CNN/DailyMail \\
\midrule
DFlash & \(64.79 \pm 0\;[3]\) & \(67.44 \pm 0\;[3]\) & \(70.31 \pm 0\;[3]\) & \(61.04 \pm 0\;[3]\) \\
\bottomrule
\end{tabular}\par\medskip
\begin{minipage}{\linewidth}\footnotesize
At each draft position, count rounds that proposed the position and accepted every earlier proposal. Among these eligible rounds, count how often the current proposal passes. Divide passes by eligible rounds and multiply by 100. Report a result only when the eligible count is positive.
The bold mean is the best observed value in each benchmark column. The comparison covers the recorded implementations under their stated configurations. At least two eligible method entries are required. Ties at the displayed precision share bold formatting. Bold formatting conveys an observed comparison and makes no statistical-significance claim. Brackets show the available repetition count. Response differences remain part of the interpretation. Higher values are preferred. Tree and linear proposal denominators form separate comparison groups.
Brackets give eligible repetition counts. Standard deviation appears when at least two repetitions are available. Each repetition uses three prompts. The evidence records preserve observation counts and measurement scopes. Empty cells lack an eligible measured value.
\end{minipage}
\end{table*}

\begin{table*}[!htp]
\centering\small
\setlength{\tabcolsep}{3pt}
\caption{Conditional acceptance rate by position; Position 5 (\% at each draft position).}
\label{tab:appendix-metric-conditional-acceptance-5}
\textit{Mathematics}\par\smallskip
\begin{tabular}{@{}>{\raggedright\arraybackslash}p{3.7cm}rrr@{}}
\toprule
Method & GSM8K & MATH-500 & AIME 2025 \\
\midrule
DFlash & \(91.11 \pm 0\;[3]\) & \(79.51 \pm 0\;[3]\) & \(79.49 \pm 0\;[3]\) \\
\bottomrule
\end{tabular}\par\medskip
\textit{Code}\par\smallskip
\begin{tabular}{@{}>{\raggedright\arraybackslash}p{3.7cm}rrr@{}}
\toprule
Method & HumanEval & MBPP & LiveCodeBench \\
\midrule
DFlash & \(82.61 \pm 0\;[3]\) & \(87.8 \pm 0\;[3]\) & \(73.68 \pm 0\;[3]\) \\
\bottomrule
\end{tabular}\par\medskip
\textit{Conversation, instruction following, and summarization}\par\smallskip
\begin{tabular}{@{}>{\raggedright\arraybackslash}p{3.7cm}rrrr@{}}
\toprule
Method & MT-Bench & Alpaca & Arena-Hard & CNN/DailyMail \\
\midrule
DFlash & \(65.22 \pm 0\;[3]\) & \(60.34 \pm 0\;[3]\) & \(65.56 \pm 0\;[3]\) & \(53.19 \pm 0\;[3]\) \\
\bottomrule
\end{tabular}\par\medskip
\begin{minipage}{\linewidth}\footnotesize
At each draft position, count rounds that proposed the position and accepted every earlier proposal. Among these eligible rounds, count how often the current proposal passes. Divide passes by eligible rounds and multiply by 100. Report a result only when the eligible count is positive.
The bold mean is the best observed value in each benchmark column. The comparison covers the recorded implementations under their stated configurations. At least two eligible method entries are required. Ties at the displayed precision share bold formatting. Bold formatting conveys an observed comparison and makes no statistical-significance claim. Brackets show the available repetition count. Response differences remain part of the interpretation. Higher values are preferred. Tree and linear proposal denominators form separate comparison groups.
Brackets give eligible repetition counts. Standard deviation appears when at least two repetitions are available. Each repetition uses three prompts. The evidence records preserve observation counts and measurement scopes. Empty cells lack an eligible measured value.
\end{minipage}
\end{table*}

\begin{table*}[!htp]
\centering\small
\setlength{\tabcolsep}{3pt}
\caption{Conditional acceptance rate by position; Position 6 (\% at each draft position).}
\label{tab:appendix-metric-conditional-acceptance-6}
\textit{Mathematics}\par\smallskip
\begin{tabular}{@{}>{\raggedright\arraybackslash}p{3.7cm}rrr@{}}
\toprule
Method & GSM8K & MATH-500 & AIME 2025 \\
\midrule
DFlash & \(90.24 \pm 0\;[3]\) & \(75.26 \pm 0\;[3]\) & \(83.23 \pm 0\;[3]\) \\
\bottomrule
\end{tabular}\par\medskip
\textit{Code}\par\smallskip
\begin{tabular}{@{}>{\raggedright\arraybackslash}p{3.7cm}rrr@{}}
\toprule
Method & HumanEval & MBPP & LiveCodeBench \\
\midrule
DFlash & \(78.95 \pm 0\;[3]\) & \(77.78 \pm 0\;[3]\) & \(69.05 \pm 0\;[3]\) \\
\bottomrule
\end{tabular}\par\medskip
\textit{Conversation, instruction following, and summarization}\par\smallskip
\begin{tabular}{@{}>{\raggedright\arraybackslash}p{3.7cm}rrrr@{}}
\toprule
Method & MT-Bench & Alpaca & Arena-Hard & CNN/DailyMail \\
\midrule
DFlash & \(53.33 \pm 0\;[3]\) & \(60 \pm 0\;[3]\) & \(62.71 \pm 0\;[3]\) & \(64 \pm 0\;[3]\) \\
\bottomrule
\end{tabular}\par\medskip
\begin{minipage}{\linewidth}\footnotesize
At each draft position, count rounds that proposed the position and accepted every earlier proposal. Among these eligible rounds, count how often the current proposal passes. Divide passes by eligible rounds and multiply by 100. Report a result only when the eligible count is positive.
The bold mean is the best observed value in each benchmark column. The comparison covers the recorded implementations under their stated configurations. At least two eligible method entries are required. Ties at the displayed precision share bold formatting. Bold formatting conveys an observed comparison and makes no statistical-significance claim. Brackets show the available repetition count. Response differences remain part of the interpretation. Higher values are preferred. Tree and linear proposal denominators form separate comparison groups.
Brackets give eligible repetition counts. Standard deviation appears when at least two repetitions are available. Each repetition uses three prompts. The evidence records preserve observation counts and measurement scopes. Empty cells lack an eligible measured value.
\end{minipage}
\end{table*}

\begin{table*}[!htp]
\centering\small
\setlength{\tabcolsep}{3pt}
\caption{Conditional acceptance rate by position; Position 7 (\% at each draft position).}
\label{tab:appendix-metric-conditional-acceptance-7}
\textit{Mathematics}\par\smallskip
\begin{tabular}{@{}>{\raggedright\arraybackslash}p{3.7cm}rrr@{}}
\toprule
Method & GSM8K & MATH-500 & AIME 2025 \\
\midrule
DFlash & \(94.59 \pm 0\;[3]\) & \(73.97 \pm 0\;[3]\) & \(87.98 \pm 0\;[3]\) \\
\bottomrule
\end{tabular}\par\medskip
\textit{Code}\par\smallskip
\begin{tabular}{@{}>{\raggedright\arraybackslash}p{3.7cm}rrr@{}}
\toprule
Method & HumanEval & MBPP & LiveCodeBench \\
\midrule
DFlash & \(66.67 \pm 0\;[3]\) & \(89.29 \pm 0\;[3]\) & \(62.07 \pm 0\;[3]\) \\
\bottomrule
\end{tabular}\par\medskip
\textit{Conversation, instruction following, and summarization}\par\smallskip
\begin{tabular}{@{}>{\raggedright\arraybackslash}p{3.7cm}rrrr@{}}
\toprule
Method & MT-Bench & Alpaca & Arena-Hard & CNN/DailyMail \\
\midrule
DFlash & \(46.88 \pm 0\;[3]\) & \(76.19 \pm 0\;[3]\) & \(72.97 \pm 0\;[3]\) & \(56.25 \pm 0\;[3]\) \\
\bottomrule
\end{tabular}\par\medskip
\begin{minipage}{\linewidth}\footnotesize
At each draft position, count rounds that proposed the position and accepted every earlier proposal. Among these eligible rounds, count how often the current proposal passes. Divide passes by eligible rounds and multiply by 100. Report a result only when the eligible count is positive.
The bold mean is the best observed value in each benchmark column. The comparison covers the recorded implementations under their stated configurations. At least two eligible method entries are required. Ties at the displayed precision share bold formatting. Bold formatting conveys an observed comparison and makes no statistical-significance claim. Brackets show the available repetition count. Response differences remain part of the interpretation. Higher values are preferred. Tree and linear proposal denominators form separate comparison groups.
Brackets give eligible repetition counts. Standard deviation appears when at least two repetitions are available. Each repetition uses three prompts. The evidence records preserve observation counts and measurement scopes. Empty cells lack an eligible measured value.
\end{minipage}
\end{table*}

\begin{table*}[!htp]
\centering\small
\setlength{\tabcolsep}{3pt}
\caption{Conditional acceptance rate by position; Position 8 (\% at each draft position).}
\label{tab:appendix-metric-conditional-acceptance-8}
\textit{Mathematics}\par\smallskip
\begin{tabular}{@{}>{\raggedright\arraybackslash}p{3.7cm}rrr@{}}
\toprule
Method & GSM8K & MATH-500 & AIME 2025 \\
\midrule
DFlash & \(94.29 \pm 0\;[3]\) & \(66.67 \pm 0\;[3]\) & \(84.14 \pm 0\;[3]\) \\
\bottomrule
\end{tabular}\par\medskip
\textit{Code}\par\smallskip
\begin{tabular}{@{}>{\raggedright\arraybackslash}p{3.7cm}rrr@{}}
\toprule
Method & HumanEval & MBPP & LiveCodeBench \\
\midrule
DFlash & \(60 \pm 0\;[3]\) & \(88 \pm 0\;[3]\) & \(72.22 \pm 0\;[3]\) \\
\bottomrule
\end{tabular}\par\medskip
\textit{Conversation, instruction following, and summarization}\par\smallskip
\begin{tabular}{@{}>{\raggedright\arraybackslash}p{3.7cm}rrrr@{}}
\toprule
Method & MT-Bench & Alpaca & Arena-Hard & CNN/DailyMail \\
\midrule
DFlash & \(46.67 \pm 0\;[3]\) & \(75 \pm 0\;[3]\) & \(70.37 \pm 0\;[3]\) & \(66.67 \pm 0\;[3]\) \\
\bottomrule
\end{tabular}\par\medskip
\begin{minipage}{\linewidth}\footnotesize
At each draft position, count rounds that proposed the position and accepted every earlier proposal. Among these eligible rounds, count how often the current proposal passes. Divide passes by eligible rounds and multiply by 100. Report a result only when the eligible count is positive.
The bold mean is the best observed value in each benchmark column. The comparison covers the recorded implementations under their stated configurations. At least two eligible method entries are required. Ties at the displayed precision share bold formatting. Bold formatting conveys an observed comparison and makes no statistical-significance claim. Brackets show the available repetition count. Response differences remain part of the interpretation. Higher values are preferred. Tree and linear proposal denominators form separate comparison groups.
Brackets give eligible repetition counts. Standard deviation appears when at least two repetitions are available. Each repetition uses three prompts. The evidence records preserve observation counts and measurement scopes. Empty cells lack an eligible measured value.
\end{minipage}
\end{table*}

\begin{table*}[!htp]
\centering\small
\setlength{\tabcolsep}{3pt}
\caption{Conditional acceptance rate by position; Position 9 (\% at each draft position).}
\label{tab:appendix-metric-conditional-acceptance-9}
\textit{Mathematics}\par\smallskip
\begin{tabular}{@{}>{\raggedright\arraybackslash}p{3.7cm}rrr@{}}
\toprule
Method & GSM8K & MATH-500 & AIME 2025 \\
\midrule
DFlash & \(90.91 \pm 0\;[3]\) & \(63.89 \pm 0\;[3]\) & \(82.72 \pm 0\;[3]\) \\
\bottomrule
\end{tabular}\par\medskip
\textit{Code}\par\smallskip
\begin{tabular}{@{}>{\raggedright\arraybackslash}p{3.7cm}rrr@{}}
\toprule
Method & HumanEval & MBPP & LiveCodeBench \\
\midrule
DFlash & \(83.33 \pm 0\;[3]\) & \(68.18 \pm 0\;[3]\) & \(100 \pm 0\;[3]\) \\
\bottomrule
\end{tabular}\par\medskip
\textit{Conversation, instruction following, and summarization}\par\smallskip
\begin{tabular}{@{}>{\raggedright\arraybackslash}p{3.7cm}rrrr@{}}
\toprule
Method & MT-Bench & Alpaca & Arena-Hard & CNN/DailyMail \\
\midrule
DFlash & \(100 \pm 0\;[3]\) & \(58.33 \pm 0\;[3]\) & \(73.68 \pm 0\;[3]\) & \(66.67 \pm 0\;[3]\) \\
\bottomrule
\end{tabular}\par\medskip
\begin{minipage}{\linewidth}\footnotesize
At each draft position, count rounds that proposed the position and accepted every earlier proposal. Among these eligible rounds, count how often the current proposal passes. Divide passes by eligible rounds and multiply by 100. Report a result only when the eligible count is positive.
The bold mean is the best observed value in each benchmark column. The comparison covers the recorded implementations under their stated configurations. At least two eligible method entries are required. Ties at the displayed precision share bold formatting. Bold formatting conveys an observed comparison and makes no statistical-significance claim. Brackets show the available repetition count. Response differences remain part of the interpretation. Higher values are preferred. Tree and linear proposal denominators form separate comparison groups.
Brackets give eligible repetition counts. Standard deviation appears when at least two repetitions are available. Each repetition uses three prompts. The evidence records preserve observation counts and measurement scopes. Empty cells lack an eligible measured value.
\end{minipage}
\end{table*}

\begin{table*}[!htp]
\centering\small
\setlength{\tabcolsep}{3pt}
\caption{Conditional acceptance rate by position; Tree-prefix diagnostic; Position 1 (\% at each draft position).}
\label{tab:appendix-metric-conditional-acceptance-tree-prefix-diagnostic-1}
\textit{Mathematics}\par\smallskip
\begin{tabular}{@{}>{\raggedright\arraybackslash}p{3.7cm}rrr@{}}
\toprule
Method & GSM8K & MATH-500 & AIME 2025 \\
\midrule
EAGLE-3 & \(92.53 \pm 0\;[3]\) & \(91.58 \pm 0\;[3]\) & \(87.03 \pm 0\;[3]\) \\
\bottomrule
\end{tabular}\par\medskip
\textit{Code}\par\smallskip
\begin{tabular}{@{}>{\raggedright\arraybackslash}p{3.7cm}rrr@{}}
\toprule
Method & HumanEval & MBPP & LiveCodeBench \\
\midrule
EAGLE-3 & \(85.14 \pm 0\;[3]\) & \(94.4 \pm 0\;[3]\) & \(85.59 \pm 0\;[3]\) \\
\bottomrule
\end{tabular}\par\medskip
\textit{Conversation, instruction following, and summarization}\par\smallskip
\begin{tabular}{@{}>{\raggedright\arraybackslash}p{3.7cm}rrrr@{}}
\toprule
Method & MT-Bench & Alpaca & Arena-Hard & CNN/DailyMail \\
\midrule
EAGLE-3 & \(94.93 \pm 0\;[3]\) & \(99.48 \pm 0\;[3]\) & \(79.94 \pm 0\;[3]\) & \(87.55 \pm 0\;[3]\) \\
\bottomrule
\end{tabular}\par\medskip
\begin{minipage}{\linewidth}\footnotesize
For EAGLE-3, count verification rounds whose accepted path reaches the preceding depth and has at least one proposed child at the stated depth. Among those rounds, count rounds that accept a child. Divide accepted rounds by eligible rounds and multiply by 100. A leaf has no eligible next proposal. Counts precede final output clipping. These tree-prefix diagnostics use the recorded branching tree. The EAGLE-3 paper measures n-alpha with a chain and defines n by the draft feature history. The tree values preserve their separate interpretation.
The bold mean is the best observed value in each benchmark column. The comparison covers the recorded implementations under their stated configurations. At least two eligible method entries are required. Ties at the displayed precision share bold formatting. Bold formatting conveys an observed comparison and makes no statistical-significance claim. Brackets show the available repetition count. Response differences remain part of the interpretation. Higher values are preferred. Tree and linear proposal denominators form separate comparison groups.
Brackets give eligible repetition counts. Standard deviation appears when at least two repetitions are available. Each repetition uses three prompts. The evidence records preserve observation counts and measurement scopes. Empty cells lack an eligible measured value.
\end{minipage}
\end{table*}

\begin{table*}[!htp]
\centering\small
\setlength{\tabcolsep}{3pt}
\caption{Conditional acceptance rate by position; Tree-prefix diagnostic; Position 2 (\% at each draft position).}
\label{tab:appendix-metric-conditional-acceptance-tree-prefix-diagnostic-2}
\textit{Mathematics}\par\smallskip
\begin{tabular}{@{}>{\raggedright\arraybackslash}p{3.7cm}rrr@{}}
\toprule
Method & GSM8K & MATH-500 & AIME 2025 \\
\midrule
EAGLE-3 & \(90.06 \pm 0\;[3]\) & \(85.96 \pm 0\;[3]\) & \(89.7 \pm 0\;[3]\) \\
\bottomrule
\end{tabular}\par\medskip
\textit{Code}\par\smallskip
\begin{tabular}{@{}>{\raggedright\arraybackslash}p{3.7cm}rrr@{}}
\toprule
Method & HumanEval & MBPP & LiveCodeBench \\
\midrule
EAGLE-3 & \(96.83 \pm 0\;[3]\) & \(93.22 \pm 0\;[3]\) & \(86.6 \pm 0\;[3]\) \\
\bottomrule
\end{tabular}\par\medskip
\textit{Conversation, instruction following, and summarization}\par\smallskip
\begin{tabular}{@{}>{\raggedright\arraybackslash}p{3.7cm}rrrr@{}}
\toprule
Method & MT-Bench & Alpaca & Arena-Hard & CNN/DailyMail \\
\midrule
EAGLE-3 & \(90.45 \pm 0\;[3]\) & \(92.23 \pm 0\;[3]\) & \(86.64 \pm 0\;[3]\) & \(85.65 \pm 0\;[3]\) \\
\bottomrule
\end{tabular}\par\medskip
\begin{minipage}{\linewidth}\footnotesize
For EAGLE-3, count verification rounds whose accepted path reaches the preceding depth and has at least one proposed child at the stated depth. Among those rounds, count rounds that accept a child. Divide accepted rounds by eligible rounds and multiply by 100. A leaf has no eligible next proposal. Counts precede final output clipping. These tree-prefix diagnostics use the recorded branching tree. The EAGLE-3 paper measures n-alpha with a chain and defines n by the draft feature history. The tree values preserve their separate interpretation.
The bold mean is the best observed value in each benchmark column. The comparison covers the recorded implementations under their stated configurations. At least two eligible method entries are required. Ties at the displayed precision share bold formatting. Bold formatting conveys an observed comparison and makes no statistical-significance claim. Brackets show the available repetition count. Response differences remain part of the interpretation. Higher values are preferred. Tree and linear proposal denominators form separate comparison groups.
Brackets give eligible repetition counts. Standard deviation appears when at least two repetitions are available. Each repetition uses three prompts. The evidence records preserve observation counts and measurement scopes. Empty cells lack an eligible measured value.
\end{minipage}
\end{table*}

\begin{table*}[!htp]
\centering\small
\setlength{\tabcolsep}{3pt}
\caption{Conditional acceptance rate by position; Tree-prefix diagnostic; Position 3 (\% at each draft position).}
\label{tab:appendix-metric-conditional-acceptance-tree-prefix-diagnostic-3}
\textit{Mathematics}\par\smallskip
\begin{tabular}{@{}>{\raggedright\arraybackslash}p{3.7cm}rrr@{}}
\toprule
Method & GSM8K & MATH-500 & AIME 2025 \\
\midrule
EAGLE-3 & \(90.15 \pm 0\;[3]\) & \(79.93 \pm 0\;[3]\) & \(89.95 \pm 0\;[3]\) \\
\bottomrule
\end{tabular}\par\medskip
\textit{Code}\par\smallskip
\begin{tabular}{@{}>{\raggedright\arraybackslash}p{3.7cm}rrr@{}}
\toprule
Method & HumanEval & MBPP & LiveCodeBench \\
\midrule
EAGLE-3 & \(83.61 \pm 0\;[3]\) & \(91.74 \pm 0\;[3]\) & \(89.44 \pm 0\;[3]\) \\
\bottomrule
\end{tabular}\par\medskip
\textit{Conversation, instruction following, and summarization}\par\smallskip
\begin{tabular}{@{}>{\raggedright\arraybackslash}p{3.7cm}rrrr@{}}
\toprule
Method & MT-Bench & Alpaca & Arena-Hard & CNN/DailyMail \\
\midrule
EAGLE-3 & \(92.39 \pm 0\;[3]\) & \(94.22 \pm 0\;[3]\) & \(87.39 \pm 0\;[3]\) & \(76.05 \pm 0\;[3]\) \\
\bottomrule
\end{tabular}\par\medskip
\begin{minipage}{\linewidth}\footnotesize
For EAGLE-3, count verification rounds whose accepted path reaches the preceding depth and has at least one proposed child at the stated depth. Among those rounds, count rounds that accept a child. Divide accepted rounds by eligible rounds and multiply by 100. A leaf has no eligible next proposal. Counts precede final output clipping. These tree-prefix diagnostics use the recorded branching tree. The EAGLE-3 paper measures n-alpha with a chain and defines n by the draft feature history. The tree values preserve their separate interpretation.
The bold mean is the best observed value in each benchmark column. The comparison covers the recorded implementations under their stated configurations. At least two eligible method entries are required. Ties at the displayed precision share bold formatting. Bold formatting conveys an observed comparison and makes no statistical-significance claim. Brackets show the available repetition count. Response differences remain part of the interpretation. Higher values are preferred. Tree and linear proposal denominators form separate comparison groups.
Brackets give eligible repetition counts. Standard deviation appears when at least two repetitions are available. Each repetition uses three prompts. The evidence records preserve observation counts and measurement scopes. Empty cells lack an eligible measured value.
\end{minipage}
\end{table*}

\begin{table*}[!htp]
\centering\small
\setlength{\tabcolsep}{3pt}
\caption{Draft latency, verification latency and round latency; Round latency (ms/round).}
\label{tab:appendix-metric-component-time-round-latency-none}
\textit{Mathematics}\par\smallskip
\begin{tabular}{@{}>{\raggedright\arraybackslash}p{3.7cm}rrr@{}}
\toprule
Method & GSM8K & MATH-500 & AIME 2025 \\
\midrule
EAGLE-3 & \(33.45 \pm 0.09\;[3]\) & \(33.43 \pm 0.11\;[3]\) & \(33.56 \pm 0.054\;[3]\) \\
DFlash &  &  &  \\
\bottomrule
\end{tabular}\par\medskip
\textit{Code}\par\smallskip
\begin{tabular}{@{}>{\raggedright\arraybackslash}p{3.7cm}rrr@{}}
\toprule
Method & HumanEval & MBPP & LiveCodeBench \\
\midrule
EAGLE-3 & \(32.94 \pm 0.29\;[3]\) & \(33.18 \pm 0.2\;[3]\) & \(33.37 \pm 0.2\;[3]\) \\
DFlash &  &  &  \\
\bottomrule
\end{tabular}\par\medskip
\textit{Conversation, instruction following, and summarization}\par\smallskip
\begin{tabular}{@{}>{\raggedright\arraybackslash}p{3.7cm}rrrr@{}}
\toprule
Method & MT-Bench & Alpaca & Arena-Hard & CNN/DailyMail \\
\midrule
EAGLE-3 & \(33.44 \pm 0.21\;[3]\) & \(33.36 \pm 0.2\;[3]\) & \(33.47 \pm 0.023\;[3]\) & \(33.49 \pm 0.11\;[3]\) \\
DFlash & \(31.25\;[1]\) & \(31.32\;[1]\) & \(31.36\;[1]\) & \(31.41\;[1]\) \\
\bottomrule
\end{tabular}\par\medskip
\begin{minipage}{\linewidth}\footnotesize
In a separate timing trace, measure a full cycle from drafting through verification and the final state update. Measure the time spent drafting, checking and moving data within that cycle. Record which operations overlap and average times over the observed cycles. Round latency uses host monotonic timestamps from drafting through the final state update. Draft and verification latency use CUDA events around their recorded operations. EAGLE-3's draft interval covers proposal generation. Draft extension remains inside the full round. P-EAGLE links each draft to its later verification, and its consecutive full-round intervals can overlap. Diagnostic measurements include synchronization and recording overhead. The evidence identifies each operation and records the overlap counts.
Component times support an explanation of measured runtime. Each method performs different work per verification round. A shorter component can still accompany a longer whole response. The table therefore reports component durations with equal emphasis.
Brackets give eligible repetition counts. Standard deviation appears when at least two repetitions are available. Each repetition uses three prompts. The evidence records preserve observation counts and measurement scopes. Empty cells lack an eligible measured value.
\end{minipage}
\end{table*}

\begin{table*}[!htp]
\centering\small
\setlength{\tabcolsep}{3pt}
\caption{Draft latency, verification latency and round latency; Draft latency (ms/round).}
\label{tab:appendix-metric-component-time-draft-latency-none}
\textit{Mathematics}\par\smallskip
\begin{tabular}{@{}>{\raggedright\arraybackslash}p{3.7cm}rrr@{}}
\toprule
Method & GSM8K & MATH-500 & AIME 2025 \\
\midrule
EAGLE-3 & \(4.365 \pm 0.053\;[3]\) & \(4.359 \pm 0.042\;[3]\) & \(4.343 \pm 0.014\;[3]\) \\
DFlash &  &  &  \\
\bottomrule
\end{tabular}\par\medskip
\textit{Code}\par\smallskip
\begin{tabular}{@{}>{\raggedright\arraybackslash}p{3.7cm}rrr@{}}
\toprule
Method & HumanEval & MBPP & LiveCodeBench \\
\midrule
EAGLE-3 & \(4.179 \pm 0.058\;[3]\) & \(4.301 \pm 0.07\;[3]\) & \(4.343 \pm 0.077\;[3]\) \\
DFlash &  &  &  \\
\bottomrule
\end{tabular}\par\medskip
\textit{Conversation, instruction following, and summarization}\par\smallskip
\begin{tabular}{@{}>{\raggedright\arraybackslash}p{3.7cm}rrrr@{}}
\toprule
Method & MT-Bench & Alpaca & Arena-Hard & CNN/DailyMail \\
\midrule
EAGLE-3 & \(4.355 \pm 0.072\;[3]\) & \(4.367 \pm 0.06\;[3]\) & \(4.357 \pm 0.011\;[3]\) & \(4.347 \pm 0.026\;[3]\) \\
DFlash & \(5.777\;[1]\) & \(5.861\;[1]\) & \(5.832\;[1]\) & \(5.856\;[1]\) \\
\bottomrule
\end{tabular}\par\medskip
\begin{minipage}{\linewidth}\footnotesize
In a separate timing trace, measure a full cycle from drafting through verification and the final state update. Measure the time spent drafting, checking and moving data within that cycle. Record which operations overlap and average times over the observed cycles. Round latency uses host monotonic timestamps from drafting through the final state update. Draft and verification latency use CUDA events around their recorded operations. EAGLE-3's draft interval covers proposal generation. Draft extension remains inside the full round. P-EAGLE links each draft to its later verification, and its consecutive full-round intervals can overlap. Diagnostic measurements include synchronization and recording overhead. The evidence identifies each operation and records the overlap counts.
Component times support an explanation of measured runtime. Each method performs different work per verification round. A shorter component can still accompany a longer whole response. The table therefore reports component durations with equal emphasis.
Brackets give eligible repetition counts. Standard deviation appears when at least two repetitions are available. Each repetition uses three prompts. The evidence records preserve observation counts and measurement scopes. Empty cells lack an eligible measured value.
\end{minipage}
\end{table*}

\begin{table*}[!htp]
\centering\small
\setlength{\tabcolsep}{3pt}
\caption{Draft latency, verification latency and round latency; Verification latency (ms/round).}
\label{tab:appendix-metric-component-time-verification-latency-none}
\textit{Mathematics}\par\smallskip
\begin{tabular}{@{}>{\raggedright\arraybackslash}p{3.7cm}rrr@{}}
\toprule
Method & GSM8K & MATH-500 & AIME 2025 \\
\midrule
EAGLE-3 & \(24.18 \pm 0.023\;[3]\) & \(24.21 \pm 0.021\;[3]\) & \(24.33 \pm 0.0039\;[3]\) \\
DFlash &  &  &  \\
\bottomrule
\end{tabular}\par\medskip
\textit{Code}\par\smallskip
\begin{tabular}{@{}>{\raggedright\arraybackslash}p{3.7cm}rrr@{}}
\toprule
Method & HumanEval & MBPP & LiveCodeBench \\
\midrule
EAGLE-3 & \(24.12 \pm 0.026\;[3]\) & \(24.11 \pm 0.028\;[3]\) & \(24.23 \pm 0.031\;[3]\) \\
DFlash &  &  &  \\
\bottomrule
\end{tabular}\par\medskip
\textit{Conversation, instruction following, and summarization}\par\smallskip
\begin{tabular}{@{}>{\raggedright\arraybackslash}p{3.7cm}rrrr@{}}
\toprule
Method & MT-Bench & Alpaca & Arena-Hard & CNN/DailyMail \\
\midrule
EAGLE-3 & \(24.21 \pm 0.031\;[3]\) & \(24.14 \pm 0.041\;[3]\) & \(24.26 \pm 0.0095\;[3]\) & \(24.32 \pm 0.019\;[3]\) \\
DFlash & \(23.62\;[1]\) & \(23.52\;[1]\) & \(23.59\;[1]\) & \(23.63\;[1]\) \\
\bottomrule
\end{tabular}\par\medskip
\begin{minipage}{\linewidth}\footnotesize
In a separate timing trace, measure a full cycle from drafting through verification and the final state update. Measure the time spent drafting, checking and moving data within that cycle. Record which operations overlap and average times over the observed cycles. Round latency uses host monotonic timestamps from drafting through the final state update. Draft and verification latency use CUDA events around their recorded operations. EAGLE-3's draft interval covers proposal generation. Draft extension remains inside the full round. P-EAGLE links each draft to its later verification, and its consecutive full-round intervals can overlap. Diagnostic measurements include synchronization and recording overhead. The evidence identifies each operation and records the overlap counts.
Component times support an explanation of measured runtime. Each method performs different work per verification round. A shorter component can still accompany a longer whole response. The table therefore reports component durations with equal emphasis.
Brackets give eligible repetition counts. Standard deviation appears when at least two repetitions are available. Each repetition uses three prompts. The evidence records preserve observation counts and measurement scopes. Empty cells lack an eligible measured value.
\end{minipage}
\end{table*}

\end{document}
